\chapter{清初：征服、摄政与帝国整合}
\label{ch:early-qing}
\begin{figure}[htbp]
  \centering
  \begin{tikzpicture}[
    x=1cm,y=1cm,
    date/.style={font=\scriptsize\bfseries, text=dynasty},
    event/.style={draw=timeline, fill=paper, rounded corners=2pt, align=center,
      inner sep=3pt, font=\scriptsize, text=ink},
    note/.style={font=\scriptsize\color{source}, align=center}
  ]
    \draw[timeline, line width=1.2pt] (0,0) -- (11.8,0);
    \foreach \x/\year in {0/1644,2.0/1645,4.0/1661,5.8/1673,7.7/1683,9.8/1691,11.8/1720}
      {\draw[timeline] (\x,.12) -- (\x,-.12); \node[date, below] at (\x,-.18) {\year};}
    \node[event, above] at (0, .42) {入关、北京\\政权建立};
    \node[event, below] at (2.0,-.62) {南京失守；\\南方多中心延续};
    \node[event, above] at (4.0,.42) {顺治帝去世；\\康熙即位};
    \node[event, below] at (5.8,-.62) {撤藩与\\三藩战争};
    \node[event, above] at (7.7,.42) {澎湖战事；\\台湾接收程序};
    \node[event, below] at (9.8,-.62) {多伦会盟；\\北方关系重组};
    \node[event, above] at (11.8,.42) {拉萨危机后的\\军行与制度调整};
    \node[note] at (5.9,-1.35)
      {早期清朝不是一条单线：南方战事、海峡接收、蒙古关系与高原军政的节奏并不相同。};
  \end{tikzpicture}
  \caption{早期清朝的编年定位（1644--1720，简表）。年点只标出本章讨论的政权、战争、海疆与边地节点，不是完整年表，也不把接收、驻防、会盟或制度调整视为立即完成的区域控制。日期与程序主要据《清世祖实录》《清圣祖实录》、内阁与军政档；南明、郑氏、蒙古、藏文及海洋材料须与清廷编年互校。}
  \label{fig:early-qing-timeline}
\end{figure}

\section{章节罗盘：从入关胜利到可治理的帝国}
\label{sec:early-qing-compass}

本章不把1644年写成一个已经完成的统一时刻。山海关、北京与南方战场之间，清廷必须把八旗的军事突破转化为能够收粮、发令、任官和维持秩序的朝廷；而降官、地方士绅、驻防军与居民进入这一秩序的条件并不相同。征服的速度、暴力与行政接收将并列呈现，而不会被一个“定鼎”概念遮蔽。

叙事随后转向摄政、亲政、撤藩、海疆和北部、西部边地。这里的关键问题不是中央是否发布命令，而是兵饷、迁徙、册封、设治和当地中介如何使命令在不同地区取得有限效力。本章结尾将把清初的整合看作一份持续支出的政治账：它提供了盛世帝国可以扩张的框架，也留下了驻军、身份和边疆供给的长期成本。


\section{编年主线：1644—1722}

本章从关隘、京城、南方战场、海峡和边地的相继变化推进；同一年的命令与战报不会被当作各地同时完成的政治转向。

\begin{event}[山海关：打开道路，不等于完成征服（1644）]
\eventanchor{EQ-1644-SHANHAIGUAN}{清军与吴三桂在山海关击败李自成军（1644）}，使八旗取得进入华北的通道。清方军报、关宁军务档和地方记事能够确认会战与随后的行军，却不能把求援文书原貌、各军兵额或每一次战场转折写成已经没有争议的事实。
\end{event}

\subsection*{征服展开为道路、官署与供给}

关隘一开，首先需要回答的并不是谁已经拥有“天下”，而是谁能让队伍穿过下一座城门，谁能为驻军找到粮食，谁又肯替新来的权力抄写、转递和执行文书。北京的接收因而是山海关战事的后续难题，而非其自然结局。\eventref{EQ-1644-BEIJING-ENTRY}{多尔衮进入北京}之后，仓储、官署、漕粮、旧官和惊惶未定的居民都在同一座城市里；摄政政权可以宣布招纳、任命与征发，却仍须借助明代留下的行政技艺和地方中介，才可能把军队的占有变成持续的调度。顺治帝驻跸北京和此后的亲政，固然使皇帝、部院与礼制中心逐渐会合；盛京的旗务与王朝记忆、前线不断索取的兵饷和人手，却使这个中心始终有两种方向。后来\eventref{EQ-1661-KANGXI-ACCESSION}{玄烨以幼帝即位}，辅政安排和皇帝重新掌握人事，也应放回这条尚未平定的军情与供给线来读，而不能只当作宫廷权力的短剧。

南方的时间尤其不服从北京所写的整齐节拍。南京失守后，弘光政权的崩解改变了长江下游的枢纽，却没有使浙东、福建、两广和西南在同一刻失去行动者。流动的南明朝廷、仍有粮道和旧部的将领、守城者与逃散者，分别在水道、山路、县城和港口重新估量风险。清廷的军报能够记录朝廷收到的捷报、俘获和任命；它不能代替一座城池易手后乡村如何纳粮、家族如何逃避征发，或降附者是否仍能保护旧有网络。发式服制的命令把政治归属置于身体可见之处，恰说明新政权需要反复辨认和迫使服从；它不是各地已经服从的证明。直到永历朝廷覆亡、夔东据点被攻破，战争的政治名义才逐渐收束，而道路、驻军、流民与税源的重建仍留在胜利叙事之外。

海峡让这份未完成性更清楚。郑氏的船队、岛屿据点和贸易关系，使福建沿岸的粮食、消息和人力不可能只由陆上府县来计算。迁界命令意在切断郑氏与岸上的联系，给渔盐、耕作、亲族与港口生计带来的中断却不会因谕旨而呈现为一种单一结果；复界和海禁的调整同样需要逐处察看。\eventref{EQ-1661-ZHENG-TAIWAN-LANDING}{郑成功登陆台湾}以后，围城、航路与补给把台湾推入新的军事组合；\eventref{EQ-1663-XIAMEN-JINMEN}{厦门、金门易手}又迫使郑氏重估大陆据点与岛内土地、港口之间的轻重。\eventref{EQ-1683-PENGHU}{澎湖战事}和随后\eventref{EQ-1684-TAIWAN-PREFECTURE}{设府}，分别打开了海面上的军事通道与福建行政链的延伸，但渡海船只、驻防粮饷、移民土地关系及原住民社群同官府的接触，并不会随着投降或设府同时改写。所谓接收，在海上首先是一连串需要维持的往返。

北方、东北与高原也不宜被当作征服后才出现的边缘。\eventref{EQ-1688-GALDAN-KHALKHA}{喀尔喀诸部在噶尔丹进攻下南徙求援}时，水草、牲畜、台站和安置已把清廷卷入草原的选择；\eventref{EQ-1691-DOLONNOR}{多伦诺尔会盟}所提供的是保护、封爵和盟旗关系的制度形式，不能抹去牧地使用和部际协商的后续。阿尔巴津围困后的\eventref{EQ-1689-NERCHINSK}{尼布楚谈判}，把河名、地名和翻译写成条约文字，解决的是当时军事僵持中的一项关系，而不是画出一条立刻被两岸居民共同遵守的现代线。入藏道路的险阻在一七一八年的覆败中显露无遗；\eventref{EQ-1720-LHASA-RESTORATION}{一七二〇年进入拉萨与迎立第七世达赖}，则依赖青海蒙古盟军、驮运、向导、沿途补给和拉萨当地多重关系。边地不是中央命令抵达得较晚的空白，而是让命令必须经过更多行动者、语言与交通条件才能取得有限效力的地方。

本章因而把诏令、军报、条约和任命与地方后果分开叙述。官书可以确证朝廷曾怎样命名敌我、部署军队、册封关系或规定制度；一份地方志、城守记、契约或跨境记录则往往只照亮一个城镇、港口、寺院、牧地或家庭的局部处境。它们不应彼此互相注销，也不能被拼成无所不知的全景。特别是兵数、伤亡、迁徙规模、航路细节和“归附”的范围，常由交战一方、行政机关或后出的编纂者带着各自目的保存下来。下列编年所追踪的，正是命令如何沿着关口、河道、驿路、草原和海峡移动，又在何处因粮秣、天气、翻译、地方选择和材料的沉默而改变了形状。

\subsection{1644—1652：越过关隘，南方仍在移动}

顺治元年四月，山海关把大顺、关宁军和清摄政政权的危机压在同一处关隘。李自成进入北京后，关宁军与大顺在兵饷、统属和互信上已经破裂；多尔衮利用吴三桂的处境，使八旗获得越关的行动空间。会战打开的是华北道路，并不是天下已经归属。五月入北京后，清军接收的首先是仓储、官署、粮道和观望中的官员；九月顺治帝自盛京抵京，皇帝驻跸、部院和礼制中心才开始汇合。冬季的安民、招纳、征发和整顿命令，显示一支行军军队正在试图成为能收发文书的朝廷。

这条线很快在江南折断成多条战线。1645年南京失守、弘光帝被俘，清廷又重申薙发易服；命令把发式当作政治归属的可见标记，却在扬州、江阴等地同守城、逃亡、征发和杀戮纠缠。城陷、围城和严重暴力是征服史不能略去的事实；《扬州十日记》《江阴城守纪》、地方笔记与方志所保存的日数、伤亡和见闻，则不能被改写成全境的精确统计。弘光的覆亡也没有使南方同时结束：鲁王集团、隆武、绍武和永历的政治中心仍在浙闽、两广和西南移动。

1646年至1650年，清军的攻取、南明朝廷的流转和将领的反复选择交织在一起。隆武、绍武政权先后覆亡，永历名义在两广和西南继续被不同军队借用；1648年金声桓、李成栋改奉南明，次年南昌再失、1650年广州再陷，又说明先前的“接收”并未冻结地方的选择。城池、道路、饷源和港口并不以同一种速度改易归属。清方军报能够钉住其所报的攻守与处分；南明奏疏、地方记事和城守记录保留的是败亡朝廷与城市共同体的另一种时间，二者都不能代替整个战区。

顺治七年多尔衮去世，次年顺治帝亲政，摄政枢纽随之改组；皇帝名义的加强却仍要经过宗室、旗务、满汉官僚和前线军饷才能抵达地方。1652年李定国在桂林、衡州一带的战事迫使清廷重新估量湘桂道路与机动兵力，西南的胜负并未随北京的亲政而预先决定。到这一年，读者所能看见的不是一条完成的疆界，而是仪式、降附、逃散和再选择不断彼此牵动。

\paragraph{材料位置。} 《清世祖实录》、内国史院满文档案和题本记录新政权的命令、接收与军务；南明奏疏、《永历实录》及地方城守材料提供对手和地方的时间线。它们可以相互校正事件顺序，却不能让求援文书原貌、各军兵额、攻城伤亡或每一处执行情形获得并不存在的精确性。

\section{从关门到京师}

顺治元年四月，山海关不是一扇一经开启便可直达天下的门，而是一处把几种危机压在一起的关隘。李自成入北京之后，关宁军同大顺之间的兵饷、统属与互信已经破裂；多尔衮则把这一裂缝变成八旗越过关隘的机会。清军与吴三桂部在山海关击败大顺军，因而取得进入华北的通道。这一结果可以立为编年事实；但若把后来流行的“乞师”故事逐句当作战前文件，把各军人数、会师次序和战场上的每一次转折都写成定论，证据便承受不住这样的重压。

《清世祖实录》中的军报把胜利纳入摄政政权的行动叙述，清初内国史院保存的满文关宁军务档案则让人看见命令、呈报和接收如何在文书链中移动。二者首先是清廷留下的行政记录：它们能证明朝廷收到什么消息、以何种语句处分军政，却不能自动代替关内各军的全部经验。后出的《清史稿》适合用来互校官名和大事顺序，但它是民国初年未定稿的官修史书，不是顺治四月的现场记录。吴三桂求援文书的原貌、各方兵额以及战场过程，仍须保留为由清廷叙事和后出文本共同塑形的问题。

五月入北京，也不应被压缩成“占领京城”四字。大顺撤离、明廷解体之后，城中留下的是仓储、官署、粮饷、逃散的人群和等待选择的官员。多尔衮率军入城，接收的是明朝京师的行政与礼制空间；到冬季，摄政王又以安民、征发、招纳明官和整顿官署的命令，尝试把军队占有的城池变成可以传递文书的治所。诏令的发布是可靠的，题本中可见的交接也证明部分官署重新运转；可是受任者是否到任、一纸命令是否越过城门和驿路、地方是否照办，不能从诏书的语气倒推出去。

北京地方记事和《甲申传信录》保存了另一种尺度：它们所见的是城内消息、秩序与恐惧，而非摄政政权所要展示的承接。它们同样有作者位置、传抄与见闻范围的限制，却正可防止以中央档案的沉默替代城市生活。顺治帝九月自盛京抵达北京、举行入城和定都仪式，标志着皇帝驻跸与部院中心向北京集中；盛京却仍保有满洲政治和象征上的重量。仪式建立的是持续的中央名义，并不等于各地已经承认，尤其不等于南方战场已经结束。

\section{摄政并非一条畅通的命令线}

多尔衮的摄政，是清初征服得以调动兵马、官员和名义资源的枢纽，但不是一架无阻运转的机器。京师需要粮饷，部院需要熟悉旧制的书吏和降官，前线又不断索要兵丁、漕运和军需。清廷档案最擅长呈现这条由奏报、批示、部议到任命的路径；地方档案、方志和南方对手的材料则不断提醒人们，路径在何处中断、被拖延，或被战事改写。把“招降”写成一个完成状态，会遮蔽降官仍须在家族、旧部、军饷和地方安全之间反复选择的处境。

顺治七年十二月，多尔衮在古北口外行猎途中去世。其死期与丧仪可以由《清世祖实录》和内国史院的遗命、丧仪档相互钉住；死因以及死后政治清算的动机，却不能只照后来定性的语言复述。次年顺治帝亲政，摄政安排撤除，皇帝周边的任官与决策网络重新组合。然而，署在皇帝名下的命令仍要依靠宗室、旗务、满汉官僚与南方将领落实。摄政结束，不是中央能力忽然圆满的时刻。

这一点在顺治帝去世后更为清楚。顺治十八年，八岁的玄烨即位，索尼、苏克萨哈、遏必隆、鳌拜辅政。继位与四辅政大臣的名单见于《清世祖实录》、起居注和遗诏档；至康熙四年鳌拜影响上升、康熙六年索尼去世、康熙八年鳌拜被拘捕，则应分别阅读《清圣祖实录》、起居注与内阁案卷。尤其不能把1669年的结果倒灌回此前：后来的鳌拜案和人物传记最容易把本来由议政、旗务、人事和军费共同构成的辅政政治，写成一人独揽权力的必然前史。康熙拘捕鳌拜使皇帝能够更直接地统摄人事军政，却也仍须在既有制度与旗人贵族网络中行使这种权力。

\section{南京失守以后：城池、朝廷与身体}

清军南下的速度使京师接收和江南战争看似连成一线，实际却是一系列由水陆节点、粮道和地方选择构成的战区。顺治二年春夏之交，清军沿运河、淮扬推进，攻取南京；弘光帝朱由崧不久在芜湖一带被俘。清方的《清世祖实录》与《清史稿·世祖本纪》可以确立城陷、押解这一主干，南明弘光朝奏疏、《南明史料》所收文书和江宁地方记事则保留了另一面：朝廷如何失去护卫与指挥，官员和军队怎样逃散、观望或改易归属。前者是胜利政权收到的战报，后者多是败亡朝廷及城市中的急迫文字；两类材料不能互相注销，也不应被合成为一种全知叙述。

南京失去首都地位，并没有令“明亡”在南方成为一个同时发生的事件。弘光名义瓦解后，鲁王集团在浙东和海上活动，隆武朝廷在福建建立，绍武与永历又在两广相继争夺名义中心。清军取得杭州、绍兴等城，或攻入广州，并不自动说明一省的道路、岛屿、山地和县城都已归于同一控制。清廷军报常以府城、俘获和招降排列战果；《福建通志》《广东通志》及各地府县志可以补足部分地名和后续秩序，却大多成书较晚、也会沿袭胜败叙事。南明文集和《永历实录》把流动朝廷、将领联络与退守路线连成另一条线索，但它们也服务于保存正统、追责或追忆的需要。史实的难处不在于挑选哪一方“说真话”，而在于说明每种文本首先能够证明什么。

顺治二年六月重申的薙发易服命令，使政治归属进入身体的可见处。实录中的命令证明清廷以发式服制辨认新旧政治身份；兵科题本和江南府县志的相关条目，则可见命令在若干地方与征发、守城和治安纠缠在一起。它不能证明全国在同一日、以同一种方式服从。有人遵从，有人逃亡，有人把发式当作拒绝交付城池或粮饷的标识；不同府县的速度、强度与后果必须逐地追问。

因此，扬州、江阴和1650年的广州特别不宜用一个整齐的死亡数字盖棺。扬州失陷、江阴围城失守、广州再陷，以及攻城后发生严重杀戮、掳掠和损毁，都是必须写入征服史的事实。可是清军战报所陈的战果，不是伤亡统计；《扬州十日记》《江阴城守纪》、广州地方笔记和后出的府县志，保存了幸存者、编纂者或地方共同体的恐惧与纪念，也各有成书、传抄、修辞和数字来源的问题。“十日”、围城天数和死亡人数不应被伪装成精确普查，更不能由一部文本推算江南或广东整体的死亡。较为诚实的写法是同时承认暴力及其长期记忆，说明现有材料不能可靠合算总数，并把不同数字标回提出它的作者、地点与文类。

\section{南方并未在一役中归入清廷}

1646年隆武帝在闽北被俘杀、绍武朝在广州结束之后，朱由榔在肇庆即位，永历年号成为两广和西南多支力量可以借用的名义。清廷军报将这看作继续清剿的对象；《永历实录》、永历朝奏疏和地方志所呈现的，却是一个不断移动、依赖地方军队、饷源与山海通道的朝廷。1648年金声桓在南昌、李成栋在广东改奉南明，说明此前的“接收”未能冻结将领与地方的选择。1649年南昌被攻取、1650年广州再陷固然削弱了这些反转的中心，却不能把战后的税源、流民和港口恢复一并写成已经解决。

永历六年，李定国攻取桂林、在衡州击败清军，迫使清廷重新估量湘桂道路、军需和机动兵力。双方的兵数、死伤和将领死因在清廷军报与南明记录之间多有不同；可以肯定战场结果，不能替任何一方把报捷数字升级为无争议的计数。此后孙可望与永历朝的冲突及其降清，削弱了南明的军事协作；清军自贵州、四川等方向入滇，永历朝廷离开昆明，终至进入缅甸。这里又有一层不能省略的证据边界：清方材料擅长记录追捕与军政调度，缅甸宫廷的《琉璃宫史》传统和欧洲旅行记所保留的接待、路线与交涉，并不必然与之重合。把缅甸只写作清军追击地图上的空白，既抹掉了东吁王朝的选择，也误解了流亡朝廷的补给条件。

1659年郑成功北上围攻南京而后撤退，显示长江下游并非已无战事；翌年清军经营云南城镇和驿路，也说明占有昆明不等于控制边地。1662年，朱由榔由缅甸被押回后处死，南明皇帝的名义中心至此终结。中缅材料对于移交过程、处死地点和责任分配仍须并读。夔东抗清据点直到1664年才被清军攻破，郑氏又在海上和台湾维持政权；这些后续并非枝节，而是对“1644年入关即全国统一”这一方便说法的校正。

从山海关到北京，从南京到西南边地，清初形成的不是一条单向延伸的疆界，而是军队、文书、仪式、降附、逃亡与再选择彼此牵动的过程。清廷档案使我们能够看见命令如何被制作和宣告；官修史书为后来者编排出帝国的时间；地方材料留下城池和社区承受的断片；南明叙事则保存失败者仍在行动的政治地理。把它们放在同一页上，不是为了把分歧磨平，而是为了让征服、摄政与整合都回到它们当时尚未完成的样子。


\subsection{事件链与材料限度}
关隘、京城与官署的材料并不构成一条无摩擦的接收线。以下事件保留各自的行动者、空间、后果与证据限度。

\paragraph{顺治元年五月（1644年6月前后）：多尔衮率军进入北京并接收明朝京城的行政与礼制空间}
\eventanchor{EQ-1644-BEIJING-ENTRY}{多尔衮率军进入北京并接收明朝京城的行政与礼制空间}。本节点叙述该事件本身。清摄政政权与北京在京师与多地军政线路相接的尺度的处境首先受大顺撤离与明廷崩溃留下军政真空，清军需获取京师粮饷、官僚和象征中心制约。《清世祖实录》顺治元年五月；《清史稿·世祖本纪》；《甲申传信录》及北京地方记事留下可核的线索；可辨的后果是清廷以北京为治所，开始将军事占领转化为对明朝制度资源的继承和重组。原有置信注记：入城与接收可靠；城内秩序、迎降者范围和掠夺状况须以不同见证互校

\paragraph{顺治元年五月（1644年6月前后）：多尔衮率军进入北京并接收明朝京城的行政与礼制空间}
\eventanchor{EQ-1644-BEIJING-ENTRY-AFTERMATH}{多尔衮率军进入北京并接收明朝京城的行政与礼制空间} 置于京师与多地军政线路相接的尺度中阅读，本节点追踪「多尔衮率军进入北京并接收明朝京城的行政与礼制空间」之后可见的余波。大顺撤离与明廷崩溃留下军政真空，清军需获取京师粮饷、官僚和象征中心构成行动者清摄政政权与北京必须面对的条件。取证从《清世祖实录》顺治元年五月；《清史稿·世祖本纪》；《甲申传信录》及北京地方记事出发，因而只能把变化谨慎地连到清廷以北京为治所，开始将军事占领转化为对明朝制度资源的继承和重组。原有置信注记：入城与接收可靠；城内秩序、迎降者范围和掠夺状况须以不同见证互校

\paragraph{顺治元年五月（1644年6月前后）：多尔衮率军进入北京并接收明朝京城的行政与礼制空间}
\eventanchor{EQ-1644-BEIJING-ENTRY-DECISION}{多尔衮率军进入北京并接收明朝京城的行政与礼制空间} 标记本节点定位围绕「多尔衮率军进入北京并接收明朝京城的行政与礼制空间」形成的处置与调度记录，涉及清摄政政权与北京。京师与多地军政线路相接的尺度不是抽象背景：大顺撤离与明廷崩溃留下军政真空，清军需获取京师粮饷、官僚和象征中心。《清世祖实录》顺治元年五月；《清史稿·世祖本纪》；《甲申传信录》及北京地方记事所见的顺序随后指向清廷以北京为治所，开始将军事占领转化为对明朝制度资源的继承和重组，但原有置信注记：入城与接收可靠；城内秩序、迎降者范围和掠夺状况须以不同见证互校

\paragraph{顺治元年五月（1644年6月前后）：多尔衮率军进入北京并接收明朝京城的行政与礼制空间}
在京师与多地军政线路相接的尺度，\eventanchor{EQ-1644-BEIJING-ENTRY-PRELUDE}{多尔衮率军进入北京并接收明朝京城的行政与礼制空间} 让清摄政政权与北京的一个选择或压力有了可查的坐标。本节点回看「多尔衮率军进入北京并接收明朝京城的行政与礼制空间」之前已可见的条件。前因可写为大顺撤离与明廷崩溃留下军政真空，清军需获取京师粮饷、官僚和象征中心；《清世祖实录》顺治元年五月；《清史稿·世祖本纪》；《甲申传信录》及北京地方记事限定了记录，后续变化是清廷以北京为治所，开始将军事占领转化为对明朝制度资源的继承和重组。原有置信注记：入城与接收可靠；城内秩序、迎降者范围和掠夺状况须以不同见证互校

\paragraph{顺治元年冬（1644年）：多尔衮以摄政王名义整顿京师官署、招降明官并发布安民与征发命令}
\eventanchor{EQ-1644-REGENT-ORDERS}{多尔衮以摄政王名义整顿京师官署、招降明官并发布安民与征发命令} 所关联的是清摄政政权在京师与多地军政线路相接的尺度中的行动。本节点叙述该事件本身。它从新政权需要恢复行政传递、维持京师粮饷并吸纳明代文官技术展开，借《清世祖实录》顺治元年十月至十二月；《清史稿·职官志》；中国第一历史档案馆藏内阁大库顺治元年题本〔京师官署交接〕进入可检验的年序，并以中央可用的部院和任官网络开始恢复，但忠诚、财力与军纪继续限制覆盖面影响下一段安排。原有置信注记：命令发布可靠；接受任官和地方执行不能从诏令直接推出

\paragraph{顺治元年冬（1644年）：多尔衮以摄政王名义整顿京师官署、招降明官并发布安民与征发命令}
京师与多地军政线路相接的尺度中的\eventanchor{EQ-1644-REGENT-ORDERS-AFTERMATH}{多尔衮以摄政王名义整顿京师官署、招降明官并发布安民与征发命令} 不能离开清摄政政权来解释。本节点追踪「多尔衮以摄政王名义整顿京师官署、招降明官并发布安民与征发命令」之后可见的余波。新政权需要恢复行政传递、维持京师粮饷并吸纳明代文官技术说明其形成的条件；《清世祖实录》顺治元年十月至十二月；《清史稿·职官志》；中国第一历史档案馆藏内阁大库顺治元年题本〔京师官署交接〕提供来源路径，中央可用的部院和任官网络开始恢复，但忠诚、财力与军纪继续限制覆盖面则是材料能够支持的近时结果。原有置信注记：命令发布可靠；接受任官和地方执行不能从诏令直接推出

\paragraph{顺治元年冬（1644年）：多尔衮以摄政王名义整顿京师官署、招降明官并发布安民与征发命令}
\eventanchor{EQ-1644-REGENT-ORDERS-DECISION}{多尔衮以摄政王名义整顿京师官署、招降明官并发布安民与征发命令} 把清摄政政权放回京师与多地军政线路相接的尺度的道路、城镇、边界或制度网络。本节点定位围绕「多尔衮以摄政王名义整顿京师官署、招降明官并发布安民与征发命令」形成的处置与调度记录。新政权需要恢复行政传递、维持京师粮饷并吸纳明代文官技术。本段采用《清世祖实录》顺治元年十月至十二月；《清史稿·职官志》；中国第一历史档案馆藏内阁大库顺治元年题本〔京师官署交接〕核对其顺序，随后可追到中央可用的部院和任官网络开始恢复，但忠诚、财力与军纪继续限制覆盖面。原有置信注记：命令发布可靠；接受任官和地方执行不能从诏令直接推出

\paragraph{顺治元年冬（1644年）：多尔衮以摄政王名义整顿京师官署、招降明官并发布安民与征发命令}
本节点回看「多尔衮以摄政王名义整顿京师官署、招降明官并发布安民与征发命令」之前已可见的条件，故\eventanchor{EQ-1644-REGENT-ORDERS-PRELUDE}{多尔衮以摄政王名义整顿京师官署、招降明官并发布安民与征发命令} 不只是一个年份标签。清摄政政权在京师与多地军政线路相接的尺度承受的压力来自新政权需要恢复行政传递、维持京师粮饷并吸纳明代文官技术。《清世祖实录》顺治元年十月至十二月；《清史稿·职官志》；中国第一历史档案馆藏内阁大库顺治元年题本〔京师官署交接〕可回查该节点；中央可用的部院和任官网络开始恢复，但忠诚、财力与军纪继续限制覆盖面是其进入后续年序的方式。原有置信注记：命令发布可靠；接受任官和地方执行不能从诏令直接推出

\paragraph{顺治元年四月（1644年5月，换算待核）：清军与吴三桂在山海关击败李自成军，打开进入华北的军事通道}
\eventanchor{EQ-1644-SHANHAIGUAN-AFTERMATH}{清军与吴三桂在山海关击败李自成军，打开进入华北的军事通道} 所见的行动者为清摄政政权、吴三桂与大顺，空间尺度为京师与多地军政线路相接的尺度。本节点追踪「清军与吴三桂在山海关击败李自成军，打开进入华北的军事通道」之后可见的余波。李自成入北京后关宁军与大顺在兵饷、统属和互信上破裂，多尔衮把危机转为入关机会使它成为具体事件链的一环；《清世祖实录》顺治元年四月山海关军报；《清史稿·世祖本纪》；《清初内国史院满文档案》顺治元年四月关宁军务档与清军获得穿越关隘的行动空间，北京政权争夺转为华北与南方的多战区征服。原有置信注记：核心战事可靠；吴三桂求援文书原貌、各军兵数和战场过程受清廷与后出叙事塑形

\paragraph{顺治元年四月（1644年5月，换算待核）：清军与吴三桂在山海关击败李自成军，打开进入华北的军事通道}
从京师与多地军政线路相接的尺度看，\eventanchor{EQ-1644-SHANHAIGUAN-DECISION}{清军与吴三桂在山海关击败李自成军，打开进入华北的军事通道} 留下本节点定位围绕「清军与吴三桂在山海关击败李自成军，打开进入华北的军事通道」形成的处置与调度记录的材料坐标。清摄政政权、吴三桂与大顺所面对的条件是李自成入北京后关宁军与大顺在兵饷、统属和互信上破裂，多尔衮把危机转为入关机会。《清世祖实录》顺治元年四月山海关军报；《清史稿·世祖本纪》；《清初内国史院满文档案》顺治元年四月关宁军务档支持的结果为清军获得穿越关隘的行动空间，北京政权争夺转为华北与南方的多战区征服。原有置信注记：核心战事可靠；吴三桂求援文书原貌、各军兵数和战场过程受清廷与后出叙事塑形

\paragraph{顺治元年四月（1644年5月，换算待核）：清军与吴三桂在山海关击败李自成军，打开进入华北的军事通道}
\eventanchor{EQ-1644-SHANHAIGUAN-PRELUDE}{清军与吴三桂在山海关击败李自成军，打开进入华北的军事通道} 记录清摄政政权、吴三桂与大顺在京师与多地军政线路相接的尺度的一次变化。本节点回看「清军与吴三桂在山海关击败李自成军，打开进入华北的军事通道」之前已可见的条件。其发生与李自成入北京后关宁军与大顺在兵饷、统属和互信上破裂，多尔衮把危机转为入关机会相连；《清世祖实录》顺治元年四月山海关军报；《清史稿·世祖本纪》；《清初内国史院满文档案》顺治元年四月关宁军务档把事件系回来源，清军获得穿越关隘的行动空间，北京政权争夺转为华北与南方的多战区征服显示压力如何向后转移。原有置信注记：核心战事可靠；吴三桂求援文书原貌、各军兵数和战场过程受清廷与后出叙事塑形

\paragraph{顺治元年九月（1644年10月前后）：顺治帝自盛京抵北京，清廷举行定都与皇帝入城仪式}
\eventanchor{EQ-1644-SHUNZHI-BEIJING}{顺治帝自盛京抵北京，清廷举行定都与皇帝入城仪式} 的地点与行动范围属于京师与多地军政线路相接的尺度，行动者是清。本节点叙述该事件本身。摄政政权虽占北京，仍需皇帝亲临与礼制重置来建立持续中央名义。读《清世祖实录》顺治元年九月；《清史稿·世祖本纪》；《清初内国史院满文档案》北京迁都礼仪档时可见其记录边界，最小后果只能写作北京成为皇帝驻跸与部院中心，盛京仍保有满洲政治和象征地位。原有置信注记：日期与仪式主干可靠；官修礼仪记录不能等同各地承认

\paragraph{顺治元年九月（1644年10月前后）：顺治帝自盛京抵北京，清廷举行定都与皇帝入城仪式}
\eventanchor{EQ-1644-SHUNZHI-BEIJING-AFTERMATH}{顺治帝自盛京抵北京，清廷举行定都与皇帝入城仪式} 是清在京师与多地军政线路相接的尺度留下的编年标记。本节点追踪「顺治帝自盛京抵北京，清廷举行定都与皇帝入城仪式」之后可见的余波。它从摄政政权虽占北京，仍需皇帝亲临与礼制重置来建立持续中央名义走来；《清世祖实录》顺治元年九月；《清史稿·世祖本纪》；《清初内国史院满文档案》北京迁都礼仪档固定可证部分，北京成为皇帝驻跸与部院中心，盛京仍保有满洲政治和象征地位构成下一步的可见结果。原有置信注记：日期与仪式主干可靠；官修礼仪记录不能等同各地承认

\paragraph{顺治元年九月（1644年10月前后）：顺治帝自盛京抵北京，清廷举行定都与皇帝入城仪式}
本节点定位围绕「顺治帝自盛京抵北京，清廷举行定都与皇帝入城仪式」形成的处置与调度记录：\eventanchor{EQ-1644-SHUNZHI-BEIJING-DECISION}{顺治帝自盛京抵北京，清廷举行定都与皇帝入城仪式}。清在京师与多地军政线路相接的尺度的处境可由摄政政权虽占北京，仍需皇帝亲临与礼制重置来建立持续中央名义说明。《清世祖实录》顺治元年九月；《清史稿·世祖本纪》；《清初内国史院满文档案》北京迁都礼仪档是本段材料入口，因而只能据此追踪到北京成为皇帝驻跸与部院中心，盛京仍保有满洲政治和象征地位。原有置信注记：日期与仪式主干可靠；官修礼仪记录不能等同各地承认

\paragraph{顺治元年九月（1644年10月前后）：顺治帝自盛京抵北京，清廷举行定都与皇帝入城仪式}
\eventanchor{EQ-1644-SHUNZHI-BEIJING-PRELUDE}{顺治帝自盛京抵北京，清廷举行定都与皇帝入城仪式} 发生在京师与多地军政线路相接的尺度，牵动清。本节点回看「顺治帝自盛京抵北京，清廷举行定都与皇帝入城仪式」之前已可见的条件。摄政政权虽占北京，仍需皇帝亲临与礼制重置来建立持续中央名义是其结构性条件；《清世祖实录》顺治元年九月；《清史稿·世祖本纪》；《清初内国史院满文档案》北京迁都礼仪档留下的证据把读者带到北京成为皇帝驻跸与部院中心，盛京仍保有满洲政治和象征地位。原有置信注记：日期与仪式主干可靠；官修礼仪记录不能等同各地承认



\subsection{事件链与材料限度}
南方战场的城池、江路与地方社会没有同速改变。下列节点保存征服、迁徙与接收之间的材料距离。

\paragraph{顺治二年六月（1645年）：清廷重申薙发易服命令，在江南等地引发服从、逃亡与抵抗}
\eventanchor{EQ-1645-HAIR-ORDER}{清廷重申薙发易服命令，在江南等地引发服从、逃亡与抵抗}。本节点叙述该事件本身。清与江南各地在省、府县与地方社会相互牵动的尺度的处境首先受清廷以发式服制识别新旧政治归属，军事推进又使命令与治安相互捆绑制约。《清世祖实录》顺治二年六月；《清史稿·舆服志》；内阁大库顺治二年兵科题本与江南府县志薙发条留下可核的线索；可辨的后果是服制成为忠降识别和集体动员的触发点，若干江南城镇的冲突加剧征服成本。原有置信注记：命令文本可靠；各地执行速度和社会反应高度不一，不能以诏书概括全国

\paragraph{顺治二年六月（1645年）：清廷重申薙发易服命令，在江南等地引发服从、逃亡与抵抗}
\eventanchor{EQ-1645-HAIR-ORDER-AFTERMATH}{清廷重申薙发易服命令，在江南等地引发服从、逃亡与抵抗} 置于省、府县与地方社会相互牵动的尺度中阅读，本节点追踪「清廷重申薙发易服命令，在江南等地引发服从、逃亡与抵抗」之后可见的余波。清廷以发式服制识别新旧政治归属，军事推进又使命令与治安相互捆绑构成行动者清与江南各地必须面对的条件。取证从《清世祖实录》顺治二年六月；《清史稿·舆服志》；内阁大库顺治二年兵科题本与江南府县志薙发条出发，因而只能把变化谨慎地连到服制成为忠降识别和集体动员的触发点，若干江南城镇的冲突加剧征服成本。原有置信注记：命令文本可靠；各地执行速度和社会反应高度不一，不能以诏书概括全国

\paragraph{顺治二年六月（1645年）：清廷重申薙发易服命令，在江南等地引发服从、逃亡与抵抗}
\eventanchor{EQ-1645-HAIR-ORDER-DECISION}{清廷重申薙发易服命令，在江南等地引发服从、逃亡与抵抗} 标记本节点定位围绕「清廷重申薙发易服命令，在江南等地引发服从、逃亡与抵抗」形成的处置与调度记录，涉及清与江南各地。省、府县与地方社会相互牵动的尺度不是抽象背景：清廷以发式服制识别新旧政治归属，军事推进又使命令与治安相互捆绑。《清世祖实录》顺治二年六月；《清史稿·舆服志》；内阁大库顺治二年兵科题本与江南府县志薙发条所见的顺序随后指向服制成为忠降识别和集体动员的触发点，若干江南城镇的冲突加剧征服成本，但原有置信注记：命令文本可靠；各地执行速度和社会反应高度不一，不能以诏书概括全国

\paragraph{顺治二年六月（1645年）：清廷重申薙发易服命令，在江南等地引发服从、逃亡与抵抗}
在省、府县与地方社会相互牵动的尺度，\eventanchor{EQ-1645-HAIR-ORDER-PRELUDE}{清廷重申薙发易服命令，在江南等地引发服从、逃亡与抵抗} 让清与江南各地的一个选择或压力有了可查的坐标。本节点回看「清廷重申薙发易服命令，在江南等地引发服从、逃亡与抵抗」之前已可见的条件。前因可写为清廷以发式服制识别新旧政治归属，军事推进又使命令与治安相互捆绑；《清世祖实录》顺治二年六月；《清史稿·舆服志》；内阁大库顺治二年兵科题本与江南府县志薙发条限定了记录，后续变化是服制成为忠降识别和集体动员的触发点，若干江南城镇的冲突加剧征服成本。原有置信注记：命令文本可靠；各地执行速度和社会反应高度不一，不能以诏书概括全国

\paragraph{顺治二年五月（1645年）：弘光帝朱由崧在芜湖一带被俘并被押往北京}
\eventanchor{EQ-1645-HONGGUANG-CAPTURE}{弘光帝朱由崧在芜湖一带被俘并被押往北京} 所关联的是清与南明弘光政权在京师与多地军政线路相接的尺度中的行动。本节点叙述该事件本身。它从南京失守后弘光朝没有持续的军事护卫与统一指挥，逃亡路线被清军和地方武装切断展开，借《清世祖实录》顺治二年五月；《清史稿·弘光帝本纪》；《明季南略》及江南地方志相关条目进入可检验的年序，并以以弘光名义维系的朝廷失能，隆武、永历政权在东南和西南重建名义中心影响下一段安排。原有置信注记：被俘及押解主干可靠；参与者责任和移交地点在纪传、地方记事间有异

\paragraph{顺治二年五月（1645年）：弘光帝朱由崧在芜湖一带被俘并被押往北京}
京师与多地军政线路相接的尺度中的\eventanchor{EQ-1645-HONGGUANG-CAPTURE-AFTERMATH}{弘光帝朱由崧在芜湖一带被俘并被押往北京} 不能离开清与南明弘光政权来解释。本节点追踪「弘光帝朱由崧在芜湖一带被俘并被押往北京」之后可见的余波。南京失守后弘光朝没有持续的军事护卫与统一指挥，逃亡路线被清军和地方武装切断说明其形成的条件；《清世祖实录》顺治二年五月；《清史稿·弘光帝本纪》；《明季南略》及江南地方志相关条目提供来源路径，以弘光名义维系的朝廷失能，隆武、永历政权在东南和西南重建名义中心则是材料能够支持的近时结果。原有置信注记：被俘及押解主干可靠；参与者责任和移交地点在纪传、地方记事间有异

\paragraph{顺治二年五月（1645年）：弘光帝朱由崧在芜湖一带被俘并被押往北京}
\eventanchor{EQ-1645-HONGGUANG-CAPTURE-DECISION}{弘光帝朱由崧在芜湖一带被俘并被押往北京} 把清与南明弘光政权放回京师与多地军政线路相接的尺度的道路、城镇、边界或制度网络。本节点定位围绕「弘光帝朱由崧在芜湖一带被俘并被押往北京」形成的处置与调度记录。南京失守后弘光朝没有持续的军事护卫与统一指挥，逃亡路线被清军和地方武装切断。本段采用《清世祖实录》顺治二年五月；《清史稿·弘光帝本纪》；《明季南略》及江南地方志相关条目核对其顺序，随后可追到以弘光名义维系的朝廷失能，隆武、永历政权在东南和西南重建名义中心。原有置信注记：被俘及押解主干可靠；参与者责任和移交地点在纪传、地方记事间有异

\paragraph{顺治二年五月（1645年）：弘光帝朱由崧在芜湖一带被俘并被押往北京}
本节点回看「弘光帝朱由崧在芜湖一带被俘并被押往北京」之前已可见的条件，故\eventanchor{EQ-1645-HONGGUANG-CAPTURE-PRELUDE}{弘光帝朱由崧在芜湖一带被俘并被押往北京} 不只是一个年份标签。清与南明弘光政权在京师与多地军政线路相接的尺度承受的压力来自南京失守后弘光朝没有持续的军事护卫与统一指挥，逃亡路线被清军和地方武装切断。《清世祖实录》顺治二年五月；《清史稿·弘光帝本纪》；《明季南略》及江南地方志相关条目可回查该节点；以弘光名义维系的朝廷失能，隆武、永历政权在东南和西南重建名义中心是其进入后续年序的方式。原有置信注记：被俘及押解主干可靠；参与者责任和移交地点在纪传、地方记事间有异

\paragraph{顺治二年闰六月至八月（1645年）：江阴在薙发命令和清军围攻中失守，地方记忆长期以守城与屠城叙述该事件}
\eventanchor{EQ-1645-JIANGYIN}{江阴在薙发命令和清军围攻中失守，地方记忆长期以守城与屠城叙述该事件} 所见的行动者为清军与江阴地方守军，空间尺度为省、府县与地方社会相互牵动的尺度。本节点叙述该事件本身。服制命令与守备危机把行政服从转成军事对抗，清军需确保沿江交通和粮道使它成为具体事件链的一环；《清世祖实录》顺治二年江南军报；《清史稿·世祖本纪》；《江阴县志》及《江阴城守纪》与清廷以军事压制恢复沿江节点；江阴成为地方忠义记忆，伤亡规模不能由单一叙事定量。原有置信注记：城陷主干可靠；围城天数、伤亡和人物言行多依赖地方后出文本

\paragraph{顺治二年闰六月至八月（1645年）：江阴在薙发命令和清军围攻中失守，地方记忆长期以守城与屠城叙述该事件}
从省、府县与地方社会相互牵动的尺度看，\eventanchor{EQ-1645-JIANGYIN-AFTERMATH}{江阴在薙发命令和清军围攻中失守，地方记忆长期以守城与屠城叙述该事件} 留下本节点追踪「江阴在薙发命令和清军围攻中失守，地方记忆长期以守城与屠城叙述该事件」之后可见的余波的材料坐标。清军与江阴地方守军所面对的条件是服制命令与守备危机把行政服从转成军事对抗，清军需确保沿江交通和粮道。《清世祖实录》顺治二年江南军报；《清史稿·世祖本纪》；《江阴县志》及《江阴城守纪》支持的结果为清廷以军事压制恢复沿江节点；江阴成为地方忠义记忆，伤亡规模不能由单一叙事定量。原有置信注记：城陷主干可靠；围城天数、伤亡和人物言行多依赖地方后出文本

\paragraph{顺治二年闰六月至八月（1645年）：江阴在薙发命令和清军围攻中失守，地方记忆长期以守城与屠城叙述该事件}
\eventanchor{EQ-1645-JIANGYIN-DECISION}{江阴在薙发命令和清军围攻中失守，地方记忆长期以守城与屠城叙述该事件} 记录清军与江阴地方守军在省、府县与地方社会相互牵动的尺度的一次变化。本节点定位围绕「江阴在薙发命令和清军围攻中失守，地方记忆长期以守城与屠城叙述该事件」形成的处置与调度记录。其发生与服制命令与守备危机把行政服从转成军事对抗，清军需确保沿江交通和粮道相连；《清世祖实录》顺治二年江南军报；《清史稿·世祖本纪》；《江阴县志》及《江阴城守纪》把事件系回来源，清廷以军事压制恢复沿江节点；江阴成为地方忠义记忆，伤亡规模不能由单一叙事定量显示压力如何向后转移。原有置信注记：城陷主干可靠；围城天数、伤亡和人物言行多依赖地方后出文本

\paragraph{顺治二年闰六月至八月（1645年）：江阴在薙发命令和清军围攻中失守，地方记忆长期以守城与屠城叙述该事件}
\eventanchor{EQ-1645-JIANGYIN-PRELUDE}{江阴在薙发命令和清军围攻中失守，地方记忆长期以守城与屠城叙述该事件} 的地点与行动范围属于省、府县与地方社会相互牵动的尺度，行动者是清军与江阴地方守军。本节点回看「江阴在薙发命令和清军围攻中失守，地方记忆长期以守城与屠城叙述该事件」之前已可见的条件。服制命令与守备危机把行政服从转成军事对抗，清军需确保沿江交通和粮道。读《清世祖实录》顺治二年江南军报；《清史稿·世祖本纪》；《江阴县志》及《江阴城守纪》时可见其记录边界，最小后果只能写作清廷以军事压制恢复沿江节点；江阴成为地方忠义记忆，伤亡规模不能由单一叙事定量。原有置信注记：城陷主干可靠；围城天数、伤亡和人物言行多依赖地方后出文本

\paragraph{顺治二年四月（1645年5月前后）：多铎军攻取南京，南明弘光朝的都城失守}
\eventanchor{EQ-1645-NANJING-TAKEN}{多铎军攻取南京，南明弘光朝的都城失守} 是清与南明弘光政权在京师与多地军政线路相接的尺度留下的编年标记。本节点叙述该事件本身。它从清军沿运河与淮扬推进，弘光朝党争、军饷困难和江防失守相互叠加走来；《清世祖实录》顺治二年四月；《清史稿·世祖本纪》；《南明史料》所收弘光朝奏疏与江宁地方记事固定可证部分，南明失去首都和财政文化中心，江浙、闽粤的抵抗及另立政权随之分化构成下一步的可见结果。原有置信注记：城陷日期可靠；守军数量、降附动机和战后暴力须并读南明与地方材料

\paragraph{顺治二年四月（1645年5月前后）：多铎军攻取南京，南明弘光朝的都城失守}
本节点追踪「多铎军攻取南京，南明弘光朝的都城失守」之后可见的余波：\eventanchor{EQ-1645-NANJING-TAKEN-AFTERMATH}{多铎军攻取南京，南明弘光朝的都城失守}。清与南明弘光政权在京师与多地军政线路相接的尺度的处境可由清军沿运河与淮扬推进，弘光朝党争、军饷困难和江防失守相互叠加说明。《清世祖实录》顺治二年四月；《清史稿·世祖本纪》；《南明史料》所收弘光朝奏疏与江宁地方记事是本段材料入口，因而只能据此追踪到南明失去首都和财政文化中心，江浙、闽粤的抵抗及另立政权随之分化。原有置信注记：城陷日期可靠；守军数量、降附动机和战后暴力须并读南明与地方材料

\paragraph{顺治二年四月（1645年5月前后）：多铎军攻取南京，南明弘光朝的都城失守}
\eventanchor{EQ-1645-NANJING-TAKEN-DECISION}{多铎军攻取南京，南明弘光朝的都城失守} 发生在京师与多地军政线路相接的尺度，牵动清与南明弘光政权。本节点定位围绕「多铎军攻取南京，南明弘光朝的都城失守」形成的处置与调度记录。清军沿运河与淮扬推进，弘光朝党争、军饷困难和江防失守相互叠加是其结构性条件；《清世祖实录》顺治二年四月；《清史稿·世祖本纪》；《南明史料》所收弘光朝奏疏与江宁地方记事留下的证据把读者带到南明失去首都和财政文化中心，江浙、闽粤的抵抗及另立政权随之分化。原有置信注记：城陷日期可靠；守军数量、降附动机和战后暴力须并读南明与地方材料

\paragraph{顺治二年四月（1645年5月前后）：多铎军攻取南京，南明弘光朝的都城失守}
\eventanchor{EQ-1645-NANJING-TAKEN-PRELUDE}{多铎军攻取南京，南明弘光朝的都城失守} 对应本节点回看「多铎军攻取南京，南明弘光朝的都城失守」之前已可见的条件。在京师与多地军政线路相接的尺度，清与南明弘光政权无法绕开清军沿运河与淮扬推进，弘光朝党争、军饷困难和江防失守相互叠加。依据《清世祖实录》顺治二年四月；《清史稿·世祖本纪》；《南明史料》所收弘光朝奏疏与江宁地方记事，本段把后续影响限定为南明失去首都和财政文化中心，江浙、闽粤的抵抗及另立政权随之分化。原有置信注记：城陷日期可靠；守军数量、降附动机和战后暴力须并读南明与地方材料

\paragraph{顺治二年四—五月（1645年）：扬州攻陷后发生大规模杀戮与掳掠，规模和持续时间存在重大争议}
\eventanchor{EQ-1645-YANGZHOU}{扬州攻陷后发生大规模杀戮与掳掠，规模和持续时间存在重大争议}。本节点叙述该事件本身。清军与扬州在京师与多地军政线路相接的尺度的处境首先受扬州守军扼守清军南下通道，攻城后的军纪、报复和掠夺处于战时权力失控与奖惩机制中制约。《清世祖实录》顺治二年四月扬州军报；《清史稿·史可法传》；王秀楚《扬州十日记》与《扬州府志》留下可核的线索；可辨的后果是扬州经验在江南传播，既促成恐惧性降附，也强化若干城镇抵抗及后世纪念政治。原有置信注记：扬州失陷可靠；死亡数字和“十日”叙事须标为后出见证并与地方材料互证

\paragraph{顺治二年四—五月（1645年）：扬州攻陷后发生大规模杀戮与掳掠，规模和持续时间存在重大争议}
\eventanchor{EQ-1645-YANGZHOU-AFTERMATH}{扬州攻陷后发生大规模杀戮与掳掠，规模和持续时间存在重大争议} 置于京师与多地军政线路相接的尺度中阅读，本节点追踪「扬州攻陷后发生大规模杀戮与掳掠，规模和持续时间存在重大争议」之后可见的余波。扬州守军扼守清军南下通道，攻城后的军纪、报复和掠夺处于战时权力失控与奖惩机制中构成行动者清军与扬州必须面对的条件。取证从《清世祖实录》顺治二年四月扬州军报；《清史稿·史可法传》；王秀楚《扬州十日记》与《扬州府志》出发，因而只能把变化谨慎地连到扬州经验在江南传播，既促成恐惧性降附，也强化若干城镇抵抗及后世纪念政治。原有置信注记：扬州失陷可靠；死亡数字和“十日”叙事须标为后出见证并与地方材料互证

\paragraph{顺治二年四—五月（1645年）：扬州攻陷后发生大规模杀戮与掳掠，规模和持续时间存在重大争议}
\eventanchor{EQ-1645-YANGZHOU-DECISION}{扬州攻陷后发生大规模杀戮与掳掠，规模和持续时间存在重大争议} 标记本节点定位围绕「扬州攻陷后发生大规模杀戮与掳掠，规模和持续时间存在重大争议」形成的处置与调度记录，涉及清军与扬州。京师与多地军政线路相接的尺度不是抽象背景：扬州守军扼守清军南下通道，攻城后的军纪、报复和掠夺处于战时权力失控与奖惩机制中。《清世祖实录》顺治二年四月扬州军报；《清史稿·史可法传》；王秀楚《扬州十日记》与《扬州府志》所见的顺序随后指向扬州经验在江南传播，既促成恐惧性降附，也强化若干城镇抵抗及后世纪念政治，但原有置信注记：扬州失陷可靠；死亡数字和“十日”叙事须标为后出见证并与地方材料互证

\paragraph{顺治二年四—五月（1645年）：扬州攻陷后发生大规模杀戮与掳掠，规模和持续时间存在重大争议}
在京师与多地军政线路相接的尺度，\eventanchor{EQ-1645-YANGZHOU-PRELUDE}{扬州攻陷后发生大规模杀戮与掳掠，规模和持续时间存在重大争议} 让清军与扬州的一个选择或压力有了可查的坐标。本节点回看「扬州攻陷后发生大规模杀戮与掳掠，规模和持续时间存在重大争议」之前已可见的条件。前因可写为扬州守军扼守清军南下通道，攻城后的军纪、报复和掠夺处于战时权力失控与奖惩机制中；《清世祖实录》顺治二年四月扬州军报；《清史稿·史可法传》；王秀楚《扬州十日记》与《扬州府志》限定了记录，后续变化是扬州经验在江南传播，既促成恐惧性降附，也强化若干城镇抵抗及后世纪念政治。原有置信注记：扬州失陷可靠；死亡数字和“十日”叙事须标为后出见证并与地方材料互证

\paragraph{顺治二年六月至十月（1645年）：清军进取杭州、绍兴等地，南明鲁王政权转入浙东及海上活动}
\eventanchor{EQ-1645-ZHEJIANG-CONQUEST}{清军进取杭州、绍兴等地，南明鲁王政权转入浙东及海上活动} 所关联的是清与南明鲁王集团、浙江地方势力在港口、海峡与跨海航路相连的尺度中的行动。本节点叙述该事件本身。它从南京失守后浙东军民与明宗室试图保持独立指挥，清军沿水陆节点分割其资源展开，借《清世祖实录》顺治二年六月至十月；《清史稿·鲁王传》；《浙江通志》与鲁王政权文书辑录进入可检验的年序，并以清廷取得浙江大部城镇，鲁王集团与郑氏网络相连，海疆战争成为长期问题影响下一段安排。原有置信注记：城市失守与鲁王转移可靠；各府实际控制和海岛联系须逐地核查

\paragraph{顺治二年六月至十月（1645年）：清军进取杭州、绍兴等地，南明鲁王政权转入浙东及海上活动}
港口、海峡与跨海航路相连的尺度中的\eventanchor{EQ-1645-ZHEJIANG-CONQUEST-AFTERMATH}{清军进取杭州、绍兴等地，南明鲁王政权转入浙东及海上活动} 不能离开清与南明鲁王集团、浙江地方势力来解释。本节点追踪「清军进取杭州、绍兴等地，南明鲁王政权转入浙东及海上活动」之后可见的余波。南京失守后浙东军民与明宗室试图保持独立指挥，清军沿水陆节点分割其资源说明其形成的条件；《清世祖实录》顺治二年六月至十月；《清史稿·鲁王传》；《浙江通志》与鲁王政权文书辑录提供来源路径，清廷取得浙江大部城镇，鲁王集团与郑氏网络相连，海疆战争成为长期问题则是材料能够支持的近时结果。原有置信注记：城市失守与鲁王转移可靠；各府实际控制和海岛联系须逐地核查

\paragraph{顺治二年六月至十月（1645年）：清军进取杭州、绍兴等地，南明鲁王政权转入浙东及海上活动}
\eventanchor{EQ-1645-ZHEJIANG-CONQUEST-DECISION}{清军进取杭州、绍兴等地，南明鲁王政权转入浙东及海上活动} 把清与南明鲁王集团、浙江地方势力放回港口、海峡与跨海航路相连的尺度的道路、城镇、边界或制度网络。本节点定位围绕「清军进取杭州、绍兴等地，南明鲁王政权转入浙东及海上活动」形成的处置与调度记录。南京失守后浙东军民与明宗室试图保持独立指挥，清军沿水陆节点分割其资源。本段采用《清世祖实录》顺治二年六月至十月；《清史稿·鲁王传》；《浙江通志》与鲁王政权文书辑录核对其顺序，随后可追到清廷取得浙江大部城镇，鲁王集团与郑氏网络相连，海疆战争成为长期问题。原有置信注记：城市失守与鲁王转移可靠；各府实际控制和海岛联系须逐地核查

\paragraph{顺治二年六月至十月（1645年）：清军进取杭州、绍兴等地，南明鲁王政权转入浙东及海上活动}
本节点回看「清军进取杭州、绍兴等地，南明鲁王政权转入浙东及海上活动」之前已可见的条件，故\eventanchor{EQ-1645-ZHEJIANG-CONQUEST-PRELUDE}{清军进取杭州、绍兴等地，南明鲁王政权转入浙东及海上活动} 不只是一个年份标签。清与南明鲁王集团、浙江地方势力在港口、海峡与跨海航路相连的尺度承受的压力来自南京失守后浙东军民与明宗室试图保持独立指挥，清军沿水陆节点分割其资源。《清世祖实录》顺治二年六月至十月；《清史稿·鲁王传》；《浙江通志》与鲁王政权文书辑录可回查该节点；清廷取得浙江大部城镇，鲁王集团与郑氏网络相连，海疆战争成为长期问题是其进入后续年序的方式。原有置信注记：城市失守与鲁王转移可靠；各府实际控制和海岛联系须逐地核查

\paragraph{顺治三年八月（1646年）：清军在闽北俘杀隆武帝朱聿键，隆武朝终结}
\eventanchor{EQ-1646-LONGWU-CAPTURE}{清军在闽北俘杀隆武帝朱聿键，隆武朝终结} 所见的行动者为清与南明隆武政权，空间尺度为京师与多地军政线路相接的尺度。本节点叙述该事件本身。清军由浙赣入闽，隆武政权军饷和将领协调不足且难以掌握山海通道使它成为具体事件链的一环；《清世祖实录》顺治三年八月；《清史稿·隆武帝本纪》；《福建通志》与黄道周文集相关记载与闽中明廷名义被打断，郑氏网络和地方武装成为东南抗清的重要行动者。原有置信注记：被俘死亡主干可靠；行止路线与地方助力记载不一

\paragraph{顺治三年八月（1646年）：清军在闽北俘杀隆武帝朱聿键，隆武朝终结}
从京师与多地军政线路相接的尺度看，\eventanchor{EQ-1646-LONGWU-CAPTURE-AFTERMATH}{清军在闽北俘杀隆武帝朱聿键，隆武朝终结} 留下本节点追踪「清军在闽北俘杀隆武帝朱聿键，隆武朝终结」之后可见的余波的材料坐标。清与南明隆武政权所面对的条件是清军由浙赣入闽，隆武政权军饷和将领协调不足且难以掌握山海通道。《清世祖实录》顺治三年八月；《清史稿·隆武帝本纪》；《福建通志》与黄道周文集相关记载支持的结果为闽中明廷名义被打断，郑氏网络和地方武装成为东南抗清的重要行动者。原有置信注记：被俘死亡主干可靠；行止路线与地方助力记载不一

\paragraph{顺治三年八月（1646年）：清军在闽北俘杀隆武帝朱聿键，隆武朝终结}
\eventanchor{EQ-1646-LONGWU-CAPTURE-DECISION}{清军在闽北俘杀隆武帝朱聿键，隆武朝终结} 记录清与南明隆武政权在京师与多地军政线路相接的尺度的一次变化。本节点定位围绕「清军在闽北俘杀隆武帝朱聿键，隆武朝终结」形成的处置与调度记录。其发生与清军由浙赣入闽，隆武政权军饷和将领协调不足且难以掌握山海通道相连；《清世祖实录》顺治三年八月；《清史稿·隆武帝本纪》；《福建通志》与黄道周文集相关记载把事件系回来源，闽中明廷名义被打断，郑氏网络和地方武装成为东南抗清的重要行动者显示压力如何向后转移。原有置信注记：被俘死亡主干可靠；行止路线与地方助力记载不一

\paragraph{顺治三年八月（1646年）：清军在闽北俘杀隆武帝朱聿键，隆武朝终结}
\eventanchor{EQ-1646-LONGWU-CAPTURE-PRELUDE}{清军在闽北俘杀隆武帝朱聿键，隆武朝终结} 的地点与行动范围属于京师与多地军政线路相接的尺度，行动者是清与南明隆武政权。本节点回看「清军在闽北俘杀隆武帝朱聿键，隆武朝终结」之前已可见的条件。清军由浙赣入闽，隆武政权军饷和将领协调不足且难以掌握山海通道。读《清世祖实录》顺治三年八月；《清史稿·隆武帝本纪》；《福建通志》与黄道周文集相关记载时可见其记录边界，最小后果只能写作闽中明廷名义被打断，郑氏网络和地方武装成为东南抗清的重要行动者。原有置信注记：被俘死亡主干可靠；行止路线与地方助力记载不一

\paragraph{顺治三年十一月（1646年）：清军攻入广州，绍武帝朱聿𨮁被杀，广州短暂的绍武朝结束}
\eventanchor{EQ-1646-SHAOWU-GUANGZHOU}{清军攻入广州，绍武帝朱聿𨮁被杀，广州短暂的绍武朝结束} 是清与南明绍武政权在京师与多地军政线路相接的尺度留下的编年标记。本节点叙述该事件本身。它从隆武朝覆灭后明宗室在广东争夺名义中心，内部冲突削弱对清军进攻的协调走来；《清世祖实录》顺治三年十一月；《清史稿·绍武帝本纪》；《广东通志》及南明文集固定可证部分，永历政权成为西南主要明廷名义，广东沿海郑氏与地方势力继续影响战局构成下一步的可见结果。原有置信注记：广州失守可靠；绍武与永历两朝冲突细节主要来自后出南明叙事

\paragraph{顺治三年十一月（1646年）：清军攻入广州，绍武帝朱聿𨮁被杀，广州短暂的绍武朝结束}
本节点追踪「清军攻入广州，绍武帝朱聿𨮁被杀，广州短暂的绍武朝结束」之后可见的余波：\eventanchor{EQ-1646-SHAOWU-GUANGZHOU-AFTERMATH}{清军攻入广州，绍武帝朱聿𨮁被杀，广州短暂的绍武朝结束}。清与南明绍武政权在京师与多地军政线路相接的尺度的处境可由隆武朝覆灭后明宗室在广东争夺名义中心，内部冲突削弱对清军进攻的协调说明。《清世祖实录》顺治三年十一月；《清史稿·绍武帝本纪》；《广东通志》及南明文集是本段材料入口，因而只能据此追踪到永历政权成为西南主要明廷名义，广东沿海郑氏与地方势力继续影响战局。原有置信注记：广州失守可靠；绍武与永历两朝冲突细节主要来自后出南明叙事

\paragraph{顺治三年十一月（1646年）：清军攻入广州，绍武帝朱聿𨮁被杀，广州短暂的绍武朝结束}
\eventanchor{EQ-1646-SHAOWU-GUANGZHOU-DECISION}{清军攻入广州，绍武帝朱聿𨮁被杀，广州短暂的绍武朝结束} 发生在京师与多地军政线路相接的尺度，牵动清与南明绍武政权。本节点定位围绕「清军攻入广州，绍武帝朱聿𨮁被杀，广州短暂的绍武朝结束」形成的处置与调度记录。隆武朝覆灭后明宗室在广东争夺名义中心，内部冲突削弱对清军进攻的协调是其结构性条件；《清世祖实录》顺治三年十一月；《清史稿·绍武帝本纪》；《广东通志》及南明文集留下的证据把读者带到永历政权成为西南主要明廷名义，广东沿海郑氏与地方势力继续影响战局。原有置信注记：广州失守可靠；绍武与永历两朝冲突细节主要来自后出南明叙事

\paragraph{顺治三年十一月（1646年）：清军攻入广州，绍武帝朱聿𨮁被杀，广州短暂的绍武朝结束}
\eventanchor{EQ-1646-SHAOWU-GUANGZHOU-PRELUDE}{清军攻入广州，绍武帝朱聿𨮁被杀，广州短暂的绍武朝结束} 对应本节点回看「清军攻入广州，绍武帝朱聿𨮁被杀，广州短暂的绍武朝结束」之前已可见的条件。在京师与多地军政线路相接的尺度，清与南明绍武政权无法绕开隆武朝覆灭后明宗室在广东争夺名义中心，内部冲突削弱对清军进攻的协调。依据《清世祖实录》顺治三年十一月；《清史稿·绍武帝本纪》；《广东通志》及南明文集，本段把后续影响限定为永历政权成为西南主要明廷名义，广东沿海郑氏与地方势力继续影响战局。原有置信注记：广州失守可靠；绍武与永历两朝冲突细节主要来自后出南明叙事

\paragraph{永历元年／顺治三年（1646年冬）：朱由榔在肇庆即位为永历帝，建立延续至西南的南明朝廷}
\eventanchor{EQ-1646-YONGLI-ZHAOQING}{朱由榔在肇庆即位为永历帝，建立延续至西南的南明朝廷}。本节点叙述该事件本身。南明永历政权与清在京师与多地军政线路相接的尺度的处境首先受绍武朝覆灭及明宗室分散，迫使支持者在两广重建朝廷；西南军政资源提供退守条件制约。《清世祖实录》顺治三年两广军报；《清史稿·永历帝本纪》；《永历实录》与《肇庆府志》留下可核的线索；可辨的后果是永历年号成为西南抗清行动的共同名义，清军须把两广、贵州、云南与缅甸路线连为战场。原有置信注记：即位时间大体可靠；改元与礼仪细节应按永历文书和清廷军报分别处理

\paragraph{永历元年／顺治三年（1646年冬）：朱由榔在肇庆即位为永历帝，建立延续至西南的南明朝廷}
\eventanchor{EQ-1646-YONGLI-ZHAOQING-AFTERMATH}{朱由榔在肇庆即位为永历帝，建立延续至西南的南明朝廷} 置于京师与多地军政线路相接的尺度中阅读，本节点追踪「朱由榔在肇庆即位为永历帝，建立延续至西南的南明朝廷」之后可见的余波。绍武朝覆灭及明宗室分散，迫使支持者在两广重建朝廷；西南军政资源提供退守条件构成行动者南明永历政权与清必须面对的条件。取证从《清世祖实录》顺治三年两广军报；《清史稿·永历帝本纪》；《永历实录》与《肇庆府志》出发，因而只能把变化谨慎地连到永历年号成为西南抗清行动的共同名义，清军须把两广、贵州、云南与缅甸路线连为战场。原有置信注记：即位时间大体可靠；改元与礼仪细节应按永历文书和清廷军报分别处理

\paragraph{永历元年／顺治三年（1646年冬）：朱由榔在肇庆即位为永历帝，建立延续至西南的南明朝廷}
\eventanchor{EQ-1646-YONGLI-ZHAOQING-DECISION}{朱由榔在肇庆即位为永历帝，建立延续至西南的南明朝廷} 标记本节点定位围绕「朱由榔在肇庆即位为永历帝，建立延续至西南的南明朝廷」形成的处置与调度记录，涉及南明永历政权与清。京师与多地军政线路相接的尺度不是抽象背景：绍武朝覆灭及明宗室分散，迫使支持者在两广重建朝廷；西南军政资源提供退守条件。《清世祖实录》顺治三年两广军报；《清史稿·永历帝本纪》；《永历实录》与《肇庆府志》所见的顺序随后指向永历年号成为西南抗清行动的共同名义，清军须把两广、贵州、云南与缅甸路线连为战场，但原有置信注记：即位时间大体可靠；改元与礼仪细节应按永历文书和清廷军报分别处理

\paragraph{永历元年／顺治三年（1646年冬）：朱由榔在肇庆即位为永历帝，建立延续至西南的南明朝廷}
在京师与多地军政线路相接的尺度，\eventanchor{EQ-1646-YONGLI-ZHAOQING-PRELUDE}{朱由榔在肇庆即位为永历帝，建立延续至西南的南明朝廷} 让南明永历政权与清的一个选择或压力有了可查的坐标。本节点回看「朱由榔在肇庆即位为永历帝，建立延续至西南的南明朝廷」之前已可见的条件。前因可写为绍武朝覆灭及明宗室分散，迫使支持者在两广重建朝廷；西南军政资源提供退守条件；《清世祖实录》顺治三年两广军报；《清史稿·永历帝本纪》；《永历实录》与《肇庆府志》限定了记录，后续变化是永历年号成为西南抗清行动的共同名义，清军须把两广、贵州、云南与缅甸路线连为战场。原有置信注记：即位时间大体可靠；改元与礼仪细节应按永历文书和清廷军报分别处理

\paragraph{顺治五年（1648年）：南昌守将金声桓反清并改奉南明，江西战局重新打开}
\eventanchor{EQ-1648-JIN-SHENGHUAN}{南昌守将金声桓反清并改奉南明，江西战局重新打开} 所关联的是清、江西与金声桓集团在京师与多地军政线路相接的尺度中的行动。本节点叙述该事件本身。它从清初任用降将以快速覆盖江南，但军饷、地位和地方网络未稳定；南明仍可提供合法性资源展开，借《清世祖实录》顺治五年江西军报；《清史稿·金声桓传》；《南昌府志》与南明永历朝文书进入可检验的年序，并以清廷被迫向江西投入围城和调兵资源，显示占领区行政接收不等于长期服从影响下一段安排。原有置信注记：反叛及南昌控制可靠；部众规模和改奉动机多为交战各方归责语言

\paragraph{顺治五年（1648年）：南昌守将金声桓反清并改奉南明，江西战局重新打开}
京师与多地军政线路相接的尺度中的\eventanchor{EQ-1648-JIN-SHENGHUAN-AFTERMATH}{南昌守将金声桓反清并改奉南明，江西战局重新打开} 不能离开清、江西与金声桓集团来解释。本节点追踪「南昌守将金声桓反清并改奉南明，江西战局重新打开」之后可见的余波。清初任用降将以快速覆盖江南，但军饷、地位和地方网络未稳定；南明仍可提供合法性资源说明其形成的条件；《清世祖实录》顺治五年江西军报；《清史稿·金声桓传》；《南昌府志》与南明永历朝文书提供来源路径，清廷被迫向江西投入围城和调兵资源，显示占领区行政接收不等于长期服从则是材料能够支持的近时结果。原有置信注记：反叛及南昌控制可靠；部众规模和改奉动机多为交战各方归责语言

\paragraph{顺治五年（1648年）：南昌守将金声桓反清并改奉南明，江西战局重新打开}
\eventanchor{EQ-1648-JIN-SHENGHUAN-DECISION}{南昌守将金声桓反清并改奉南明，江西战局重新打开} 把清、江西与金声桓集团放回京师与多地军政线路相接的尺度的道路、城镇、边界或制度网络。本节点定位围绕「南昌守将金声桓反清并改奉南明，江西战局重新打开」形成的处置与调度记录。清初任用降将以快速覆盖江南，但军饷、地位和地方网络未稳定；南明仍可提供合法性资源。本段采用《清世祖实录》顺治五年江西军报；《清史稿·金声桓传》；《南昌府志》与南明永历朝文书核对其顺序，随后可追到清廷被迫向江西投入围城和调兵资源，显示占领区行政接收不等于长期服从。原有置信注记：反叛及南昌控制可靠；部众规模和改奉动机多为交战各方归责语言

\paragraph{顺治五年（1648年）：南昌守将金声桓反清并改奉南明，江西战局重新打开}
本节点回看「南昌守将金声桓反清并改奉南明，江西战局重新打开」之前已可见的条件，故\eventanchor{EQ-1648-JIN-SHENGHUAN-PRELUDE}{南昌守将金声桓反清并改奉南明，江西战局重新打开} 不只是一个年份标签。清、江西与金声桓集团在京师与多地军政线路相接的尺度承受的压力来自清初任用降将以快速覆盖江南，但军饷、地位和地方网络未稳定；南明仍可提供合法性资源。《清世祖实录》顺治五年江西军报；《清史稿·金声桓传》；《南昌府志》与南明永历朝文书可回查该节点；清廷被迫向江西投入围城和调兵资源，显示占领区行政接收不等于长期服从是其进入后续年序的方式。原有置信注记：反叛及南昌控制可靠；部众规模和改奉动机多为交战各方归责语言

\paragraph{顺治五年（1648年）：广东清将李成栋反清，广州一度转入永历朝控制}
\eventanchor{EQ-1648-LI-CHENGDONG}{广东清将李成栋反清，广州一度转入永历朝控制} 所见的行动者为清、广东与李成栋集团，空间尺度为省、府县与地方社会相互牵动的尺度。本节点叙述该事件本身。清军依靠原明将领维系两广，军功分配与中央控制的矛盾给永历朝重新联络的空间使它成为具体事件链的一环；《清世祖实录》顺治五年广东军报；《清史稿·李成栋传》；《广东通志》及永历朝诏敕辑录与两广战场扩大并与江西、闽海相连，清廷须重新部署南方军饷和将领体系。原有置信注记：反叛与广州易手可靠；个人动机和地方支持范围有强烈阵营叙事差异

\paragraph{顺治五年（1648年）：广东清将李成栋反清，广州一度转入永历朝控制}
从省、府县与地方社会相互牵动的尺度看，\eventanchor{EQ-1648-LI-CHENGDONG-AFTERMATH}{广东清将李成栋反清，广州一度转入永历朝控制} 留下本节点追踪「广东清将李成栋反清，广州一度转入永历朝控制」之后可见的余波的材料坐标。清、广东与李成栋集团所面对的条件是清军依靠原明将领维系两广，军功分配与中央控制的矛盾给永历朝重新联络的空间。《清世祖实录》顺治五年广东军报；《清史稿·李成栋传》；《广东通志》及永历朝诏敕辑录支持的结果为两广战场扩大并与江西、闽海相连，清廷须重新部署南方军饷和将领体系。原有置信注记：反叛与广州易手可靠；个人动机和地方支持范围有强烈阵营叙事差异

\paragraph{顺治五年（1648年）：广东清将李成栋反清，广州一度转入永历朝控制}
\eventanchor{EQ-1648-LI-CHENGDONG-DECISION}{广东清将李成栋反清，广州一度转入永历朝控制} 记录清、广东与李成栋集团在省、府县与地方社会相互牵动的尺度的一次变化。本节点定位围绕「广东清将李成栋反清，广州一度转入永历朝控制」形成的处置与调度记录。其发生与清军依靠原明将领维系两广，军功分配与中央控制的矛盾给永历朝重新联络的空间相连；《清世祖实录》顺治五年广东军报；《清史稿·李成栋传》；《广东通志》及永历朝诏敕辑录把事件系回来源，两广战场扩大并与江西、闽海相连，清廷须重新部署南方军饷和将领体系显示压力如何向后转移。原有置信注记：反叛与广州易手可靠；个人动机和地方支持范围有强烈阵营叙事差异

\paragraph{顺治五年（1648年）：广东清将李成栋反清，广州一度转入永历朝控制}
\eventanchor{EQ-1648-LI-CHENGDONG-PRELUDE}{广东清将李成栋反清，广州一度转入永历朝控制} 的地点与行动范围属于省、府县与地方社会相互牵动的尺度，行动者是清、广东与李成栋集团。本节点回看「广东清将李成栋反清，广州一度转入永历朝控制」之前已可见的条件。清军依靠原明将领维系两广，军功分配与中央控制的矛盾给永历朝重新联络的空间。读《清世祖实录》顺治五年广东军报；《清史稿·李成栋传》；《广东通志》及永历朝诏敕辑录时可见其记录边界，最小后果只能写作两广战场扩大并与江西、闽海相连，清廷须重新部署南方军饷和将领体系。原有置信注记：反叛与广州易手可靠；个人动机和地方支持范围有强烈阵营叙事差异

\paragraph{顺治六年正月至三月（1649年）：清军围攻并攻取南昌，金声桓反叛失败}
\eventanchor{EQ-1649-NANCHANG}{清军围攻并攻取南昌，金声桓反叛失败} 是清与金声桓集团在京师与多地军政线路相接的尺度留下的编年标记。本节点叙述该事件本身。它从金声桓据南昌切断清廷赣江通道，清军需以持续军需和外围据点完成围攻走来；《清世祖实录》顺治六年正月至三月；《清史稿·金声桓传》；内阁大库顺治六年江西军需题本与《南昌府志》固定可证部分，江西反叛中心被击破，清廷恢复通道但须面对战后税源、流民和秩序重建构成下一步的可见结果。原有置信注记：围城结局可靠；攻城损失、城内处置和战后人口变化须以地方资料复核

\paragraph{顺治六年正月至三月（1649年）：清军围攻并攻取南昌，金声桓反叛失败}
本节点追踪「清军围攻并攻取南昌，金声桓反叛失败」之后可见的余波：\eventanchor{EQ-1649-NANCHANG-AFTERMATH}{清军围攻并攻取南昌，金声桓反叛失败}。清与金声桓集团在京师与多地军政线路相接的尺度的处境可由金声桓据南昌切断清廷赣江通道，清军需以持续军需和外围据点完成围攻说明。《清世祖实录》顺治六年正月至三月；《清史稿·金声桓传》；内阁大库顺治六年江西军需题本与《南昌府志》是本段材料入口，因而只能据此追踪到江西反叛中心被击破，清廷恢复通道但须面对战后税源、流民和秩序重建。原有置信注记：围城结局可靠；攻城损失、城内处置和战后人口变化须以地方资料复核

\paragraph{顺治六年正月至三月（1649年）：清军围攻并攻取南昌，金声桓反叛失败}
\eventanchor{EQ-1649-NANCHANG-DECISION}{清军围攻并攻取南昌，金声桓反叛失败} 发生在京师与多地军政线路相接的尺度，牵动清与金声桓集团。本节点定位围绕「清军围攻并攻取南昌，金声桓反叛失败」形成的处置与调度记录。金声桓据南昌切断清廷赣江通道，清军需以持续军需和外围据点完成围攻是其结构性条件；《清世祖实录》顺治六年正月至三月；《清史稿·金声桓传》；内阁大库顺治六年江西军需题本与《南昌府志》留下的证据把读者带到江西反叛中心被击破，清廷恢复通道但须面对战后税源、流民和秩序重建。原有置信注记：围城结局可靠；攻城损失、城内处置和战后人口变化须以地方资料复核

\paragraph{顺治六年正月至三月（1649年）：清军围攻并攻取南昌，金声桓反叛失败}
\eventanchor{EQ-1649-NANCHANG-PRELUDE}{清军围攻并攻取南昌，金声桓反叛失败} 对应本节点回看「清军围攻并攻取南昌，金声桓反叛失败」之前已可见的条件。在京师与多地军政线路相接的尺度，清与金声桓集团无法绕开金声桓据南昌切断清廷赣江通道，清军需以持续军需和外围据点完成围攻。依据《清世祖实录》顺治六年正月至三月；《清史稿·金声桓传》；内阁大库顺治六年江西军需题本与《南昌府志》，本段把后续影响限定为江西反叛中心被击破，清廷恢复通道但须面对战后税源、流民和秩序重建。原有置信注记：围城结局可靠；攻城损失、城内处置和战后人口变化须以地方资料复核

\paragraph{顺治七年十一月（1650年）：清军重新攻取广州，战后杀戮与损毁规模成为广东地方记忆中的重大争议}
\eventanchor{EQ-1650-GUANGZHOU}{清军重新攻取广州，战后杀戮与损毁规模成为广东地方记忆中的重大争议}。本节点叙述该事件本身。清与永历朝广东守军在省、府县与地方社会相互牵动的尺度的处境首先受李成栋反叛后广东成为永历朝资源节点，清军重夺广州兼具军事和海防目的制约。《清世祖实录》顺治七年十一月；《清史稿·世祖本纪》；《广东通志》、广州地方笔记与内阁大库广东军报留下可核的线索；可辨的后果是永历朝在两广基础削弱，广州损失与人口流动影响地方财政、宗族和港口恢复。原有置信注记：广州再陷可靠；死亡数字和“庚寅之劫”持续时间不能从清军战报单定

\paragraph{顺治七年十一月（1650年）：清军重新攻取广州，战后杀戮与损毁规模成为广东地方记忆中的重大争议}
\eventanchor{EQ-1650-GUANGZHOU-AFTERMATH}{清军重新攻取广州，战后杀戮与损毁规模成为广东地方记忆中的重大争议} 置于省、府县与地方社会相互牵动的尺度中阅读，本节点追踪「清军重新攻取广州，战后杀戮与损毁规模成为广东地方记忆中的重大争议」之后可见的余波。李成栋反叛后广东成为永历朝资源节点，清军重夺广州兼具军事和海防目的构成行动者清与永历朝广东守军必须面对的条件。取证从《清世祖实录》顺治七年十一月；《清史稿·世祖本纪》；《广东通志》、广州地方笔记与内阁大库广东军报出发，因而只能把变化谨慎地连到永历朝在两广基础削弱，广州损失与人口流动影响地方财政、宗族和港口恢复。原有置信注记：广州再陷可靠；死亡数字和“庚寅之劫”持续时间不能从清军战报单定

\paragraph{顺治七年十一月（1650年）：清军重新攻取广州，战后杀戮与损毁规模成为广东地方记忆中的重大争议}
\eventanchor{EQ-1650-GUANGZHOU-DECISION}{清军重新攻取广州，战后杀戮与损毁规模成为广东地方记忆中的重大争议} 标记本节点定位围绕「清军重新攻取广州，战后杀戮与损毁规模成为广东地方记忆中的重大争议」形成的处置与调度记录，涉及清与永历朝广东守军。省、府县与地方社会相互牵动的尺度不是抽象背景：李成栋反叛后广东成为永历朝资源节点，清军重夺广州兼具军事和海防目的。《清世祖实录》顺治七年十一月；《清史稿·世祖本纪》；《广东通志》、广州地方笔记与内阁大库广东军报所见的顺序随后指向永历朝在两广基础削弱，广州损失与人口流动影响地方财政、宗族和港口恢复，但原有置信注记：广州再陷可靠；死亡数字和“庚寅之劫”持续时间不能从清军战报单定

\paragraph{顺治七年十一月（1650年）：清军重新攻取广州，战后杀戮与损毁规模成为广东地方记忆中的重大争议}
在省、府县与地方社会相互牵动的尺度，\eventanchor{EQ-1650-GUANGZHOU-PRELUDE}{清军重新攻取广州，战后杀戮与损毁规模成为广东地方记忆中的重大争议} 让清与永历朝广东守军的一个选择或压力有了可查的坐标。本节点回看「清军重新攻取广州，战后杀戮与损毁规模成为广东地方记忆中的重大争议」之前已可见的条件。前因可写为李成栋反叛后广东成为永历朝资源节点，清军重夺广州兼具军事和海防目的；《清世祖实录》顺治七年十一月；《清史稿·世祖本纪》；《广东通志》、广州地方笔记与内阁大库广东军报限定了记录，后续变化是永历朝在两广基础削弱，广州损失与人口流动影响地方财政、宗族和港口恢复。原有置信注记：广州再陷可靠；死亡数字和“庚寅之劫”持续时间不能从清军战报单定

\paragraph{永历六年／顺治九年（1652年）：李定国攻克桂林，孔有德自尽，清在广西的军事据点受重创}
\eventanchor{EQ-1652-DINGGUO-GUILIN}{李定国攻克桂林，孔有德自尽，清在广西的军事据点受重创} 所关联的是南明永历朝、清与李定国军在京师与多地军政线路相接的尺度中的行动。本节点叙述该事件本身。它从两广清军战线长、补给困难且将领协调不足，永历军利用机动兵力突击桂林展开，借《清世祖实录》顺治九年广西军报；《清史稿·孔有德传》；《永历实录》与《广西通志》进入可检验的年序，并以永历朝获得声望与物资，但未能稳定保持两广；清廷随即加强西南调兵影响下一段安排。原有置信注记：桂林战果可靠；双方兵数和孔有德死因细节应保留不同文本说法

\paragraph{永历六年／顺治九年（1652年）：李定国攻克桂林，孔有德自尽，清在广西的军事据点受重创}
京师与多地军政线路相接的尺度中的\eventanchor{EQ-1652-DINGGUO-GUILIN-AFTERMATH}{李定国攻克桂林，孔有德自尽，清在广西的军事据点受重创} 不能离开南明永历朝、清与李定国军来解释。本节点追踪「李定国攻克桂林，孔有德自尽，清在广西的军事据点受重创」之后可见的余波。两广清军战线长、补给困难且将领协调不足，永历军利用机动兵力突击桂林说明其形成的条件；《清世祖实录》顺治九年广西军报；《清史稿·孔有德传》；《永历实录》与《广西通志》提供来源路径，永历朝获得声望与物资，但未能稳定保持两广；清廷随即加强西南调兵则是材料能够支持的近时结果。原有置信注记：桂林战果可靠；双方兵数和孔有德死因细节应保留不同文本说法

\paragraph{永历六年／顺治九年（1652年）：李定国攻克桂林，孔有德自尽，清在广西的军事据点受重创}
\eventanchor{EQ-1652-DINGGUO-GUILIN-DECISION}{李定国攻克桂林，孔有德自尽，清在广西的军事据点受重创} 把南明永历朝、清与李定国军放回京师与多地军政线路相接的尺度的道路、城镇、边界或制度网络。本节点定位围绕「李定国攻克桂林，孔有德自尽，清在广西的军事据点受重创」形成的处置与调度记录。两广清军战线长、补给困难且将领协调不足，永历军利用机动兵力突击桂林。本段采用《清世祖实录》顺治九年广西军报；《清史稿·孔有德传》；《永历实录》与《广西通志》核对其顺序，随后可追到永历朝获得声望与物资，但未能稳定保持两广；清廷随即加强西南调兵。原有置信注记：桂林战果可靠；双方兵数和孔有德死因细节应保留不同文本说法

\paragraph{永历六年／顺治九年（1652年）：李定国攻克桂林，孔有德自尽，清在广西的军事据点受重创}
本节点回看「李定国攻克桂林，孔有德自尽，清在广西的军事据点受重创」之前已可见的条件，故\eventanchor{EQ-1652-DINGGUO-GUILIN-PRELUDE}{李定国攻克桂林，孔有德自尽，清在广西的军事据点受重创} 不只是一个年份标签。南明永历朝、清与李定国军在京师与多地军政线路相接的尺度承受的压力来自两广清军战线长、补给困难且将领协调不足，永历军利用机动兵力突击桂林。《清世祖实录》顺治九年广西军报；《清史稿·孔有德传》；《永历实录》与《广西通志》可回查该节点；永历朝获得声望与物资，但未能稳定保持两广；清廷随即加强西南调兵是其进入后续年序的方式。原有置信注记：桂林战果可靠；双方兵数和孔有德死因细节应保留不同文本说法

\paragraph{永历六年／顺治九年夏（1652年）：李定国在衡州击败尼堪军，清军西南攻势暂受挫}
\eventanchor{EQ-1652-DINGGUO-HENGZHOU}{李定国在衡州击败尼堪军，清军西南攻势暂受挫} 所见的行动者为南明永历朝与清，空间尺度为京师与多地军政线路相接的尺度。本节点叙述该事件本身。清军试图由湖南压迫两广和贵州，李定国依托本地军队与道路选择反击使它成为具体事件链的一环；《清世祖实录》顺治九年湖南军报；《清史稿·尼堪传》；《永历实录》与《衡州府志》与清廷认识到西南征服不能只靠快速突进，转向更大规模调兵、饷运和分区推进。原有置信注记：战役结局大体可靠；地点、伤亡与尼堪死亡过程在中清叙事与南明记录间有异

\paragraph{永历六年／顺治九年夏（1652年）：李定国在衡州击败尼堪军，清军西南攻势暂受挫}
从京师与多地军政线路相接的尺度看，\eventanchor{EQ-1652-DINGGUO-HENGZHOU-AFTERMATH}{李定国在衡州击败尼堪军，清军西南攻势暂受挫} 留下本节点追踪「李定国在衡州击败尼堪军，清军西南攻势暂受挫」之后可见的余波的材料坐标。南明永历朝与清所面对的条件是清军试图由湖南压迫两广和贵州，李定国依托本地军队与道路选择反击。《清世祖实录》顺治九年湖南军报；《清史稿·尼堪传》；《永历实录》与《衡州府志》支持的结果为清廷认识到西南征服不能只靠快速突进，转向更大规模调兵、饷运和分区推进。原有置信注记：战役结局大体可靠；地点、伤亡与尼堪死亡过程在中清叙事与南明记录间有异

\paragraph{永历六年／顺治九年夏（1652年）：李定国在衡州击败尼堪军，清军西南攻势暂受挫}
\eventanchor{EQ-1652-DINGGUO-HENGZHOU-DECISION}{李定国在衡州击败尼堪军，清军西南攻势暂受挫} 记录南明永历朝与清在京师与多地军政线路相接的尺度的一次变化。本节点定位围绕「李定国在衡州击败尼堪军，清军西南攻势暂受挫」形成的处置与调度记录。其发生与清军试图由湖南压迫两广和贵州，李定国依托本地军队与道路选择反击相连；《清世祖实录》顺治九年湖南军报；《清史稿·尼堪传》；《永历实录》与《衡州府志》把事件系回来源，清廷认识到西南征服不能只靠快速突进，转向更大规模调兵、饷运和分区推进显示压力如何向后转移。原有置信注记：战役结局大体可靠；地点、伤亡与尼堪死亡过程在中清叙事与南明记录间有异

\paragraph{永历六年／顺治九年夏（1652年）：李定国在衡州击败尼堪军，清军西南攻势暂受挫}
\eventanchor{EQ-1652-DINGGUO-HENGZHOU-PRELUDE}{李定国在衡州击败尼堪军，清军西南攻势暂受挫} 的地点与行动范围属于京师与多地军政线路相接的尺度，行动者是南明永历朝与清。本节点回看「李定国在衡州击败尼堪军，清军西南攻势暂受挫」之前已可见的条件。清军试图由湖南压迫两广和贵州，李定国依托本地军队与道路选择反击。读《清世祖实录》顺治九年湖南军报；《清史稿·尼堪传》；《永历实录》与《衡州府志》时可见其记录边界，最小后果只能写作清廷认识到西南征服不能只靠快速突进，转向更大规模调兵、饷运和分区推进。原有置信注记：战役结局大体可靠；地点、伤亡与尼堪死亡过程在中清叙事与南明记录间有异

\paragraph{永历十一年／顺治十四年（1657年）：永历朝在贵州、云南的孙可望与李定国集团矛盾公开化，孙可望转降清廷}
\eventanchor{EQ-1657-YONGLI-GUIZHOU}{永历朝在贵州、云南的孙可望与李定国集团矛盾公开化，孙可望转降清廷} 是南明永历朝、清与孙可望集团在京师与多地军政线路相接的尺度留下的编年标记。本节点叙述该事件本身。它从长期战争中的饷源、官爵和指挥权竞争削弱永历朝合作，清廷提供孙可望转换阵营的条件走来；《清世祖实录》顺治十四年贵州军报；《清史稿·孙可望传》；《永历实录》与贵州地方志固定可证部分，清军获得西南军情与向导资源，永历朝失去大批可动员兵力，云南退路日益孤立构成下一步的可见结果。原有置信注记：孙可望降清可靠；冲突责任、军队规模与永历帝处境受各方政治文本影响

\paragraph{永历十一年／顺治十四年（1657年）：永历朝在贵州、云南的孙可望与李定国集团矛盾公开化，孙可望转降清廷}
本节点追踪「永历朝在贵州、云南的孙可望与李定国集团矛盾公开化，孙可望转降清廷」之后可见的余波：\eventanchor{EQ-1657-YONGLI-GUIZHOU-AFTERMATH}{永历朝在贵州、云南的孙可望与李定国集团矛盾公开化，孙可望转降清廷}。南明永历朝、清与孙可望集团在京师与多地军政线路相接的尺度的处境可由长期战争中的饷源、官爵和指挥权竞争削弱永历朝合作，清廷提供孙可望转换阵营的条件说明。《清世祖实录》顺治十四年贵州军报；《清史稿·孙可望传》；《永历实录》与贵州地方志是本段材料入口，因而只能据此追踪到清军获得西南军情与向导资源，永历朝失去大批可动员兵力，云南退路日益孤立。原有置信注记：孙可望降清可靠；冲突责任、军队规模与永历帝处境受各方政治文本影响

\paragraph{永历十一年／顺治十四年（1657年）：永历朝在贵州、云南的孙可望与李定国集团矛盾公开化，孙可望转降清廷}
\eventanchor{EQ-1657-YONGLI-GUIZHOU-DECISION}{永历朝在贵州、云南的孙可望与李定国集团矛盾公开化，孙可望转降清廷} 发生在京师与多地军政线路相接的尺度，牵动南明永历朝、清与孙可望集团。本节点定位围绕「永历朝在贵州、云南的孙可望与李定国集团矛盾公开化，孙可望转降清廷」形成的处置与调度记录。长期战争中的饷源、官爵和指挥权竞争削弱永历朝合作，清廷提供孙可望转换阵营的条件是其结构性条件；《清世祖实录》顺治十四年贵州军报；《清史稿·孙可望传》；《永历实录》与贵州地方志留下的证据把读者带到清军获得西南军情与向导资源，永历朝失去大批可动员兵力，云南退路日益孤立。原有置信注记：孙可望降清可靠；冲突责任、军队规模与永历帝处境受各方政治文本影响

\paragraph{永历十一年／顺治十四年（1657年）：永历朝在贵州、云南的孙可望与李定国集团矛盾公开化，孙可望转降清廷}
\eventanchor{EQ-1657-YONGLI-GUIZHOU-PRELUDE}{永历朝在贵州、云南的孙可望与李定国集团矛盾公开化，孙可望转降清廷} 对应本节点回看「永历朝在贵州、云南的孙可望与李定国集团矛盾公开化，孙可望转降清廷」之前已可见的条件。在京师与多地军政线路相接的尺度，南明永历朝、清与孙可望集团无法绕开长期战争中的饷源、官爵和指挥权竞争削弱永历朝合作，清廷提供孙可望转换阵营的条件。依据《清世祖实录》顺治十四年贵州军报；《清史稿·孙可望传》；《永历实录》与贵州地方志，本段把后续影响限定为清军获得西南军情与向导资源，永历朝失去大批可动员兵力，云南退路日益孤立。原有置信注记：孙可望降清可靠；冲突责任、军队规模与永历帝处境受各方政治文本影响

\paragraph{顺治十五年（1658年）：清军由贵州、四川等方向进逼云南，永历朝廷撤离昆明并向西南退却}
\eventanchor{EQ-1658-QING-YUNNAN}{清军由贵州、四川等方向进逼云南，永历朝廷撤离昆明并向西南退却}。本节点叙述该事件本身。清与南明永历朝在省、府县与地方社会相互牵动的尺度的处境首先受孙可望降清削弱南明防线，清军依托多路会攻和地方土司、降将信息推进制约。《清世祖实录》顺治十五年云南军报；《清史稿·洪承畴传》；《云南通志》与永历朝奏疏辑录留下可核的线索；可辨的后果是永历政权转为流动朝廷，清军虽占昆明等城，仍须在边地和交通线上持续作战。原有置信注记：清军入滇和永历撤离可靠；各路行军与地方归附应分区而非视作全省一次完成

\paragraph{顺治十五年（1658年）：清军由贵州、四川等方向进逼云南，永历朝廷撤离昆明并向西南退却}
\eventanchor{EQ-1658-QING-YUNNAN-AFTERMATH}{清军由贵州、四川等方向进逼云南，永历朝廷撤离昆明并向西南退却} 置于省、府县与地方社会相互牵动的尺度中阅读，本节点追踪「清军由贵州、四川等方向进逼云南，永历朝廷撤离昆明并向西南退却」之后可见的余波。孙可望降清削弱南明防线，清军依托多路会攻和地方土司、降将信息推进构成行动者清与南明永历朝必须面对的条件。取证从《清世祖实录》顺治十五年云南军报；《清史稿·洪承畴传》；《云南通志》与永历朝奏疏辑录出发，因而只能把变化谨慎地连到永历政权转为流动朝廷，清军虽占昆明等城，仍须在边地和交通线上持续作战。原有置信注记：清军入滇和永历撤离可靠；各路行军与地方归附应分区而非视作全省一次完成

\paragraph{顺治十五年（1658年）：清军由贵州、四川等方向进逼云南，永历朝廷撤离昆明并向西南退却}
\eventanchor{EQ-1658-QING-YUNNAN-DECISION}{清军由贵州、四川等方向进逼云南，永历朝廷撤离昆明并向西南退却} 标记本节点定位围绕「清军由贵州、四川等方向进逼云南，永历朝廷撤离昆明并向西南退却」形成的处置与调度记录，涉及清与南明永历朝。省、府县与地方社会相互牵动的尺度不是抽象背景：孙可望降清削弱南明防线，清军依托多路会攻和地方土司、降将信息推进。《清世祖实录》顺治十五年云南军报；《清史稿·洪承畴传》；《云南通志》与永历朝奏疏辑录所见的顺序随后指向永历政权转为流动朝廷，清军虽占昆明等城，仍须在边地和交通线上持续作战，但原有置信注记：清军入滇和永历撤离可靠；各路行军与地方归附应分区而非视作全省一次完成

\paragraph{顺治十五年（1658年）：清军由贵州、四川等方向进逼云南，永历朝廷撤离昆明并向西南退却}
在省、府县与地方社会相互牵动的尺度，\eventanchor{EQ-1658-QING-YUNNAN-PRELUDE}{清军由贵州、四川等方向进逼云南，永历朝廷撤离昆明并向西南退却} 让清与南明永历朝的一个选择或压力有了可查的坐标。本节点回看「清军由贵州、四川等方向进逼云南，永历朝廷撤离昆明并向西南退却」之前已可见的条件。前因可写为孙可望降清削弱南明防线，清军依托多路会攻和地方土司、降将信息推进；《清世祖实录》顺治十五年云南军报；《清史稿·洪承畴传》；《云南通志》与永历朝奏疏辑录限定了记录，后续变化是永历政权转为流动朝廷，清军虽占昆明等城，仍须在边地和交通线上持续作战。原有置信注记：清军入滇和永历撤离可靠；各路行军与地方归附应分区而非视作全省一次完成

\paragraph{永历十三年／顺治十六年末（1659年）：永历帝一行进入缅甸，西南战争扩展为跨境避难和外交问题}
\eventanchor{EQ-1659-YONGLI-MYANMAR}{永历帝一行进入缅甸，西南战争扩展为跨境避难和外交问题} 所关联的是南明永历朝、缅甸东吁王朝与清在边地道路、驻防与多方社群交错的尺度中的行动。本节点叙述该事件本身。它从清军推进使永历朝失去云南腹地，缅甸王朝面临接纳流亡朝廷与避免清军压力的两难展开，借《清世祖实录》顺治十六年云南军报；《清史稿·永历帝本纪》；《琉璃宫史》缅文传统及欧洲旅行记译注进入可检验的年序，并以永历政权补给与外交选择受缅甸宫廷制约，清廷将追捕与边境交涉结合影响下一段安排。原有置信注记：入缅事实可靠；路线、接待条件与缅甸朝廷决策须并读汉文和缅文/欧洲记载

\paragraph{永历十三年／顺治十六年末（1659年）：永历帝一行进入缅甸，西南战争扩展为跨境避难和外交问题}
边地道路、驻防与多方社群交错的尺度中的\eventanchor{EQ-1659-YONGLI-MYANMAR-AFTERMATH}{永历帝一行进入缅甸，西南战争扩展为跨境避难和外交问题} 不能离开南明永历朝、缅甸东吁王朝与清来解释。本节点追踪「永历帝一行进入缅甸，西南战争扩展为跨境避难和外交问题」之后可见的余波。清军推进使永历朝失去云南腹地，缅甸王朝面临接纳流亡朝廷与避免清军压力的两难说明其形成的条件；《清世祖实录》顺治十六年云南军报；《清史稿·永历帝本纪》；《琉璃宫史》缅文传统及欧洲旅行记译注提供来源路径，永历政权补给与外交选择受缅甸宫廷制约，清廷将追捕与边境交涉结合则是材料能够支持的近时结果。原有置信注记：入缅事实可靠；路线、接待条件与缅甸朝廷决策须并读汉文和缅文/欧洲记载

\paragraph{永历十三年／顺治十六年末（1659年）：永历帝一行进入缅甸，西南战争扩展为跨境避难和外交问题}
\eventanchor{EQ-1659-YONGLI-MYANMAR-DECISION}{永历帝一行进入缅甸，西南战争扩展为跨境避难和外交问题} 把南明永历朝、缅甸东吁王朝与清放回边地道路、驻防与多方社群交错的尺度的道路、城镇、边界或制度网络。本节点定位围绕「永历帝一行进入缅甸，西南战争扩展为跨境避难和外交问题」形成的处置与调度记录。清军推进使永历朝失去云南腹地，缅甸王朝面临接纳流亡朝廷与避免清军压力的两难。本段采用《清世祖实录》顺治十六年云南军报；《清史稿·永历帝本纪》；《琉璃宫史》缅文传统及欧洲旅行记译注核对其顺序，随后可追到永历政权补给与外交选择受缅甸宫廷制约，清廷将追捕与边境交涉结合。原有置信注记：入缅事实可靠；路线、接待条件与缅甸朝廷决策须并读汉文和缅文/欧洲记载

\paragraph{永历十三年／顺治十六年末（1659年）：永历帝一行进入缅甸，西南战争扩展为跨境避难和外交问题}
本节点回看「永历帝一行进入缅甸，西南战争扩展为跨境避难和外交问题」之前已可见的条件，故\eventanchor{EQ-1659-YONGLI-MYANMAR-PRELUDE}{永历帝一行进入缅甸，西南战争扩展为跨境避难和外交问题} 不只是一个年份标签。南明永历朝、缅甸东吁王朝与清在边地道路、驻防与多方社群交错的尺度承受的压力来自清军推进使永历朝失去云南腹地，缅甸王朝面临接纳流亡朝廷与避免清军压力的两难。《清世祖实录》顺治十六年云南军报；《清史稿·永历帝本纪》；《琉璃宫史》缅文传统及欧洲旅行记译注可回查该节点；永历政权补给与外交选择受缅甸宫廷制约，清廷将追捕与边境交涉结合是其进入后续年序的方式。原有置信注记：入缅事实可靠；路线、接待条件与缅甸朝廷决策须并读汉文和缅文/欧洲记载

\paragraph{顺治十七年（1660年）：清军继续经营云南城镇与驿路，永历朝残部向中缅边地退缩}
\eventanchor{EQ-1660-QING-YUNNAN-CONSOLIDATION}{清军继续经营云南城镇与驿路，永历朝残部向中缅边地退缩} 所见的行动者为清与南明永历朝残部，空间尺度为京师与多地军政线路相接的尺度。本节点叙述该事件本身。占领昆明后清军仍须供应远征军、维持道路并争取地方中介，永历残部利用边地机动使它成为具体事件链的一环；《清世祖实录》顺治十七年云南军报；《清史稿·吴三桂传》；内阁大库顺治十七年云南粮饷题本与《云南通志》与云南成为吴三桂等将领驻军和财政调度基地，也为后来的藩镇结构埋下条件。原有置信注记：清军控制主要城镇大体可靠；土司、矿区和山地实际控制不能以府城易手推定

\paragraph{顺治十七年（1660年）：清军继续经营云南城镇与驿路，永历朝残部向中缅边地退缩}
从京师与多地军政线路相接的尺度看，\eventanchor{EQ-1660-QING-YUNNAN-CONSOLIDATION-AFTERMATH}{清军继续经营云南城镇与驿路，永历朝残部向中缅边地退缩} 留下本节点追踪「清军继续经营云南城镇与驿路，永历朝残部向中缅边地退缩」之后可见的余波的材料坐标。清与南明永历朝残部所面对的条件是占领昆明后清军仍须供应远征军、维持道路并争取地方中介，永历残部利用边地机动。《清世祖实录》顺治十七年云南军报；《清史稿·吴三桂传》；内阁大库顺治十七年云南粮饷题本与《云南通志》支持的结果为云南成为吴三桂等将领驻军和财政调度基地，也为后来的藩镇结构埋下条件。原有置信注记：清军控制主要城镇大体可靠；土司、矿区和山地实际控制不能以府城易手推定

\paragraph{顺治十七年（1660年）：清军继续经营云南城镇与驿路，永历朝残部向中缅边地退缩}
\eventanchor{EQ-1660-QING-YUNNAN-CONSOLIDATION-DECISION}{清军继续经营云南城镇与驿路，永历朝残部向中缅边地退缩} 记录清与南明永历朝残部在京师与多地军政线路相接的尺度的一次变化。本节点定位围绕「清军继续经营云南城镇与驿路，永历朝残部向中缅边地退缩」形成的处置与调度记录。其发生与占领昆明后清军仍须供应远征军、维持道路并争取地方中介，永历残部利用边地机动相连；《清世祖实录》顺治十七年云南军报；《清史稿·吴三桂传》；内阁大库顺治十七年云南粮饷题本与《云南通志》把事件系回来源，云南成为吴三桂等将领驻军和财政调度基地，也为后来的藩镇结构埋下条件显示压力如何向后转移。原有置信注记：清军控制主要城镇大体可靠；土司、矿区和山地实际控制不能以府城易手推定

\paragraph{顺治十七年（1660年）：清军继续经营云南城镇与驿路，永历朝残部向中缅边地退缩}
\eventanchor{EQ-1660-QING-YUNNAN-CONSOLIDATION-PRELUDE}{清军继续经营云南城镇与驿路，永历朝残部向中缅边地退缩} 的地点与行动范围属于京师与多地军政线路相接的尺度，行动者是清与南明永历朝残部。本节点回看「清军继续经营云南城镇与驿路，永历朝残部向中缅边地退缩」之前已可见的条件。占领昆明后清军仍须供应远征军、维持道路并争取地方中介，永历残部利用边地机动。读《清世祖实录》顺治十七年云南军报；《清史稿·吴三桂传》；内阁大库顺治十七年云南粮饷题本与《云南通志》时可见其记录边界，最小后果只能写作云南成为吴三桂等将领驻军和财政调度基地，也为后来的藩镇结构埋下条件。原有置信注记：清军控制主要城镇大体可靠；土司、矿区和山地实际控制不能以府城易手推定

\paragraph{康熙元年四月（1662年）：吴三桂军将自缅甸押回的永历帝朱由榔处死，南明朝廷名义中心结束}
\eventanchor{EQ-1662-YONGLI-EXECUTION}{吴三桂军将自缅甸押回的永历帝朱由榔处死，南明朝廷名义中心结束} 是清、南明永历朝与缅甸在边地道路、驻防与多方社群交错的尺度留下的编年标记。本节点叙述该事件本身。它从永历朝长期依赖流亡和边地补给，缅甸宫廷在清军压力下改变安置策略，吴三桂完成追捕走来；《清圣祖实录》康熙元年四月；《清史稿·永历帝本纪》；《琉璃宫史》缅文传统与《云南通志》固定可证部分，南明皇帝名义不复存在，但西南残部、郑氏与地方记忆仍继续影响清初整合构成下一步的可见结果。原有置信注记：被押回和处死主干可靠；缅甸移交过程、处死地点与责任叙事须比对中缅材料

\paragraph{康熙元年四月（1662年）：吴三桂军将自缅甸押回的永历帝朱由榔处死，南明朝廷名义中心结束}
本节点追踪「吴三桂军将自缅甸押回的永历帝朱由榔处死，南明朝廷名义中心结束」之后可见的余波：\eventanchor{EQ-1662-YONGLI-EXECUTION-AFTERMATH}{吴三桂军将自缅甸押回的永历帝朱由榔处死，南明朝廷名义中心结束}。清、南明永历朝与缅甸在边地道路、驻防与多方社群交错的尺度的处境可由永历朝长期依赖流亡和边地补给，缅甸宫廷在清军压力下改变安置策略，吴三桂完成追捕说明。《清圣祖实录》康熙元年四月；《清史稿·永历帝本纪》；《琉璃宫史》缅文传统与《云南通志》是本段材料入口，因而只能据此追踪到南明皇帝名义不复存在，但西南残部、郑氏与地方记忆仍继续影响清初整合。原有置信注记：被押回和处死主干可靠；缅甸移交过程、处死地点与责任叙事须比对中缅材料

\paragraph{康熙元年四月（1662年）：吴三桂军将自缅甸押回的永历帝朱由榔处死，南明朝廷名义中心结束}
\eventanchor{EQ-1662-YONGLI-EXECUTION-DECISION}{吴三桂军将自缅甸押回的永历帝朱由榔处死，南明朝廷名义中心结束} 发生在边地道路、驻防与多方社群交错的尺度，牵动清、南明永历朝与缅甸。本节点定位围绕「吴三桂军将自缅甸押回的永历帝朱由榔处死，南明朝廷名义中心结束」形成的处置与调度记录。永历朝长期依赖流亡和边地补给，缅甸宫廷在清军压力下改变安置策略，吴三桂完成追捕是其结构性条件；《清圣祖实录》康熙元年四月；《清史稿·永历帝本纪》；《琉璃宫史》缅文传统与《云南通志》留下的证据把读者带到南明皇帝名义不复存在，但西南残部、郑氏与地方记忆仍继续影响清初整合。原有置信注记：被押回和处死主干可靠；缅甸移交过程、处死地点与责任叙事须比对中缅材料

\paragraph{康熙元年四月（1662年）：吴三桂军将自缅甸押回的永历帝朱由榔处死，南明朝廷名义中心结束}
\eventanchor{EQ-1662-YONGLI-EXECUTION-PRELUDE}{吴三桂军将自缅甸押回的永历帝朱由榔处死，南明朝廷名义中心结束} 对应本节点回看「吴三桂军将自缅甸押回的永历帝朱由榔处死，南明朝廷名义中心结束」之前已可见的条件。在边地道路、驻防与多方社群交错的尺度，清、南明永历朝与缅甸无法绕开永历朝长期依赖流亡和边地补给，缅甸宫廷在清军压力下改变安置策略，吴三桂完成追捕。依据《清圣祖实录》康熙元年四月；《清史稿·永历帝本纪》；《琉璃宫史》缅文传统与《云南通志》，本段把后续影响限定为南明皇帝名义不复存在，但西南残部、郑氏与地方记忆仍继续影响清初整合。原有置信注记：被押回和处死主干可靠；缅甸移交过程、处死地点与责任叙事须比对中缅材料

\paragraph{康熙三年（1664年）：清军攻破夔东据点，李来亨等抗清武装失败，长江上游的明遗民军事中心瓦解}
\eventanchor{EQ-1664-KUIZHOU}{清军攻破夔东据点，李来亨等抗清武装失败，长江上游的明遗民军事中心瓦解}。本节点叙述该事件本身。清与夔东抗清武装在京师与多地军政线路相接的尺度的处境首先受永历朝覆灭后夔东武装仍凭山地、峡江与地方网络抵抗，清廷需清理长江航运安全隐患制约。《清圣祖实录》康熙三年湖广四川军报；《清史稿·李来亨传》；《夔州府志》与南明遗民文集留下可核的线索；可辨的后果是清廷加强川东鄂西通道控制，残余反清力量更难与外部政权形成连线。原有置信注记：夔东结局可靠；各部关系、山地据点与居民伤亡材料零散

\paragraph{康熙三年（1664年）：清军攻破夔东据点，李来亨等抗清武装失败，长江上游的明遗民军事中心瓦解}
\eventanchor{EQ-1664-KUIZHOU-AFTERMATH}{清军攻破夔东据点，李来亨等抗清武装失败，长江上游的明遗民军事中心瓦解} 置于京师与多地军政线路相接的尺度中阅读，本节点追踪「清军攻破夔东据点，李来亨等抗清武装失败，长江上游的明遗民军事中心瓦解」之后可见的余波。永历朝覆灭后夔东武装仍凭山地、峡江与地方网络抵抗，清廷需清理长江航运安全隐患构成行动者清与夔东抗清武装必须面对的条件。取证从《清圣祖实录》康熙三年湖广四川军报；《清史稿·李来亨传》；《夔州府志》与南明遗民文集出发，因而只能把变化谨慎地连到清廷加强川东鄂西通道控制，残余反清力量更难与外部政权形成连线。原有置信注记：夔东结局可靠；各部关系、山地据点与居民伤亡材料零散

\paragraph{康熙三年（1664年）：清军攻破夔东据点，李来亨等抗清武装失败，长江上游的明遗民军事中心瓦解}
\eventanchor{EQ-1664-KUIZHOU-DECISION}{清军攻破夔东据点，李来亨等抗清武装失败，长江上游的明遗民军事中心瓦解} 标记本节点定位围绕「清军攻破夔东据点，李来亨等抗清武装失败，长江上游的明遗民军事中心瓦解」形成的处置与调度记录，涉及清与夔东抗清武装。京师与多地军政线路相接的尺度不是抽象背景：永历朝覆灭后夔东武装仍凭山地、峡江与地方网络抵抗，清廷需清理长江航运安全隐患。《清圣祖实录》康熙三年湖广四川军报；《清史稿·李来亨传》；《夔州府志》与南明遗民文集所见的顺序随后指向清廷加强川东鄂西通道控制，残余反清力量更难与外部政权形成连线，但原有置信注记：夔东结局可靠；各部关系、山地据点与居民伤亡材料零散

\paragraph{康熙三年（1664年）：清军攻破夔东据点，李来亨等抗清武装失败，长江上游的明遗民军事中心瓦解}
在京师与多地军政线路相接的尺度，\eventanchor{EQ-1664-KUIZHOU-PRELUDE}{清军攻破夔东据点，李来亨等抗清武装失败，长江上游的明遗民军事中心瓦解} 让清与夔东抗清武装的一个选择或压力有了可查的坐标。本节点回看「清军攻破夔东据点，李来亨等抗清武装失败，长江上游的明遗民军事中心瓦解」之前已可见的条件。前因可写为永历朝覆灭后夔东武装仍凭山地、峡江与地方网络抵抗，清廷需清理长江航运安全隐患；《清圣祖实录》康熙三年湖广四川军报；《清史稿·李来亨传》；《夔州府志》与南明遗民文集限定了记录，后续变化是清廷加强川东鄂西通道控制，残余反清力量更难与外部政权形成连线。原有置信注记：夔东结局可靠；各部关系、山地据点与居民伤亡材料零散

\begin{figure}[htbp]
  \centering
  \begin{tikzpicture}[
    x=1cm,y=1cm,
    place/.style={draw=timeline, fill=paper, rounded corners=2pt, align=center,
      inner sep=3pt, font=\small},
    qing/.style={place, draw=dynasty, fill=dynasty!13},
    ming/.style={place, draw=timeline, fill=timeline!13},
    note/.style={align=center, font=\scriptsize\color{source}}
  ]
    \fill[paper] (-.6,-3.9) rectangle (7.0,3.8);
    \draw[dynasty, line width=.8pt] (-.6,3.8) -- (7.0,3.8);
    \node[font=\bfseries\large, text=dynasty] at (3.2,3.35)
      {1644--1646：入关、北京与南方多中心};

    \node[qing] (shanhai) at (.5,2.35) {山海关\\1644};
    \node[qing] (beijing) at (2.25,1.7) {北京\\摄政与定都};
    \node[qing] (nanjing) at (3.65,.65) {南京\\弘光朝失守\\1645};
    \node[ming] (zhejiang) at (5.55,.45) {浙东、海上\\鲁王集团};
    \node[ming] (fuzhou) at (5.35,-1.25) {福州\\隆武朝\\1646};
    \node[ming] (guangzhou) at (2.55,-2.15) {广州\\绍武朝\\1646};
    \node[ming] (zhaoqing) at (.55,-3.0) {肇庆\\永历朝\\1646冬};

    \draw[->, dynasty, line width=1.15pt] (shanhai) -- node[above,sloped,note]
      {清军与吴三桂部的入关通道} (beijing);
    \draw[->, dynasty, line width=1.15pt] (beijing) -- node[above,sloped,note]
      {军行、接收与文书网络（示意）} (nanjing);
    \draw[->, dynasty, line width=1.15pt] (nanjing) -- node[above,note]
      {浙东、沿海战区} (zhejiang);
    \draw[->, dynasty, line width=1.15pt] (nanjing) -- node[right,note]
      {闽北战区} (fuzhou);
    \draw[->, dynasty, line width=1.15pt] (nanjing) -- node[left,note]
      {两广战区} (guangzhou);
    \draw[->, timeline, dashed, line width=.95pt] (guangzhou) -- node[below,sloped,note]
      {南明名义与人员流动（非疆界）} (zhaoqing);
    \draw[timeline, dashed] (fuzhou) -- (zhejiang);

    \node[draw=source, fill=paper, rounded corners=3pt, align=center,
      text=source, font=\small] at (4.55,-3.0)
      {城池、朝廷中心与路线\\不等于整片区域的持续控制};
    \draw[source, dashed] (4.1,-2.55) -- (3.05,-2.25);
  \end{tikzpicture}
  \caption{清军入关、北京接收与南方竞争性中心（1644--1646，示意）。朱褐实线为账本中可锚定的主要军行和接收通道；铜色虚线为南明中心及其联络的示意，均非连续控制线。依据《清世祖实录》顺治元年四、五、九月及顺治二、三年军报、《清初内国史院满文档案》关宁军务与北京迁都礼仪档（满、汉文行政链），并参照《甲申传信录》、江宁地方记事、《福建通志》《广东通志》、弘光朝奏疏与《永历实录》。范围仅及山海关--北京及江南、浙闽、两广的若干节点；路线不按比例，未绘政权边界，也不表示乡村、岛屿、山地或府县已被一方稳定控制。}
  \label{fig:qing-entry-competing-centers}
\end{figure}

\subsection{1653—1672：海上退守与皇权重组}

福建的港湾没有因为府县已换了旗号便归入同一种秩序。顺治十年（1653年）以后，郑成功以厦门、金门为泊船、转运和集结的支点，在漳州、泉州、同安之间反复争夺水陆相接的地方。清廷在南方取得城池和官署，首先意味着可以设置守备、征催钱粮和接收投诚；它并不使沿海的商人、船户、耕作者和旧有军政关系同时改变。郑氏所能倚赖的，也不是一条抽象的“海上生命线”，而是港口外可停泊的水面、通往腹地的渡口、愿意或被迫供给的村落，以及仍可传递消息的亲族和贸易关系。陆上的城防若能截住这些接点，岛屿便难以养兵；郑军若能保持它们，清方的府县接收便始终有边界。

顺治十六年（1659年）的北上南京，把这种未竟的南方接收推到长江下游。郑军试图把海上的机动与江南的明遗民联系、粮赋资源结合起来，却须在内河水道、城外营地和城防之间维持船只、粮食与联络。撤退并非只是一场战役的结果：它显示一支能够穿入长江的舰队，仍难以在远离岛屿据点的地方承受围城和久驻的耗费。南京周围的行动不能用作江南人心的总表；有人提供消息、有人避战、有人被征发，也有人首先面对的是保城、保粮或逃离的选择。清方战报、杨英《从征实录》、地方记事和荷兰方面的海上消息能够彼此校正北伐与撤军的大致次序，却不足以把内应的范围、兵员损失或每座村镇的立场写成整齐数字。

因此，1661年的迁界令与郑成功登陆台湾必须放在同一张水陆网络中理解。清廷命闽粤居民内迁，所要切断的是郑氏从岸上获得粮食、人口和情报的可能；这是一种以迁徙制造军事距离的省级治理，并不是海面从此封闭的证明。谕旨、督抚文书和地方志可看见命令的层层下达、若干撤迁与后来的复业线索，却不能由此推算整段海岸的迁徙距离、死亡人数或秘密贸易规模。渔盐、田地、渡海、家族和市集原本相互牵连，禁令把风险分配得并不均匀：一处港口的沉寂不能代表另一处，留在原地、迁往内陆和改换生计的人，也不会在官府档案中留下同样清楚的痕迹。

台湾在此时不是南京撤军后自然出现的“退路”，而是争夺粮食、人口和航路的另一处战场。郑军围攻热兰遮城，至1662年初荷兰守军投降；城堡、贸易节点和周边土地随之易手。围城使荷兰东印度公司的记录保存了守军补给、城内决策和交接的一个视角，投降条约也能钉住转换的时点；它们都不能代替城外汉人居民、被俘者或原住民社群的叙述。郑成功旋即去世，继承之争又使新基地必须重新安置军政、宗族和海商网络。1663年厦门、金门易手，郑氏主力转守台湾，海峡两边的压力并未消失，只是更多地落到岛内土地、港口与跨海经营所能供给的限度上。

北京的权力重组同样要穿过这些地方性的供给和文书链。1661年顺治帝去世，八岁的玄烨即位，索尼、苏克萨哈、遏必隆、鳌拜辅政；1665年前后鳌拜影响扩大，1667年索尼去世，原有的均衡更难维持。辅政者处置的并非只是一座宫廷中的职位，还包括旗务、人事、军费以及各省报上来的战情和钱粮。1669年康熙拘捕鳌拜，扩大了皇帝直接统摄军政人事的条件，却不能倒过来证明此后每一项调拨都畅通无阻。1669年案件的案卷和后出的传记尤其容易把结局写成早已注定的个人胜负；满汉题本与起居注留下的议事程序，反而提醒我们当时的约束并未因后来处分而消失。

1670年颁行十六条圣谕，朝廷尝试把伦理、赋役和治安的语言交给地方官绅与社区。这与沿海隔离和战时转运一样，显示中央试图以可传达的规范接续有限的官僚、驻军和财力；但诏令的颁布只能证明国家说了什么，不能证明某县怎样宣讲、何人听见，或社会秩序因而如何改变。各省方志、基层诉讼和契约偶尔保存这类接触的局部痕迹，保存的又多是有文书能力的官、绅与产权关系，不能替无名的役夫、流徙者和船户发言。到1672年，南方的接收、海峡的战争和省级的筹饷仍未收束为稳定秩序；它们留下的驻军、道路、债务与地方中介关系，也构成次年撤藩冲突得以扩大的条件。

\paragraph{材料位置。} 《康熙起居注》、满汉题本和鳌拜案档案呈现议事、处分与命令的文书链；《清史稿》的后出人物叙事则带有1669年以后胜利的视角。沿海一侧，实录和军报须与郑氏文书、地方志、投降条约及荷兰东印度公司热兰遮城日志并读：它们能够连接政策、围城和港口，却各自遮蔽不同社群。中央重组与海岸战争共享财政、交通和信息空间，仍不能彼此化约。

\section{摄政的终结与皇权的重新组织}

古北口外的行猎行程在顺治七年十二月突然中断：多尔衮去世，原先围绕摄政王组织的军政枢纽随之改组（EQ-1650-DORGON-DEATH）。《清世祖实录》和内国史院的丧仪、遗命档可以确认死期和制度性的后果；死因以及死后清算的动机，则不能照抄后来人物传记的定性。次年顺治帝亲政，撤去摄政安排、重置皇帝周边的任官和决策网络（EQ-1651-SHUNZHI-PERSONAL-RULE）。亲政诏令和题本首先证明命令的发出和人事的程序，并不证明南方将领、降官、旗务与部院已经无条件服从。皇帝名义的强化仍须经过同一套军饷、道路和文书链才能抵达地方。

顺治十八年，八岁的玄烨即位，索尼、苏克萨哈、遏必隆、鳌拜辅政（EQ-1661-KANGXI-ACCESSION）。继位和四人名单可由《清世祖实录》、起居注、内国史院遗诏档交叉确立；顺治死因与遗诏形成过程仍有材料和后世传说的层次差异。康熙四年辅政大臣依议政与旗务网络处置军政，鳌拜的影响扩大（EQ-1665-REGENTS-CONSOLIDATION）；索尼于1667年去世，使原有均衡更加脆弱（EQ-1667-SONI-DEATH）。《康熙起居注》与满汉题本呈现的是当时的议事和程序，而《清史稿·鳌拜传》及鳌拜案档案带着1669年胜利后的视角。把后来被拘捕的结局提前写成“独揽大权”，会抹掉辅政本来具有的集体性和相互制约。

康熙八年五月拘捕鳌拜并处分其党羽，的确改变了皇帝直接统摄人事与军政的条件（EQ-1669-OBOI-ARREST）。拘捕、审理和处分可由实录、起居注及内阁案卷相互验证；宫中擒捕的戏剧性细节、罪状的完整性和少年皇帝的单独谋略，则都是胜利叙事最容易放大的部分。康熙九年颁布十六条圣谕（EQ-1670-SACRED-EDICT），又把中央的伦理、赋役和秩序语言交给地方官绅与社区。圣谕的文本和宣讲制度是可见的，识读范围、宣讲频率与社会效果却必须由各省方志、基层文书和具体时段检验。它说明国家试图以训导补足武力和官僚覆盖的不足，而不是一条可以直接量得的服从曲线。


\subsection{事件链与材料限度}
摄政、宗室、部院与地方官之间的文书和人事安排保留了不同速度。

\paragraph{顺治七年十二月（1650年）：摄政王多尔衮在古北口外行猎途中去世，摄政政治结构随之改组}
\eventanchor{EQ-1650-DORGON-DEATH}{摄政王多尔衮在古北口外行猎途中去世，摄政政治结构随之改组}。本节点叙述该事件本身。清在京师与多地军政线路相接的尺度的处境首先受多尔衮集中军政与皇室事务，围绕摄政权的宗室和旗人竞争在其死后失去调停中心制约。《清世祖实录》顺治七年十二月；《清史稿·多尔衮传》；《清初内国史院满文档案》摄政王丧仪与遗命档留下可核的线索；可辨的后果是顺治帝及近臣得以重组朝廷，摄政集团的名誉与权力被重新界定。原有置信注记：死亡日期可靠；死因、身后清算与权力斗争不能只依后来的政治定性

\paragraph{顺治七年十二月（1650年）：摄政王多尔衮在古北口外行猎途中去世，摄政政治结构随之改组}
\eventanchor{EQ-1650-DORGON-DEATH-AFTERMATH}{摄政王多尔衮在古北口外行猎途中去世，摄政政治结构随之改组} 置于京师与多地军政线路相接的尺度中阅读，本节点追踪「摄政王多尔衮在古北口外行猎途中去世，摄政政治结构随之改组」之后可见的余波。多尔衮集中军政与皇室事务，围绕摄政权的宗室和旗人竞争在其死后失去调停中心构成行动者清必须面对的条件。取证从《清世祖实录》顺治七年十二月；《清史稿·多尔衮传》；《清初内国史院满文档案》摄政王丧仪与遗命档出发，因而只能把变化谨慎地连到顺治帝及近臣得以重组朝廷，摄政集团的名誉与权力被重新界定。原有置信注记：死亡日期可靠；死因、身后清算与权力斗争不能只依后来的政治定性

\paragraph{顺治七年十二月（1650年）：摄政王多尔衮在古北口外行猎途中去世，摄政政治结构随之改组}
\eventanchor{EQ-1650-DORGON-DEATH-DECISION}{摄政王多尔衮在古北口外行猎途中去世，摄政政治结构随之改组} 标记本节点定位围绕「摄政王多尔衮在古北口外行猎途中去世，摄政政治结构随之改组」形成的处置与调度记录，涉及清。京师与多地军政线路相接的尺度不是抽象背景：多尔衮集中军政与皇室事务，围绕摄政权的宗室和旗人竞争在其死后失去调停中心。《清世祖实录》顺治七年十二月；《清史稿·多尔衮传》；《清初内国史院满文档案》摄政王丧仪与遗命档所见的顺序随后指向顺治帝及近臣得以重组朝廷，摄政集团的名誉与权力被重新界定，但原有置信注记：死亡日期可靠；死因、身后清算与权力斗争不能只依后来的政治定性

\paragraph{顺治七年十二月（1650年）：摄政王多尔衮在古北口外行猎途中去世，摄政政治结构随之改组}
在京师与多地军政线路相接的尺度，\eventanchor{EQ-1650-DORGON-DEATH-PRELUDE}{摄政王多尔衮在古北口外行猎途中去世，摄政政治结构随之改组} 让清的一个选择或压力有了可查的坐标。本节点回看「摄政王多尔衮在古北口外行猎途中去世，摄政政治结构随之改组」之前已可见的条件。前因可写为多尔衮集中军政与皇室事务，围绕摄政权的宗室和旗人竞争在其死后失去调停中心；《清世祖实录》顺治七年十二月；《清史稿·多尔衮传》；《清初内国史院满文档案》摄政王丧仪与遗命档限定了记录，后续变化是顺治帝及近臣得以重组朝廷，摄政集团的名誉与权力被重新界定。原有置信注记：死亡日期可靠；死因、身后清算与权力斗争不能只依后来的政治定性

\paragraph{顺治八年（1651年）：顺治帝亲政，撤除摄政安排并重置皇帝周边的决策网络}
\eventanchor{EQ-1651-SHUNZHI-PERSONAL-RULE}{顺治帝亲政，撤除摄政安排并重置皇帝周边的决策网络} 所关联的是清在京师与多地军政线路相接的尺度中的行动。本节点叙述该事件本身。它从多尔衮去世使摄政权威无法延续，皇帝及支持者以礼制和任免重建指挥链展开，借《清世祖实录》顺治八年正月；《清史稿·世祖本纪》；《顺治朝内阁档》亲政诏令与任官题本进入可检验的年序，并以中央命令署名和人事结构改变，但南方征服与财政仍要求依靠地方将领和降官影响下一段安排。原有置信注记：亲政事实可靠；实际决策仍受宗室、旗务和满汉官僚共同制约

\paragraph{顺治八年（1651年）：顺治帝亲政，撤除摄政安排并重置皇帝周边的决策网络}
京师与多地军政线路相接的尺度中的\eventanchor{EQ-1651-SHUNZHI-PERSONAL-RULE-AFTERMATH}{顺治帝亲政，撤除摄政安排并重置皇帝周边的决策网络} 不能离开清来解释。本节点追踪「顺治帝亲政，撤除摄政安排并重置皇帝周边的决策网络」之后可见的余波。多尔衮去世使摄政权威无法延续，皇帝及支持者以礼制和任免重建指挥链说明其形成的条件；《清世祖实录》顺治八年正月；《清史稿·世祖本纪》；《顺治朝内阁档》亲政诏令与任官题本提供来源路径，中央命令署名和人事结构改变，但南方征服与财政仍要求依靠地方将领和降官则是材料能够支持的近时结果。原有置信注记：亲政事实可靠；实际决策仍受宗室、旗务和满汉官僚共同制约

\paragraph{顺治八年（1651年）：顺治帝亲政，撤除摄政安排并重置皇帝周边的决策网络}
\eventanchor{EQ-1651-SHUNZHI-PERSONAL-RULE-DECISION}{顺治帝亲政，撤除摄政安排并重置皇帝周边的决策网络} 把清放回京师与多地军政线路相接的尺度的道路、城镇、边界或制度网络。本节点定位围绕「顺治帝亲政，撤除摄政安排并重置皇帝周边的决策网络」形成的处置与调度记录。多尔衮去世使摄政权威无法延续，皇帝及支持者以礼制和任免重建指挥链。本段采用《清世祖实录》顺治八年正月；《清史稿·世祖本纪》；《顺治朝内阁档》亲政诏令与任官题本核对其顺序，随后可追到中央命令署名和人事结构改变，但南方征服与财政仍要求依靠地方将领和降官。原有置信注记：亲政事实可靠；实际决策仍受宗室、旗务和满汉官僚共同制约

\paragraph{顺治八年（1651年）：顺治帝亲政，撤除摄政安排并重置皇帝周边的决策网络}
本节点回看「顺治帝亲政，撤除摄政安排并重置皇帝周边的决策网络」之前已可见的条件，故\eventanchor{EQ-1651-SHUNZHI-PERSONAL-RULE-PRELUDE}{顺治帝亲政，撤除摄政安排并重置皇帝周边的决策网络} 不只是一个年份标签。清在京师与多地军政线路相接的尺度承受的压力来自多尔衮去世使摄政权威无法延续，皇帝及支持者以礼制和任免重建指挥链。《清世祖实录》顺治八年正月；《清史稿·世祖本纪》；《顺治朝内阁档》亲政诏令与任官题本可回查该节点；中央命令署名和人事结构改变，但南方征服与财政仍要求依靠地方将领和降官是其进入后续年序的方式。原有置信注记：亲政事实可靠；实际决策仍受宗室、旗务和满汉官僚共同制约

\paragraph{顺治十八年正月（1661年）：顺治帝去世，八岁玄烨即位为康熙帝，索尼、苏克萨哈、遏必隆、鳌拜辅政}
\eventanchor{EQ-1661-KANGXI-ACCESSION}{顺治帝去世，八岁玄烨即位为康熙帝，索尼、苏克萨哈、遏必隆、鳌拜辅政} 所见的行动者为清，空间尺度为京师与多地军政线路相接的尺度。本节点叙述该事件本身。顺治早逝使未成年皇帝继承，宗室与旗人高层需建立可承认的集体决策安排使它成为具体事件链的一环；《清世祖实录》顺治十八年正月；《清史稿·圣祖本纪》；《康熙起居注》与内国史院遗诏档与辅政集团掌握朝政，皇权成长、旗人权力平衡和南方战事资源相互交织。原有置信注记：继位及辅政名单可靠；顺治死因和遗诏形成过程存在材料与后世叙事差异

\paragraph{顺治十八年正月（1661年）：顺治帝去世，八岁玄烨即位为康熙帝，索尼、苏克萨哈、遏必隆、鳌拜辅政}
从京师与多地军政线路相接的尺度看，\eventanchor{EQ-1661-KANGXI-ACCESSION-AFTERMATH}{顺治帝去世，八岁玄烨即位为康熙帝，索尼、苏克萨哈、遏必隆、鳌拜辅政} 留下本节点追踪「顺治帝去世，八岁玄烨即位为康熙帝，索尼、苏克萨哈、遏必隆、鳌拜辅政」之后可见的余波的材料坐标。清所面对的条件是顺治早逝使未成年皇帝继承，宗室与旗人高层需建立可承认的集体决策安排。《清世祖实录》顺治十八年正月；《清史稿·圣祖本纪》；《康熙起居注》与内国史院遗诏档支持的结果为辅政集团掌握朝政，皇权成长、旗人权力平衡和南方战事资源相互交织。原有置信注记：继位及辅政名单可靠；顺治死因和遗诏形成过程存在材料与后世叙事差异

\paragraph{顺治十八年正月（1661年）：顺治帝去世，八岁玄烨即位为康熙帝，索尼、苏克萨哈、遏必隆、鳌拜辅政}
\eventanchor{EQ-1661-KANGXI-ACCESSION-DECISION}{顺治帝去世，八岁玄烨即位为康熙帝，索尼、苏克萨哈、遏必隆、鳌拜辅政} 记录清在京师与多地军政线路相接的尺度的一次变化。本节点定位围绕「顺治帝去世，八岁玄烨即位为康熙帝，索尼、苏克萨哈、遏必隆、鳌拜辅政」形成的处置与调度记录。其发生与顺治早逝使未成年皇帝继承，宗室与旗人高层需建立可承认的集体决策安排相连；《清世祖实录》顺治十八年正月；《清史稿·圣祖本纪》；《康熙起居注》与内国史院遗诏档把事件系回来源，辅政集团掌握朝政，皇权成长、旗人权力平衡和南方战事资源相互交织显示压力如何向后转移。原有置信注记：继位及辅政名单可靠；顺治死因和遗诏形成过程存在材料与后世叙事差异

\paragraph{顺治十八年正月（1661年）：顺治帝去世，八岁玄烨即位为康熙帝，索尼、苏克萨哈、遏必隆、鳌拜辅政}
\eventanchor{EQ-1661-KANGXI-ACCESSION-PRELUDE}{顺治帝去世，八岁玄烨即位为康熙帝，索尼、苏克萨哈、遏必隆、鳌拜辅政} 的地点与行动范围属于京师与多地军政线路相接的尺度，行动者是清。本节点回看「顺治帝去世，八岁玄烨即位为康熙帝，索尼、苏克萨哈、遏必隆、鳌拜辅政」之前已可见的条件。顺治早逝使未成年皇帝继承，宗室与旗人高层需建立可承认的集体决策安排。读《清世祖实录》顺治十八年正月；《清史稿·圣祖本纪》；《康熙起居注》与内国史院遗诏档时可见其记录边界，最小后果只能写作辅政集团掌握朝政，皇权成长、旗人权力平衡和南方战事资源相互交织。原有置信注记：继位及辅政名单可靠；顺治死因和遗诏形成过程存在材料与后世叙事差异

\paragraph{康熙四年（1665年）：四辅政大臣以议政和旗务网络处理军政事务，鳌拜在政务中的影响扩大}
\eventanchor{EQ-1665-REGENTS-CONSOLIDATION}{四辅政大臣以议政和旗务网络处理军政事务，鳌拜在政务中的影响扩大} 是清在京师与多地军政线路相接的尺度留下的编年标记。本节点叙述该事件本身。它从幼帝尚未亲裁，辅政大臣需在旗人贵族、部院官僚和南方军费之间协调走来；《清圣祖实录》康熙四年政务条；《清史稿·鳌拜传》；《康熙起居注》及内阁大库满汉题本固定可证部分，中央决策呈现集体辅政与个人网络并存，为1669年权力重组提供背景构成下一步的可见结果。原有置信注记：辅政架构可靠；权力高下多由后来的鳌拜案档与《清史稿》倒叙，须避免预设结局

\paragraph{康熙四年（1665年）：四辅政大臣以议政和旗务网络处理军政事务，鳌拜在政务中的影响扩大}
本节点追踪「四辅政大臣以议政和旗务网络处理军政事务，鳌拜在政务中的影响扩大」之后可见的余波：\eventanchor{EQ-1665-REGENTS-CONSOLIDATION-AFTERMATH}{四辅政大臣以议政和旗务网络处理军政事务，鳌拜在政务中的影响扩大}。清在京师与多地军政线路相接的尺度的处境可由幼帝尚未亲裁，辅政大臣需在旗人贵族、部院官僚和南方军费之间协调说明。《清圣祖实录》康熙四年政务条；《清史稿·鳌拜传》；《康熙起居注》及内阁大库满汉题本是本段材料入口，因而只能据此追踪到中央决策呈现集体辅政与个人网络并存，为1669年权力重组提供背景。原有置信注记：辅政架构可靠；权力高下多由后来的鳌拜案档与《清史稿》倒叙，须避免预设结局

\paragraph{康熙四年（1665年）：四辅政大臣以议政和旗务网络处理军政事务，鳌拜在政务中的影响扩大}
\eventanchor{EQ-1665-REGENTS-CONSOLIDATION-DECISION}{四辅政大臣以议政和旗务网络处理军政事务，鳌拜在政务中的影响扩大} 发生在京师与多地军政线路相接的尺度，牵动清。本节点定位围绕「四辅政大臣以议政和旗务网络处理军政事务，鳌拜在政务中的影响扩大」形成的处置与调度记录。幼帝尚未亲裁，辅政大臣需在旗人贵族、部院官僚和南方军费之间协调是其结构性条件；《清圣祖实录》康熙四年政务条；《清史稿·鳌拜传》；《康熙起居注》及内阁大库满汉题本留下的证据把读者带到中央决策呈现集体辅政与个人网络并存，为1669年权力重组提供背景。原有置信注记：辅政架构可靠；权力高下多由后来的鳌拜案档与《清史稿》倒叙，须避免预设结局

\paragraph{康熙四年（1665年）：四辅政大臣以议政和旗务网络处理军政事务，鳌拜在政务中的影响扩大}
\eventanchor{EQ-1665-REGENTS-CONSOLIDATION-PRELUDE}{四辅政大臣以议政和旗务网络处理军政事务，鳌拜在政务中的影响扩大} 对应本节点回看「四辅政大臣以议政和旗务网络处理军政事务，鳌拜在政务中的影响扩大」之前已可见的条件。在京师与多地军政线路相接的尺度，清无法绕开幼帝尚未亲裁，辅政大臣需在旗人贵族、部院官僚和南方军费之间协调。依据《清圣祖实录》康熙四年政务条；《清史稿·鳌拜传》；《康熙起居注》及内阁大库满汉题本，本段把后续影响限定为中央决策呈现集体辅政与个人网络并存，为1669年权力重组提供背景。原有置信注记：辅政架构可靠；权力高下多由后来的鳌拜案档与《清史稿》倒叙，须避免预设结局

\paragraph{康熙六年（1667年）：首辅索尼去世，四辅政大臣的均衡被打破}
\eventanchor{EQ-1667-SONI-DEATH}{首辅索尼去世，四辅政大臣的均衡被打破}。本节点叙述该事件本身。清在京师与多地军政线路相接的尺度的处境首先受索尼居于辅政资历与旗人关系枢纽，其去世改变了原有制衡条件制约。《清圣祖实录》康熙六年；《清史稿·索尼传》；《康熙起居注》康熙六年丧仪与议政记录留下可核的线索；可辨的后果是鳌拜及盟友影响更显著，康熙帝近侍和太皇太后方面开始寻找收回裁决权的机会。原有置信注记：索尼死亡及官职变动可靠；其政治立场与遗言多经后出人物传记修饰

\paragraph{康熙六年（1667年）：首辅索尼去世，四辅政大臣的均衡被打破}
\eventanchor{EQ-1667-SONI-DEATH-AFTERMATH}{首辅索尼去世，四辅政大臣的均衡被打破} 置于京师与多地军政线路相接的尺度中阅读，本节点追踪「首辅索尼去世，四辅政大臣的均衡被打破」之后可见的余波。索尼居于辅政资历与旗人关系枢纽，其去世改变了原有制衡条件构成行动者清必须面对的条件。取证从《清圣祖实录》康熙六年；《清史稿·索尼传》；《康熙起居注》康熙六年丧仪与议政记录出发，因而只能把变化谨慎地连到鳌拜及盟友影响更显著，康熙帝近侍和太皇太后方面开始寻找收回裁决权的机会。原有置信注记：索尼死亡及官职变动可靠；其政治立场与遗言多经后出人物传记修饰

\paragraph{康熙六年（1667年）：首辅索尼去世，四辅政大臣的均衡被打破}
\eventanchor{EQ-1667-SONI-DEATH-DECISION}{首辅索尼去世，四辅政大臣的均衡被打破} 标记本节点定位围绕「首辅索尼去世，四辅政大臣的均衡被打破」形成的处置与调度记录，涉及清。京师与多地军政线路相接的尺度不是抽象背景：索尼居于辅政资历与旗人关系枢纽，其去世改变了原有制衡条件。《清圣祖实录》康熙六年；《清史稿·索尼传》；《康熙起居注》康熙六年丧仪与议政记录所见的顺序随后指向鳌拜及盟友影响更显著，康熙帝近侍和太皇太后方面开始寻找收回裁决权的机会，但原有置信注记：索尼死亡及官职变动可靠；其政治立场与遗言多经后出人物传记修饰

\paragraph{康熙六年（1667年）：首辅索尼去世，四辅政大臣的均衡被打破}
在京师与多地军政线路相接的尺度，\eventanchor{EQ-1667-SONI-DEATH-PRELUDE}{首辅索尼去世，四辅政大臣的均衡被打破} 让清的一个选择或压力有了可查的坐标。本节点回看「首辅索尼去世，四辅政大臣的均衡被打破」之前已可见的条件。前因可写为索尼居于辅政资历与旗人关系枢纽，其去世改变了原有制衡条件；《清圣祖实录》康熙六年；《清史稿·索尼传》；《康熙起居注》康熙六年丧仪与议政记录限定了记录，后续变化是鳌拜及盟友影响更显著，康熙帝近侍和太皇太后方面开始寻找收回裁决权的机会。原有置信注记：索尼死亡及官职变动可靠；其政治立场与遗言多经后出人物传记修饰

\paragraph{康熙八年五月（1669年）：康熙帝拘捕鳌拜并审理其党羽，结束辅政大臣实际主政的局面}
\eventanchor{EQ-1669-OBOI-ARREST}{康熙帝拘捕鳌拜并审理其党羽，结束辅政大臣实际主政的局面} 所关联的是清在京师与多地军政线路相接的尺度中的行动。本节点叙述该事件本身。它从索尼死后辅政均衡失去支点，鳌拜与苏克萨哈案件激化皇帝近臣、太皇太后和旗人贵族冲突展开，借《清圣祖实录》康熙八年五月；《清史稿·鳌拜传》；《康熙起居注》与鳌拜案内阁档进入可检验的年序，并以皇帝可直接统摄人事军政，但仍须通过议政、内阁及旗务结构行使权力影响下一段安排。原有置信注记：拘捕与处分可靠；宫中擒捕过程、罪状与少年皇帝个人谋略含胜利者叙事

\paragraph{康熙八年五月（1669年）：康熙帝拘捕鳌拜并审理其党羽，结束辅政大臣实际主政的局面}
京师与多地军政线路相接的尺度中的\eventanchor{EQ-1669-OBOI-ARREST-AFTERMATH}{康熙帝拘捕鳌拜并审理其党羽，结束辅政大臣实际主政的局面} 不能离开清来解释。本节点追踪「康熙帝拘捕鳌拜并审理其党羽，结束辅政大臣实际主政的局面」之后可见的余波。索尼死后辅政均衡失去支点，鳌拜与苏克萨哈案件激化皇帝近臣、太皇太后和旗人贵族冲突说明其形成的条件；《清圣祖实录》康熙八年五月；《清史稿·鳌拜传》；《康熙起居注》与鳌拜案内阁档提供来源路径，皇帝可直接统摄人事军政，但仍须通过议政、内阁及旗务结构行使权力则是材料能够支持的近时结果。原有置信注记：拘捕与处分可靠；宫中擒捕过程、罪状与少年皇帝个人谋略含胜利者叙事

\paragraph{康熙八年五月（1669年）：康熙帝拘捕鳌拜并审理其党羽，结束辅政大臣实际主政的局面}
\eventanchor{EQ-1669-OBOI-ARREST-DECISION}{康熙帝拘捕鳌拜并审理其党羽，结束辅政大臣实际主政的局面} 把清放回京师与多地军政线路相接的尺度的道路、城镇、边界或制度网络。本节点定位围绕「康熙帝拘捕鳌拜并审理其党羽，结束辅政大臣实际主政的局面」形成的处置与调度记录。索尼死后辅政均衡失去支点，鳌拜与苏克萨哈案件激化皇帝近臣、太皇太后和旗人贵族冲突。本段采用《清圣祖实录》康熙八年五月；《清史稿·鳌拜传》；《康熙起居注》与鳌拜案内阁档核对其顺序，随后可追到皇帝可直接统摄人事军政，但仍须通过议政、内阁及旗务结构行使权力。原有置信注记：拘捕与处分可靠；宫中擒捕过程、罪状与少年皇帝个人谋略含胜利者叙事

\paragraph{康熙八年五月（1669年）：康熙帝拘捕鳌拜并审理其党羽，结束辅政大臣实际主政的局面}
本节点回看「康熙帝拘捕鳌拜并审理其党羽，结束辅政大臣实际主政的局面」之前已可见的条件，故\eventanchor{EQ-1669-OBOI-ARREST-PRELUDE}{康熙帝拘捕鳌拜并审理其党羽，结束辅政大臣实际主政的局面} 不只是一个年份标签。清在京师与多地军政线路相接的尺度承受的压力来自索尼死后辅政均衡失去支点，鳌拜与苏克萨哈案件激化皇帝近臣、太皇太后和旗人贵族冲突。《清圣祖实录》康熙八年五月；《清史稿·鳌拜传》；《康熙起居注》与鳌拜案内阁档可回查该节点；皇帝可直接统摄人事军政，但仍须通过议政、内阁及旗务结构行使权力是其进入后续年序的方式。原有置信注记：拘捕与处分可靠；宫中擒捕过程、罪状与少年皇帝个人谋略含胜利者叙事

\paragraph{康熙九年（1670年）：康熙帝颁布十六条圣谕，规定由地方官和民众遵循的伦理、赋役与秩序训导框架}
\eventanchor{EQ-1670-SACRED-EDICT}{康熙帝颁布十六条圣谕，规定由地方官和民众遵循的伦理、赋役与秩序训导框架} 所见的行动者为清，空间尺度为省、府县与地方社会相互牵动的尺度。本节点叙述该事件本身。清廷在征服后寻求以共同伦理语言补足武力和官僚的有限覆盖，地方官绅被要求传达使它成为具体事件链的一环；《清圣祖实录》康熙九年十月；《清史稿·圣祖本纪》；《大清会典事例》圣谕条与各省地方志宣讲记录与圣谕成为后续地方教化与政治合法性资源，其接受和改写具有地区差异。原有置信注记：颁布可靠；宣讲频率、识读范围和实际社会效果不能由文本证明

\paragraph{康熙九年（1670年）：康熙帝颁布十六条圣谕，规定由地方官和民众遵循的伦理、赋役与秩序训导框架}
从省、府县与地方社会相互牵动的尺度看，\eventanchor{EQ-1670-SACRED-EDICT-AFTERMATH}{康熙帝颁布十六条圣谕，规定由地方官和民众遵循的伦理、赋役与秩序训导框架} 留下本节点追踪「康熙帝颁布十六条圣谕，规定由地方官和民众遵循的伦理、赋役与秩序训导框架」之后可见的余波的材料坐标。清所面对的条件是清廷在征服后寻求以共同伦理语言补足武力和官僚的有限覆盖，地方官绅被要求传达。《清圣祖实录》康熙九年十月；《清史稿·圣祖本纪》；《大清会典事例》圣谕条与各省地方志宣讲记录支持的结果为圣谕成为后续地方教化与政治合法性资源，其接受和改写具有地区差异。原有置信注记：颁布可靠；宣讲频率、识读范围和实际社会效果不能由文本证明

\paragraph{康熙九年（1670年）：康熙帝颁布十六条圣谕，规定由地方官和民众遵循的伦理、赋役与秩序训导框架}
\eventanchor{EQ-1670-SACRED-EDICT-DECISION}{康熙帝颁布十六条圣谕，规定由地方官和民众遵循的伦理、赋役与秩序训导框架} 记录清在省、府县与地方社会相互牵动的尺度的一次变化。本节点定位围绕「康熙帝颁布十六条圣谕，规定由地方官和民众遵循的伦理、赋役与秩序训导框架」形成的处置与调度记录。其发生与清廷在征服后寻求以共同伦理语言补足武力和官僚的有限覆盖，地方官绅被要求传达相连；《清圣祖实录》康熙九年十月；《清史稿·圣祖本纪》；《大清会典事例》圣谕条与各省地方志宣讲记录把事件系回来源，圣谕成为后续地方教化与政治合法性资源，其接受和改写具有地区差异显示压力如何向后转移。原有置信注记：颁布可靠；宣讲频率、识读范围和实际社会效果不能由文本证明

\paragraph{康熙九年（1670年）：康熙帝颁布十六条圣谕，规定由地方官和民众遵循的伦理、赋役与秩序训导框架}
\eventanchor{EQ-1670-SACRED-EDICT-PRELUDE}{康熙帝颁布十六条圣谕，规定由地方官和民众遵循的伦理、赋役与秩序训导框架} 的地点与行动范围属于省、府县与地方社会相互牵动的尺度，行动者是清。本节点回看「康熙帝颁布十六条圣谕，规定由地方官和民众遵循的伦理、赋役与秩序训导框架」之前已可见的条件。清廷在征服后寻求以共同伦理语言补足武力和官僚的有限覆盖，地方官绅被要求传达。读《清圣祖实录》康熙九年十月；《清史稿·圣祖本纪》；《大清会典事例》圣谕条与各省地方志宣讲记录时可见其记录边界，最小后果只能写作圣谕成为后续地方教化与政治合法性资源，其接受和改写具有地区差异。原有置信注记：颁布可靠；宣讲频率、识读范围和实际社会效果不能由文本证明

\section{海岸不是陆上征服的边缘}

厦门与金门之间的水道并不因为陆上府县易手而安静下来。顺治十年（1653年），郑成功以两岛为支点，把军队、商船和仍可调动的沿海关系拢在同一套经营中。岛屿提供停泊与转运的地点，港口把粮食、人员和消息送向舰队；但要把这些东西送到前线，还得越过潮汐、水道与清军据守的城镇。郑氏因此不能只守海面：漳州、泉州、同安一线的攻守，关系到腹地粮食能否抵达海边，也关系到城防能否把这条补给线截断。

顺治十一年至十三年间，围攻漳州以及在泉州、同安附近争夺水陆通道，正是这种压力的展开。郑军需要扩大粮饷与招募的基础，清军则依托府县城防和陆路援兵，使其难以迅速取得大陆府城。对港口周边的人们而言，战争不是两支舰队在远海相遇：征粮、避乱、治安的变化，会落在村落、渡口、盐场和市集不同的位置。我们可以从军报、地方志与郑氏文书看见攻守和若干地点，却很难据此排出每一处居民如何选择、何时离开或继续交易的整齐次序。商人、船民、耕作者、守城者和逃亡者面对的风险不必相同，更不能被概括成一种一致的“支持”或“抵抗”。

这也是郑氏海上力量的两面。商船和侨民网络能够输送筹饷的机会，岛屿也能让军政机构避开陆上包围；可是船只、兵员和收入的数字多见于敌对军报或后出的郑氏叙事，不能把它们当作精确的家底。较为稳妥的判断是：福建海疆已无法仅由陆上州县来控制。清廷后来把海禁、迁界与水师建设连在一起，并不是出于抽象的禁海观念，而是要处理港口、岛屿、腹地与消息网络彼此相连的难题。

顺治十六年（1659年）北上南京，使这种难题沿江而上。郑成功试图利用长江下游的粮赋与明遗民网络，把海上基地转化为更大的反清战场。围城并不只考验冲锋；军队要在内河环境中维持粮食、船只、城外营地和联络，清军则集中江防、城门与守备力量。郑军未能突破城防而撤回，显示其攻势在内河攻坚与长期补给面前难以持续。清廷记录、杨英《从征实录》与荷兰海上报告都能确认北伐、围城和撤退的主干，但内应的范围、作战计划及伤亡互有歧异；把这次行动写成几乎成功的复明，或写成毫无准备的冒险，都会超出材料所能承担的程度。

撤退没有使台湾自动成为唯一的去处，却使那里成为更迫切的战略选择。顺治十八年（1661年），清廷命福建、广东沿海迁界，意在把郑氏同岸上的粮食、人口与信息来源隔开。谕令所能直接告诉我们的，是中央选择了以强制迁徙制造隔离带；它不能替我们证明每一县迁了多远、多少人死于途中，或秘密贸易有多大规模。沿海县志和通志留下的迁撤、复界线索，使渔盐、耕作、宗族与港口关系被打断的代价可以进入叙事；但这些记载的密度和视角各异，不能把局部遭遇拼成一幅均一的海岸图景。

迁界也并未把海面与陆地彻底分开。禁令增加了往来成本，海岸居民、地方官和海上经营者仍可能在不同地点以不同方式应对：有人被迫内迁，有人守着可耕之地或亲族关系，有人继续依赖港口和航路。现有材料足以说明隔离带未能完全阻绝往来，却不足以量化这些变通，更不足以把它们叙述成全海岸一致、持续不断的走私网络。政策的直接军事目的，是压缩郑氏补给的可能；它的副作用，则是把本已被战争牵动的生计、迁徙和地方关系再度推入不确定之中。

同年四月，郑成功军登陆台湾，围攻荷兰东印度公司在热兰遮城的据点。对郑氏而言，这并非单纯向外扩张：南京行动受挫之后，需要一处能够安置军政、取得粮食并控制航路的基地；对岛上的荷兰据点、汉人移民社会与原住民社群而言，战争则把既有的土地、征发与交往关系重新置于强制之下。围城的胜负取决于守军补给能否进入、周边村社与航道由谁控制，以及攻守双方能否把人力和粮食留在城下，而不只是城墙前的一次决战。

次年正月，热兰遮城守军投降，荷兰东印度公司在台湾西南的统治至此终结。投降条约和《热兰遮城日志》能够确定交接的时点，也能让人看见围城中荷兰方面的处境；清廷军报与后出的史书则呈现北京获知这一消息的路径。它们不能替城外居民、被俘者或原住民社群作证。兵员、联盟和伤亡在中荷战报中各有立场，因而不能用任一方的数字填满这场转换。可以确认的是，郑氏接收了城堡、田地和贸易节点，台湾由此成为反清政权筹集人口与财政资源的基地；至于这种接收怎样抵达岛上不同社群、各处又怎样回应，仍须依赖保存不均的地方、契约和多语材料逐地重建。

郑成功在康熙元年五月去世，使新得的基地立刻面对继承问题。城堡已经易手，并不表示军政、宗族与海商网络已经形成稳定的服从关系；郑经最终掌权，是在这些网络重新安排权力与资源的过程中出现的结果。有关死因、继承争执和军中处置的细节，主要依赖后出的郑氏材料与荷兰报告，叙事只能把握重组的压力，不能替当事人补出一致的动机。

康熙二年（1663年），清军联合部分郑氏降将攻取厦门、金门，郑氏主力退守台湾。熟悉海上条件的降将与沿海据点帮助清军夺取近大陆的前进基地，但厦金易手并不等于清廷已经取得制海权。对郑氏而言，失去两岛迫使台湾的土地、港口和跨海贸易承担更重的军政负担；对清廷而言，迁界所制造的隔离、沿海守备与水师建设仍要持续付出社会与后勤成本。这里的制度得失并不在于哪一方终于“控制了海岸”，而在于双方都不得不以迁徙、城防、航路和粮饷来维持一场尚未结束的海上战争。


\begin{figure}[htbp]
  \centering
  \begin{tikzpicture}[
    x=1cm,y=1cm,
    place/.style={draw=timeline, fill=paper, rounded corners=2pt, align=center,
      inner sep=3pt, font=\small},
    qing/.style={place, draw=dynasty, fill=dynasty!13},
    ming/.style={place, draw=timeline, fill=timeline!13},
    sea/.style={place, draw=source, fill=source!10},
    note/.style={align=center, font=\scriptsize\color{source}}
  ]
    \fill[paper] (-.6,-3.8) rectangle (7.4,3.7);
    \draw[dynasty, line width=.8pt] (-.6,3.7) -- (7.4,3.7);
    \node[font=\bfseries\large, text=dynasty] at (3.4,3.25)
      {1645--1660：南方征服与南明、海上网络的压力};

    \node[qing] (nanjing) at (-.05,2.15) {南京\\1645};
    \node[qing] (nanchang) at (1.55,1.2) {南昌\\1648--49};
    \node[ming] (guangzhou) at (-.05,-.15) {广州\\1648--50};
    \node[ming] (xianggui) at (1.75,-.5) {湘桂\\李定国反攻\\1652};
    \node[ming] (yunnan) at (2.7,-2.2) {滇、黔\\永历朝退守\\1657--60};
    \node[ming] (myanmar) at (4.6,-2.85) {缅甸边地\\永历避难\\1659};
    \node[sea] (xiamen) at (5.1,1.65) {厦门、金门\\郑氏军商基地};
    \node[sea] (nanjingsea) at (3.9,2.45) {长江口、南京\\郑成功北伐\\1659};

    \draw[->, dynasty, line width=1.15pt] (nanjing) -- node[above,sloped,note]
      {江西战区（示意）} (nanchang);
    \draw[->, dynasty, line width=1.15pt] (nanchang) -- node[left,sloped,note]
      {清军重建攻势的节点} (guangzhou);
    \draw[->, dynasty, line width=1.15pt] (guangzhou) -- node[above,sloped,note]
      {湘桂、滇黔方向的战事} (xianggui);
    \draw[->, dynasty, line width=1.15pt] (xianggui) -- node[right,sloped,note]
      {入滇军行与补给（示意）} (yunnan);
    \draw[->, timeline, dashed, line width=1pt] (yunnan) -- node[below,sloped,note]
      {朝廷及随从继续西迁} (myanmar);
    \draw[->, timeline, dashed, line width=1pt] (xianggui) -- node[above,sloped,note]
      {南明反攻与联络\\并非固定战线} (guangzhou);
    \draw[->, source, line width=1pt] (xiamen) -- node[above,sloped,note]
      {海上补给与北上行动} (nanjingsea);
    \draw[source, dashed] (xiamen) to[out=-145,in=15] node[below,note]
      {闽南水陆通道（示意）} (guangzhou);

    \node[draw=source, fill=paper, rounded corners=3pt, align=center,
      text=source, font=\small] at (5.9,-.45)
      {实线：清军战区推进\\虚线：南明流动与反攻\\灰褐线：郑氏海上压力};
  \end{tikzpicture}
  \caption{南方征服、南明流动中心与郑氏海上压力（1645--1660，示意）。实线连接的只是账本所载南京、南昌、广州、湘桂与滇黔等战区节点；虚线和海线分别表示南明的退守／反攻及郑氏由闽南通向长江口的行动网络，不能读作疆界、补给线的全貌或海域控制范围。依据《清世祖实录》顺治二年至十七年军报、内阁大库江西军需和云南粮饷题本档，以及《永历实录》、弘光与永历朝奏疏、《南昌府志》《广东通志》《云南通志》；海上部分并读《延平王户官杨英从征实录》、郑氏文书辑录与荷兰东印度公司台湾商馆日记／海上报告，缅甸段参照缅文《琉璃宫史》传统。范围仅示华南--西南和闽南--长江口的少数城镇、山地与航路节点，未延伸为全国地图；路线约略且不按比例，尤其不以清廷军报或单一海图推定地方、岛屿或跨境地带的实际控制。}
  \label{fig:southern-conquest-southern-ming-sea}
\end{figure}

\subsection{1673—1684：撤藩、战事与跨海接收}

康熙十二年撤藩议定后，朝廷所要收回的并不只是几个藩王的名号。平西、平南、靖南三藩各有久驻军队、属下官员、土地与饷源；在云南、两广等地，它们又同商人、土司、地方绅耆和军户构成日常征收、转运与治安的关系。故而撤藩把继嗣与爵赏的问题，同兵丁去留、欠饷和地方中介的生计一并提出。吴三桂于1673年在云南反清，次年耿精忠在福建响应、王辅臣据平凉起事，郑经亦由台湾进兵福建。它们在时间上彼此牵动，却不宜合称为一套意志统一的“反清阵营”：藩镇自保、边地军人的选择、郑氏的海上经营及西北军政矛盾，各有其条件。

战场的推进尤其取决于省份能否持续供给。奏疏和军需册里频繁出现的饷银、米豆、夫役、船只与道路，并非战报的背景，而是战争得以跨省延续的条件。江西、湖广、福建、广东以及四川、陕西的转运，在不同阶段承受不同压力；一次解围或一座城的收复，也不能证明该省赋役、市场和交通已经复常。1675年前后尚之信在广州军政之间周旋，1676年耿精忠降清，清军又收复武昌等地。招抚文书能够说明朝廷愿意给予何种名分与条件，却不能据此推定降军已妥善编入绿营、旧部立即散去，或沿海贸易和港口税课已恢复。

吴三桂1677年在衡阳称帝，1678年去世后由吴世璠继承。清军随后在湖南、四川通向西南的线路上推进，既要攻城，也要维持粮运和接收既有军政网络。1681年昆明失守、吴世璠自尽，吴周政权遂告终结；这是叙述战局时可以把握的明确节点。可是驻军的安置、积欠的清理、荒地复垦和流民返乡不可能与城陷同步完成。地方志所写的“复业”“安集”常将多年过程压缩为一事，题本所报的数字又多服务于请饷、核销或论功；两者只能相互参照，不能直接换算为社会恢复的速度。

陆上战局收束后，北京才有较大的余力重新估算海峡的风险与成本。1681年郑经去世，东宁的继承冲突削弱了郑氏政权；1682年施琅受命统率福建水师，1683年澎湖战后郑克塽降清。这里的“澎湖战役”可作为跨海接收的军事转折，却不应把战报中的舰船、兵员、歼获和伤亡数逐项当作无误的统计：它们产生于请功、报捷和战时指挥的文书环境。1684年设台湾府及三县、隶属福建，并在迁界、复界与海禁上继续调整，标志着朝廷开始以府县、驻防和渡海管制处理台湾及其海域。行政建置并未使全岛立即成为同质的治理空间；汉人移民与租佃、原住民社群、熟番界务、渡台限制和既有贸易网络，仍使命令在各地呈现不同效果。

\paragraph{材料位置。} 《平定三逆方略》与《清圣祖实录》适于校定诏令、任免、战事和处分的朝廷序列，内阁大库题本、军需与仓储档案可追问某次行动所需的人力、钱粮和转运。它们首先记录的是官员呈报与中央裁决，不能单独替代地方经验。吴周檄文、郑氏文书、地方志、族谱和契约则可补看动员、迁徙、土地与贸易的后果；但前二者各有立场，后几类往往成书较晚、保存也高度不均。因而可较有把握地叙述政权更替与设治日期，对某地人口损失、物价、逃亡规模或战后“恢复”的普遍程度，则应保留材料所允许的限度。

\section{撤藩：把临时安排变成战争}

三藩并非清初版图上一块静止的颜色。平西、平南、靖南三藩拥有驻军、地盘和饷源，是征服西南和两广的工具，也使中央财政与皇权面临一个不能无限期延宕的问题。康熙十二年三月，朝廷批准或安排撤藩（EQ-1673-FEUDATORY-ABOLITION）。《清圣祖实录》、内阁大库题本与《平定三逆方略》能确立决策的文书序列；各藩请求的先后、朝议的立场以及军费核算，仍须把奏报、部议与最终处分拆开阅读。撤藩不是一句抽象的“削藩”，而是把将领家族的继承、军队的粮饷和地方中介的前途同时置于风险之中。

吴三桂于同年十一月在云南起兵（EQ-1673-WU-REVOLT），耿精忠次年在福建反叛，王辅臣据平凉起事（EQ-1674-GENG-JINGZHONG、EQ-1674-WANG-FUCHEN）。战争迅速把云南、福建和西北交通连在一起。吴军檄文中的忠明、族群和财政语言是动员资源，不可直接等同所有参与者的动机；《平定三逆方略》是清廷战时汇编，也会按凯旋叙事组织敌我。福建的耿精忠既需郑氏海上力量，又要维持陆上府县，郑经同年自台湾入闽、占据若干沿海节点（EQ-1674-ZHENG-JING-FUJIAN）。据点易手可由军报、方志和郑氏材料相互钉住，盟约的范围、征粮和实际控制却只能逐地核查。

两广的尚之信在1675年前后控制广州军政、反复周旋（EQ-1675-SHANG-ZHIXIN）。这类“反清”或“归清”的标签若不按命令、行动和时间分段，就会遮蔽港口财政、家族继承与地方军队的讨价还价。1676年耿精忠投降，清军收复武昌等长江中游节点（EQ-1676-GENG-SURRENDER、EQ-1676-WUCHANG）。赦降文书说明招抚的条件，不能自动说明降军已被安置、郑氏影响已消退；一座府城的收复，也不等于周边县乡已脱离战争。

吴三桂1677年在衡阳称帝、国号周，试图以更高的政治名义整合将领、赋役和盟友（EQ-1677-WU-ZHOU）。称帝本身可靠，礼制细节和各地响应范围却主要见于敌对军报，需同吴周檄文互证。翌年吴三桂去世，吴世璠继承（EQ-1678-WU-DEATH）；此后清军在湖南推进、清理四川通往西南的道路（EQ-1679-HUNAN-RECOVERY、EQ-1680-SICHUAN-RECOVERY）。《实录》与方略记录了军事推进，地方志、军需题本和契约材料才可能局部说明逃亡、仓储、垦复和赋役恢复。特别是四川战后人口与耕地数字，口径差异很大，不应从任何一种招民叙事中推出整齐的复苏。

康熙二十年十月，清军入昆明，吴世璠自尽，吴周政权结束（EQ-1681-KUNMING）。城陷与死亡是可靠的终局节点；围城损失、降附与后续清剿并非同日完成。多年战争中形成的驻军、饷运、道路损耗与地方恢复，仍将塑造云南、贵州和两广的治理。中央收回的是名义军政权，却继承了把这一权力落到山地、城镇和税源上的长期成本。


\section{跨海接收之后}

郑经1681年去世，东宁的继承冲突使郑克塽最终掌权（EQ-1681-ZHENG-JING-DEATH）。死期与继承结果可由清廷情报、郑氏后出材料和荷兰通信互校；宫廷政变的每一步不能因后出的叙述而获得同等确定性。三藩战争结束与东宁的不稳，给北京重新估算跨海行动提供了窗口。康熙二十一年任命施琅统率福建水师（EQ-1682-SHI-LANG），任命与筹备见于实录和军务奏折；施琅的战略意图及朝中争论须区别奏折、朱批与后出的传记，不能把成功后的理由当成当时唯一的选择。

康熙二十二年六月，施琅水师在澎湖击败郑氏舰队（EQ-1683-PENGHU）。胜败与其战略后果可以确立；舰船、兵员、死伤数字多出自胜方战报，须同郑氏叙事及海域材料交叉约束。季风、福建补给和航道知识，解释了行动为何能在那个时点集中展开，也提醒我们水师不是从地图上直线抵达台湾。八月郑克塽投降（EQ-1683-TAIWAN-SURRENDER），清廷取得政治接收权，但投降条件、官兵安置和移民、原住民社群的变化属于不同层次的过程，不能压缩为“台湾归入清朝”一项结果。

1684年设置台湾府及诸县、隶属福建，同时调整迁界与海禁，准许居民复界并逐步恢复海贸（EQ-1684-TAIWAN-PREFECTURE、EQ-1684-SEA-BAN-LIFT）。实录、地理志和福建档案能够证明制度设置和政策转向；府县与驻防的建立不等于全岛有效治理，复界诏令也不等于每一港口在同一时间开放。平原、沿海、内山与原住民社群同官府的接触程度并不相同。1721年朱一贵起事、一度进入府城，清廷继而派兵恢复控制（EQ-1721-ZHU-YIGUI）。军报可确立起事、府城易手和镇压的序列；官方“逆匪”分类不能代替参与者的社会构成或赋役诉求。地方志、督抚奏折和契约文书提示垦殖、租佃、治安及征收摩擦，却也只提供局部证词。跨海设治的脆弱性，正由这场起事暴露出来。


\subsection{事件链与材料限度}
东宁、台湾与清廷之间的转换须放在继承、航运与地方安置的不同节奏中阅读。

\paragraph{康熙二十年（1681年）：郑经去世，东宁发生继承冲突，郑克塽最终掌权}
\eventanchor{EQ-1681-ZHENG-JING-DEATH}{郑经去世，东宁发生继承冲突，郑克塽最终掌权}。本节点叙述该事件本身。东宁郑氏政权在京师与多地军政线路相接的尺度的处境首先受郑经长期在大陆作战并承受台湾财政压力，其死亡使宗族、将领和文官围绕继承重组制约。《清圣祖实录》康熙二十年台湾情报；《清史稿·郑成功传附郑经》；《台湾外纪》与荷兰商馆残存通信留下可核的线索；可辨的后果是东宁短期政治不稳，清廷得以评估以水师进攻台湾的风险和机会。原有置信注记：郑经死期和继承结果可靠；宫廷政变过程主要依郑氏后出材料与清方情报

\paragraph{康熙二十年（1681年）：郑经去世，东宁发生继承冲突，郑克塽最终掌权}
\eventanchor{EQ-1681-ZHENG-JING-DEATH-AFTERMATH}{郑经去世，东宁发生继承冲突，郑克塽最终掌权} 置于京师与多地军政线路相接的尺度中阅读，本节点追踪「郑经去世，东宁发生继承冲突，郑克塽最终掌权」之后可见的余波。郑经长期在大陆作战并承受台湾财政压力，其死亡使宗族、将领和文官围绕继承重组构成行动者东宁郑氏政权必须面对的条件。取证从《清圣祖实录》康熙二十年台湾情报；《清史稿·郑成功传附郑经》；《台湾外纪》与荷兰商馆残存通信出发，因而只能把变化谨慎地连到东宁短期政治不稳，清廷得以评估以水师进攻台湾的风险和机会。原有置信注记：郑经死期和继承结果可靠；宫廷政变过程主要依郑氏后出材料与清方情报

\paragraph{康熙二十年（1681年）：郑经去世，东宁发生继承冲突，郑克塽最终掌权}
\eventanchor{EQ-1681-ZHENG-JING-DEATH-DECISION}{郑经去世，东宁发生继承冲突，郑克塽最终掌权} 标记本节点定位围绕「郑经去世，东宁发生继承冲突，郑克塽最终掌权」形成的处置与调度记录，涉及东宁郑氏政权。京师与多地军政线路相接的尺度不是抽象背景：郑经长期在大陆作战并承受台湾财政压力，其死亡使宗族、将领和文官围绕继承重组。《清圣祖实录》康熙二十年台湾情报；《清史稿·郑成功传附郑经》；《台湾外纪》与荷兰商馆残存通信所见的顺序随后指向东宁短期政治不稳，清廷得以评估以水师进攻台湾的风险和机会，但原有置信注记：郑经死期和继承结果可靠；宫廷政变过程主要依郑氏后出材料与清方情报

\paragraph{康熙二十年（1681年）：郑经去世，东宁发生继承冲突，郑克塽最终掌权}
在京师与多地军政线路相接的尺度，\eventanchor{EQ-1681-ZHENG-JING-DEATH-PRELUDE}{郑经去世，东宁发生继承冲突，郑克塽最终掌权} 让东宁郑氏政权的一个选择或压力有了可查的坐标。本节点回看「郑经去世，东宁发生继承冲突，郑克塽最终掌权」之前已可见的条件。前因可写为郑经长期在大陆作战并承受台湾财政压力，其死亡使宗族、将领和文官围绕继承重组；《清圣祖实录》康熙二十年台湾情报；《清史稿·郑成功传附郑经》；《台湾外纪》与荷兰商馆残存通信限定了记录，后续变化是东宁短期政治不稳，清廷得以评估以水师进攻台湾的风险和机会。原有置信注记：郑经死期和继承结果可靠；宫廷政变过程主要依郑氏后出材料与清方情报

\paragraph{康熙二十一年（1682年）：清廷任命施琅统率福建水师，筹备进攻台湾}
\eventanchor{EQ-1682-SHI-LANG}{清廷任命施琅统率福建水师，筹备进攻台湾} 所关联的是清与东宁郑氏政权在港口、海峡与跨海航路相连的尺度中的行动。本节点叙述该事件本身。它从三藩平定和东宁继承危机减少内陆牵制，厦金旧将对航道与郑氏军制的了解提供条件展开，借《清圣祖实录》康熙二十一年福建海防条；《清史稿·施琅传》；《康熙朝满文朱批奏折全译》康熙二十一年台湾军务奏折进入可检验的年序，并以清廷从迁界防御转向跨海进攻，水师训练、粮械和季风仍决定行动可行性影响下一段安排。原有置信注记：任命和筹备可靠；施琅战略构想与朝廷内部争论须以奏折、朱批和后出传记区别处理

\paragraph{康熙二十一年（1682年）：清廷任命施琅统率福建水师，筹备进攻台湾}
港口、海峡与跨海航路相连的尺度中的\eventanchor{EQ-1682-SHI-LANG-AFTERMATH}{清廷任命施琅统率福建水师，筹备进攻台湾} 不能离开清与东宁郑氏政权来解释。本节点追踪「清廷任命施琅统率福建水师，筹备进攻台湾」之后可见的余波。三藩平定和东宁继承危机减少内陆牵制，厦金旧将对航道与郑氏军制的了解提供条件说明其形成的条件；《清圣祖实录》康熙二十一年福建海防条；《清史稿·施琅传》；《康熙朝满文朱批奏折全译》康熙二十一年台湾军务奏折提供来源路径，清廷从迁界防御转向跨海进攻，水师训练、粮械和季风仍决定行动可行性则是材料能够支持的近时结果。原有置信注记：任命和筹备可靠；施琅战略构想与朝廷内部争论须以奏折、朱批和后出传记区别处理

\paragraph{康熙二十一年（1682年）：清廷任命施琅统率福建水师，筹备进攻台湾}
\eventanchor{EQ-1682-SHI-LANG-DECISION}{清廷任命施琅统率福建水师，筹备进攻台湾} 把清与东宁郑氏政权放回港口、海峡与跨海航路相连的尺度的道路、城镇、边界或制度网络。本节点定位围绕「清廷任命施琅统率福建水师，筹备进攻台湾」形成的处置与调度记录。三藩平定和东宁继承危机减少内陆牵制，厦金旧将对航道与郑氏军制的了解提供条件。本段采用《清圣祖实录》康熙二十一年福建海防条；《清史稿·施琅传》；《康熙朝满文朱批奏折全译》康熙二十一年台湾军务奏折核对其顺序，随后可追到清廷从迁界防御转向跨海进攻，水师训练、粮械和季风仍决定行动可行性。原有置信注记：任命和筹备可靠；施琅战略构想与朝廷内部争论须以奏折、朱批和后出传记区别处理

\paragraph{康熙二十一年（1682年）：清廷任命施琅统率福建水师，筹备进攻台湾}
本节点回看「清廷任命施琅统率福建水师，筹备进攻台湾」之前已可见的条件，故\eventanchor{EQ-1682-SHI-LANG-PRELUDE}{清廷任命施琅统率福建水师，筹备进攻台湾} 不只是一个年份标签。清与东宁郑氏政权在港口、海峡与跨海航路相连的尺度承受的压力来自三藩平定和东宁继承危机减少内陆牵制，厦金旧将对航道与郑氏军制的了解提供条件。《清圣祖实录》康熙二十一年福建海防条；《清史稿·施琅传》；《康熙朝满文朱批奏折全译》康熙二十一年台湾军务奏折可回查该节点；清廷从迁界防御转向跨海进攻，水师训练、粮械和季风仍决定行动可行性是其进入后续年序的方式。原有置信注记：任命和筹备可靠；施琅战略构想与朝廷内部争论须以奏折、朱批和后出传记区别处理

\paragraph{康熙二十二年六月（1683年）：施琅水师在澎湖击败郑氏舰队，取得通往台湾本岛的制海优势}
\eventanchor{EQ-1683-PENGHU}{施琅水师在澎湖击败郑氏舰队，取得通往台湾本岛的制海优势} 所见的行动者为清与东宁郑氏政权，空间尺度为港口、海峡与跨海航路相连的尺度。本节点叙述该事件本身。清军利用季风窗口、福建补给和施琅对郑氏舰队的熟悉发动集中攻击，郑氏难以同时守澎湖与本岛使它成为具体事件链的一环；《清圣祖实录》康熙二十二年六月；《清史稿·施琅传》；施琅《飞报澎湖大捷疏》与台湾地方志与郑氏失去海上屏障，台湾方面转入投降谈判；澎湖伤亡和战术细节仍需多源检验。原有置信注记：澎湖战役结果可靠；舰船、兵员和死伤数字多来自胜方战报，须以郑氏和海域材料约束

\paragraph{康熙二十二年六月（1683年）：施琅水师在澎湖击败郑氏舰队，取得通往台湾本岛的制海优势}
从港口、海峡与跨海航路相连的尺度看，\eventanchor{EQ-1683-PENGHU-AFTERMATH}{施琅水师在澎湖击败郑氏舰队，取得通往台湾本岛的制海优势} 留下本节点追踪「施琅水师在澎湖击败郑氏舰队，取得通往台湾本岛的制海优势」之后可见的余波的材料坐标。清与东宁郑氏政权所面对的条件是清军利用季风窗口、福建补给和施琅对郑氏舰队的熟悉发动集中攻击，郑氏难以同时守澎湖与本岛。《清圣祖实录》康熙二十二年六月；《清史稿·施琅传》；施琅《飞报澎湖大捷疏》与台湾地方志支持的结果为郑氏失去海上屏障，台湾方面转入投降谈判；澎湖伤亡和战术细节仍需多源检验。原有置信注记：澎湖战役结果可靠；舰船、兵员和死伤数字多来自胜方战报，须以郑氏和海域材料约束

\paragraph{康熙二十二年六月（1683年）：施琅水师在澎湖击败郑氏舰队，取得通往台湾本岛的制海优势}
\eventanchor{EQ-1683-PENGHU-DECISION}{施琅水师在澎湖击败郑氏舰队，取得通往台湾本岛的制海优势} 记录清与东宁郑氏政权在港口、海峡与跨海航路相连的尺度的一次变化。本节点定位围绕「施琅水师在澎湖击败郑氏舰队，取得通往台湾本岛的制海优势」形成的处置与调度记录。其发生与清军利用季风窗口、福建补给和施琅对郑氏舰队的熟悉发动集中攻击，郑氏难以同时守澎湖与本岛相连；《清圣祖实录》康熙二十二年六月；《清史稿·施琅传》；施琅《飞报澎湖大捷疏》与台湾地方志把事件系回来源，郑氏失去海上屏障，台湾方面转入投降谈判；澎湖伤亡和战术细节仍需多源检验显示压力如何向后转移。原有置信注记：澎湖战役结果可靠；舰船、兵员和死伤数字多来自胜方战报，须以郑氏和海域材料约束

\paragraph{康熙二十二年六月（1683年）：施琅水师在澎湖击败郑氏舰队，取得通往台湾本岛的制海优势}
\eventanchor{EQ-1683-PENGHU-PRELUDE}{施琅水师在澎湖击败郑氏舰队，取得通往台湾本岛的制海优势} 的地点与行动范围属于港口、海峡与跨海航路相连的尺度，行动者是清与东宁郑氏政权。本节点回看「施琅水师在澎湖击败郑氏舰队，取得通往台湾本岛的制海优势」之前已可见的条件。清军利用季风窗口、福建补给和施琅对郑氏舰队的熟悉发动集中攻击，郑氏难以同时守澎湖与本岛。读《清圣祖实录》康熙二十二年六月；《清史稿·施琅传》；施琅《飞报澎湖大捷疏》与台湾地方志时可见其记录边界，最小后果只能写作郑氏失去海上屏障，台湾方面转入投降谈判；澎湖伤亡和战术细节仍需多源检验。原有置信注记：澎湖战役结果可靠；舰船、兵员和死伤数字多来自胜方战报，须以郑氏和海域材料约束

\paragraph{康熙二十二年八月（1683年）：郑克塽向清廷投降，东宁政权结束}
\eventanchor{EQ-1683-TAIWAN-SURRENDER}{郑克塽向清廷投降，东宁政权结束} 是清与东宁郑氏政权、台湾社会在港口、海峡与跨海航路相连的尺度留下的编年标记。本节点叙述该事件本身。它从澎湖失败切断东宁海防与外援预期，继承危机使郑氏难以继续集中抵抗走来；《清圣祖实录》康熙二十二年八月；《清史稿·郑成功传附郑克塽》；施琅奏折与《台湾府志》固定可证部分，清廷取得台湾政治接收权，后续设治、迁界调整和土地社会重组并非投降即告完成构成下一步的可见结果。原有置信注记：投降与清廷接收主干可靠；投降条件、官兵安置和原住民、移民社会变化须分层记录

\paragraph{康熙二十二年八月（1683年）：郑克塽向清廷投降，东宁政权结束}
本节点追踪「郑克塽向清廷投降，东宁政权结束」之后可见的余波：\eventanchor{EQ-1683-TAIWAN-SURRENDER-AFTERMATH}{郑克塽向清廷投降，东宁政权结束}。清与东宁郑氏政权、台湾社会在港口、海峡与跨海航路相连的尺度的处境可由澎湖失败切断东宁海防与外援预期，继承危机使郑氏难以继续集中抵抗说明。《清圣祖实录》康熙二十二年八月；《清史稿·郑成功传附郑克塽》；施琅奏折与《台湾府志》是本段材料入口，因而只能据此追踪到清廷取得台湾政治接收权，后续设治、迁界调整和土地社会重组并非投降即告完成。原有置信注记：投降与清廷接收主干可靠；投降条件、官兵安置和原住民、移民社会变化须分层记录

\paragraph{康熙二十二年八月（1683年）：郑克塽向清廷投降，东宁政权结束}
\eventanchor{EQ-1683-TAIWAN-SURRENDER-DECISION}{郑克塽向清廷投降，东宁政权结束} 发生在港口、海峡与跨海航路相连的尺度，牵动清与东宁郑氏政权、台湾社会。本节点定位围绕「郑克塽向清廷投降，东宁政权结束」形成的处置与调度记录。澎湖失败切断东宁海防与外援预期，继承危机使郑氏难以继续集中抵抗是其结构性条件；《清圣祖实录》康熙二十二年八月；《清史稿·郑成功传附郑克塽》；施琅奏折与《台湾府志》留下的证据把读者带到清廷取得台湾政治接收权，后续设治、迁界调整和土地社会重组并非投降即告完成。原有置信注记：投降与清廷接收主干可靠；投降条件、官兵安置和原住民、移民社会变化须分层记录

\paragraph{康熙二十二年八月（1683年）：郑克塽向清廷投降，东宁政权结束}
\eventanchor{EQ-1683-TAIWAN-SURRENDER-PRELUDE}{郑克塽向清廷投降，东宁政权结束} 对应本节点回看「郑克塽向清廷投降，东宁政权结束」之前已可见的条件。在港口、海峡与跨海航路相连的尺度，清与东宁郑氏政权、台湾社会无法绕开澎湖失败切断东宁海防与外援预期，继承危机使郑氏难以继续集中抵抗。依据《清圣祖实录》康熙二十二年八月；《清史稿·郑成功传附郑克塽》；施琅奏折与《台湾府志》，本段把后续影响限定为清廷取得台湾政治接收权，后续设治、迁界调整和土地社会重组并非投降即告完成。原有置信注记：投降与清廷接收主干可靠；投降条件、官兵安置和原住民、移民社会变化须分层记录

\paragraph{康熙二十三年（1684年）：清廷调整迁界与海禁政策，允许沿海居民复界并逐步恢复民间海贸}
\eventanchor{EQ-1684-SEA-BAN-LIFT}{清廷调整迁界与海禁政策，允许沿海居民复界并逐步恢复民间海贸}。本节点叙述该事件本身。清与闽粤沿海在港口、海峡与跨海航路相连的尺度的处境首先受台湾投降降低清廷对郑氏补给网络的担忧，长期迁界的税源、民生和治安代价亦需处理制约。《清圣祖实录》康熙二十三年复界海禁条；《清史稿·食货志》；《福建通志》《广东通志》沿海复界条留下可核的线索；可辨的后果是沿海耕作、渔盐和航运开始恢复，但贸易许可、海防和地方执行塑造不同港口开放节奏。原有置信注记：政策转向可靠；复界日期、港口开放和走私变化因省县而异，不能由中央谕旨推全国

\paragraph{康熙二十三年（1684年）：清廷调整迁界与海禁政策，允许沿海居民复界并逐步恢复民间海贸}
\eventanchor{EQ-1684-SEA-BAN-LIFT-AFTERMATH}{清廷调整迁界与海禁政策，允许沿海居民复界并逐步恢复民间海贸} 置于港口、海峡与跨海航路相连的尺度中阅读，本节点追踪「清廷调整迁界与海禁政策，允许沿海居民复界并逐步恢复民间海贸」之后可见的余波。台湾投降降低清廷对郑氏补给网络的担忧，长期迁界的税源、民生和治安代价亦需处理构成行动者清与闽粤沿海必须面对的条件。取证从《清圣祖实录》康熙二十三年复界海禁条；《清史稿·食货志》；《福建通志》《广东通志》沿海复界条出发，因而只能把变化谨慎地连到沿海耕作、渔盐和航运开始恢复，但贸易许可、海防和地方执行塑造不同港口开放节奏。原有置信注记：政策转向可靠；复界日期、港口开放和走私变化因省县而异，不能由中央谕旨推全国

\paragraph{康熙二十三年（1684年）：清廷调整迁界与海禁政策，允许沿海居民复界并逐步恢复民间海贸}
\eventanchor{EQ-1684-SEA-BAN-LIFT-DECISION}{清廷调整迁界与海禁政策，允许沿海居民复界并逐步恢复民间海贸} 标记本节点定位围绕「清廷调整迁界与海禁政策，允许沿海居民复界并逐步恢复民间海贸」形成的处置与调度记录，涉及清与闽粤沿海。港口、海峡与跨海航路相连的尺度不是抽象背景：台湾投降降低清廷对郑氏补给网络的担忧，长期迁界的税源、民生和治安代价亦需处理。《清圣祖实录》康熙二十三年复界海禁条；《清史稿·食货志》；《福建通志》《广东通志》沿海复界条所见的顺序随后指向沿海耕作、渔盐和航运开始恢复，但贸易许可、海防和地方执行塑造不同港口开放节奏，但原有置信注记：政策转向可靠；复界日期、港口开放和走私变化因省县而异，不能由中央谕旨推全国

\paragraph{康熙二十三年（1684年）：清廷调整迁界与海禁政策，允许沿海居民复界并逐步恢复民间海贸}
在港口、海峡与跨海航路相连的尺度，\eventanchor{EQ-1684-SEA-BAN-LIFT-PRELUDE}{清廷调整迁界与海禁政策，允许沿海居民复界并逐步恢复民间海贸} 让清与闽粤沿海的一个选择或压力有了可查的坐标。本节点回看「清廷调整迁界与海禁政策，允许沿海居民复界并逐步恢复民间海贸」之前已可见的条件。前因可写为台湾投降降低清廷对郑氏补给网络的担忧，长期迁界的税源、民生和治安代价亦需处理；《清圣祖实录》康熙二十三年复界海禁条；《清史稿·食货志》；《福建通志》《广东通志》沿海复界条限定了记录，后续变化是沿海耕作、渔盐和航运开始恢复，但贸易许可、海防和地方执行塑造不同港口开放节奏。原有置信注记：政策转向可靠；复界日期、港口开放和走私变化因省县而异，不能由中央谕旨推全国

\paragraph{康熙二十三年（1684年）：清廷设置台湾府及诸县，隶属福建省，开始以府县和驻防制度经营台湾}
\eventanchor{EQ-1684-TAIWAN-PREFECTURE}{清廷设置台湾府及诸县，隶属福建省，开始以府县和驻防制度经营台湾} 所关联的是清与台湾在港口、海峡与跨海航路相连的尺度中的行动。本节点叙述该事件本身。它从东宁投降后清廷须在弃守与设治间权衡海防、税源、移民和驻军成本展开，借《清圣祖实录》康熙二十三年台湾设治条；《清史稿·地理志》；《台湾府志》与福建省档案相关奏议进入可检验的年序，并以台湾纳入福建行政体系，但沿海、平原、内山和原住民社群与官府关系呈现不均衡接触影响下一段安排。原有置信注记：设府设县可靠；机构设置不等于全岛有效治理，原住民地区和内山控制尤须保留

\paragraph{康熙二十三年（1684年）：清廷设置台湾府及诸县，隶属福建省，开始以府县和驻防制度经营台湾}
港口、海峡与跨海航路相连的尺度中的\eventanchor{EQ-1684-TAIWAN-PREFECTURE-AFTERMATH}{清廷设置台湾府及诸县，隶属福建省，开始以府县和驻防制度经营台湾} 不能离开清与台湾来解释。本节点追踪「清廷设置台湾府及诸县，隶属福建省，开始以府县和驻防制度经营台湾」之后可见的余波。东宁投降后清廷须在弃守与设治间权衡海防、税源、移民和驻军成本说明其形成的条件；《清圣祖实录》康熙二十三年台湾设治条；《清史稿·地理志》；《台湾府志》与福建省档案相关奏议提供来源路径，台湾纳入福建行政体系，但沿海、平原、内山和原住民社群与官府关系呈现不均衡接触则是材料能够支持的近时结果。原有置信注记：设府设县可靠；机构设置不等于全岛有效治理，原住民地区和内山控制尤须保留

\paragraph{康熙二十三年（1684年）：清廷设置台湾府及诸县，隶属福建省，开始以府县和驻防制度经营台湾}
\eventanchor{EQ-1684-TAIWAN-PREFECTURE-DECISION}{清廷设置台湾府及诸县，隶属福建省，开始以府县和驻防制度经营台湾} 把清与台湾放回港口、海峡与跨海航路相连的尺度的道路、城镇、边界或制度网络。本节点定位围绕「清廷设置台湾府及诸县，隶属福建省，开始以府县和驻防制度经营台湾」形成的处置与调度记录。东宁投降后清廷须在弃守与设治间权衡海防、税源、移民和驻军成本。本段采用《清圣祖实录》康熙二十三年台湾设治条；《清史稿·地理志》；《台湾府志》与福建省档案相关奏议核对其顺序，随后可追到台湾纳入福建行政体系，但沿海、平原、内山和原住民社群与官府关系呈现不均衡接触。原有置信注记：设府设县可靠；机构设置不等于全岛有效治理，原住民地区和内山控制尤须保留

\paragraph{康熙二十三年（1684年）：清廷设置台湾府及诸县，隶属福建省，开始以府县和驻防制度经营台湾}
本节点回看「清廷设置台湾府及诸县，隶属福建省，开始以府县和驻防制度经营台湾」之前已可见的条件，故\eventanchor{EQ-1684-TAIWAN-PREFECTURE-PRELUDE}{清廷设置台湾府及诸县，隶属福建省，开始以府县和驻防制度经营台湾} 不只是一个年份标签。清与台湾在港口、海峡与跨海航路相连的尺度承受的压力来自东宁投降后清廷须在弃守与设治间权衡海防、税源、移民和驻军成本。《清圣祖实录》康熙二十三年台湾设治条；《清史稿·地理志》；《台湾府志》与福建省档案相关奏议可回查该节点；台湾纳入福建行政体系，但沿海、平原、内山和原住民社群与官府关系呈现不均衡接触是其进入后续年序的方式。原有置信注记：设府设县可靠；机构设置不等于全岛有效治理，原住民地区和内山控制尤须保留

\paragraph{康熙六十年（1721年）：朱一贵在台湾发动反清起事，一度攻入府城，清廷随后派兵恢复控制}
\eventanchor{EQ-1721-ZHU-YIGUI}{朱一贵在台湾发动反清起事，一度攻入府城，清廷随后派兵恢复控制} 所见的行动者为清与台湾朱一贵集团，空间尺度为港口、海峡与跨海航路相连的尺度。本节点叙述该事件本身。台湾设治后垦殖、租佃、族群关系与地方行政摩擦累积，官府征收和治安处置成为触发因素使它成为具体事件链的一环；《清圣祖实录》康熙六十年台湾军报；《清史稿·朱一贵传》；《台湾府志》、福建督抚奏折与契约文书与清廷暴露跨海调兵与府县治理薄弱环节，后续强化驻防和行政整顿；不能将起事简化为单一族群反应。原有置信注记：起事、府城易手与镇压主干可靠；参与群体、赋税原因和伤亡数字受官府“逆匪”分类影响

\paragraph{康熙六十年（1721年）：朱一贵在台湾发动反清起事，一度攻入府城，清廷随后派兵恢复控制}
从港口、海峡与跨海航路相连的尺度看，\eventanchor{EQ-1721-ZHU-YIGUI-AFTERMATH}{朱一贵在台湾发动反清起事，一度攻入府城，清廷随后派兵恢复控制} 留下本节点追踪「朱一贵在台湾发动反清起事，一度攻入府城，清廷随后派兵恢复控制」之后可见的余波的材料坐标。清与台湾朱一贵集团所面对的条件是台湾设治后垦殖、租佃、族群关系与地方行政摩擦累积，官府征收和治安处置成为触发因素。《清圣祖实录》康熙六十年台湾军报；《清史稿·朱一贵传》；《台湾府志》、福建督抚奏折与契约文书支持的结果为清廷暴露跨海调兵与府县治理薄弱环节，后续强化驻防和行政整顿；不能将起事简化为单一族群反应。原有置信注记：起事、府城易手与镇压主干可靠；参与群体、赋税原因和伤亡数字受官府“逆匪”分类影响

\paragraph{康熙六十年（1721年）：朱一贵在台湾发动反清起事，一度攻入府城，清廷随后派兵恢复控制}
\eventanchor{EQ-1721-ZHU-YIGUI-DECISION}{朱一贵在台湾发动反清起事，一度攻入府城，清廷随后派兵恢复控制} 记录清与台湾朱一贵集团在港口、海峡与跨海航路相连的尺度的一次变化。本节点定位围绕「朱一贵在台湾发动反清起事，一度攻入府城，清廷随后派兵恢复控制」形成的处置与调度记录。其发生与台湾设治后垦殖、租佃、族群关系与地方行政摩擦累积，官府征收和治安处置成为触发因素相连；《清圣祖实录》康熙六十年台湾军报；《清史稿·朱一贵传》；《台湾府志》、福建督抚奏折与契约文书把事件系回来源，清廷暴露跨海调兵与府县治理薄弱环节，后续强化驻防和行政整顿；不能将起事简化为单一族群反应显示压力如何向后转移。原有置信注记：起事、府城易手与镇压主干可靠；参与群体、赋税原因和伤亡数字受官府“逆匪”分类影响

\paragraph{康熙六十年（1721年）：朱一贵在台湾发动反清起事，一度攻入府城，清廷随后派兵恢复控制}
\eventanchor{EQ-1721-ZHU-YIGUI-PRELUDE}{朱一贵在台湾发动反清起事，一度攻入府城，清廷随后派兵恢复控制} 的地点与行动范围属于港口、海峡与跨海航路相连的尺度，行动者是清与台湾朱一贵集团。本节点回看「朱一贵在台湾发动反清起事，一度攻入府城，清廷随后派兵恢复控制」之前已可见的条件。台湾设治后垦殖、租佃、族群关系与地方行政摩擦累积，官府征收和治安处置成为触发因素。读《清圣祖实录》康熙六十年台湾军报；《清史稿·朱一贵传》；《台湾府志》、福建督抚奏折与契约文书时可见其记录边界，最小后果只能写作清廷暴露跨海调兵与府县治理薄弱环节，后续强化驻防和行政整顿；不能将起事简化为单一族群反应。原有置信注记：起事、府城易手与镇压主干可靠；参与群体、赋税原因和伤亡数字受官府“逆匪”分类影响

\subsection{1685—1698：黑龙江交涉与草原战争}

1685年清军围攻阿尔巴津，俄方撤离后重返，1686年清军再次围困。这里并不是由后世国界预先决定归属的一点：围城牵连黑龙江防务、河道、补给、疾病和俄方西伯利亚的驻军体系。1689年双方在尼布楚签约，将军事僵持暂时转化为多语文本安排。满文、拉丁文和俄文条约文本能够确认签约及其措辞；译名、地物对应和界线含义并未因此完全一致，后续还要经过使节、地图和地方执行才获得具体内容。

在条约谈判前后，西北草原已打开另一条成本更高的线。1688年噶尔丹进攻喀尔喀，诸部南徙并向清廷求援；1690年乌兰布通交战阻止其继续深入内蒙古，但没有立即消灭其军队。1691年多伦诺尔会盟把军事保护、封爵、补给和盟旗安排结合起来。会盟和册封不是单向完成的“归附”：牧地使用、领主自主性、安置和跨部纠纷仍须在随后材料中追踪。

1695年远征集中军粮、驼马和多路部队，运输、牲畜损耗、季节和疾病限制了行动。昭莫多战后噶尔丹西遁，1697年去世；清廷继而在1698年调整喀尔喀的盟旗、台站和赈济安排。凯旋叙事能够说明朝廷怎样排列胜负，军需奏折、蒙古文材料和俄方记录更能显出远征的距离、供给与迁徙。清军取得军事优势，并不等于准噶尔政权、跨境贸易或牧地争议已经消失。

\paragraph{材料位置。} 黑龙江军报、俄国西伯利亚衙门文件和城址考古能相互约束围城的物质空间，却不能单独裁定伤亡或谈判责任。草原方面，《清圣祖实录》、理藩院满蒙文文书、《蒙古源流》及俄国使团记录各自保存不同的政治与交通视角；“归附”一词必须放回这些文书的产生关系中阅读。

\section{黑龙江的围城与多语边界}

阿尔巴津不是一个由后世国界预先决定归属的点。1685年清军围攻该地、俄方撤离；俄方重返后，1686年清军再次围困。《清圣祖实录》的黑龙江军报显示清廷如何理解防务与行动，俄国西伯利亚衙门文件和驻军报告则记录另一套补给、疾病与领土主张。城址考古可以约束围困发生的物质空间，却不能单独裁定伤亡或谈判责任。守军人数、疾病死亡和停战时点必须保留满文、汉文与俄文材料之间的差异。

1689年七月双方在\eventanchor{EQ-1689-NERCHINSK}{尼布楚签约}。满文、拉丁文、俄文文本的存在和签约本身可靠；耶稣会译员参与的多语谈判并没有消除地名、术语和界线含义的不一致。《清史稿·邦交志》可作后出的检索线索，不能代替条约各本与俄国谈判档。条约把阿尔巴津危机暂时制度化，随后仍须由使节往来、地图制作和地方执行来赋予它具体含义。因此，尼布楚不是一条可直接投射成精确现代线的结论，而是双方在军事僵持、地理知识不对称和翻译条件中达成的文本安排。


\section{草原同盟、迁徙与远征的成本}

1688年噶尔丹进攻喀尔喀，诸部南徙并向清廷求援（EQ-1688-GALDAN-KHALKHA）。入侵与迁徙是可靠主干；部众规模、从属关系和“归附”的措辞，不能只由清廷满汉军报决定。蒙古文材料、《蒙古源流》与俄国使团记录保留不同的政治与交通视角。清廷获得介入的机会，也面对安置、牧地、赈济和对准噶尔战争的成本。1690年乌兰布通交战阻止噶尔丹继续深入内蒙古，却未立即消灭其军队（EQ-1690-ULAN-BUTUNG）。方略的胜负语言、兵力和伤亡，需要同准噶尔相关材料与俄方信息并读。

1691年多伦诺尔会盟，将军事保护、封爵、补给和盟旗安排结合起来（EQ-1691-DOLONNOR）。会盟与册封可由实录及理藩院满蒙文文书确立；“入盟旗体系”并不是一次单向完成，牧地使用、领主自主性和跨部纠纷还需随后的蒙文档案追踪。1695年清廷为新的远征集中军粮、驼马和多路部队（EQ-1695-GALDAN-CAMPAIGN）。军需奏折比凯旋叙事更能让人看见运输、牲畜损耗、季节与疾病怎样限制行动；名义集结数字不能直接当作实际行军人数。

昭莫多之战后，噶尔丹西遁，1697年去世（EQ-1696-JAOMODO、EQ-1697-GALDAN-DEATH）。会战结果可以确认，地点、伤亡、各部贡献以及噶尔丹的死因和最后所在却有不同说法。清廷宣布战事告捷，并不意味着准噶尔政权消失，亦不意味着迁徙者和牧地立刻恢复旧状。1698年调整喀尔喀盟旗、台站和赈济安排（EQ-1698-KHALKHA-ADMINISTRATION），表明战时同盟开始被转化为日常治理。谕令能够证明制度意图，实际移徙、资源选择和地方关系仍应留给长期的档案与地方材料来检验。


\subsection{事件链与材料限度}
草原的盟约、牧地、驻防和季节是行动能否成形的条件。

\paragraph{康熙二十七年（1688年）：噶尔丹进攻喀尔喀，喀尔喀诸部南徙并向清廷求援，北方草原秩序改变}
\eventanchor{EQ-1688-GALDAN-KHALKHA}{噶尔丹进攻喀尔喀，喀尔喀诸部南徙并向清廷求援，北方草原秩序改变}。本节点叙述该事件本身。准噶尔、喀尔喀蒙古与清在边地道路、驻防与多方社群交错的尺度的处境首先受草原盟主竞争、牧地与贸易路线冲突促使噶尔丹东进，喀尔喀领主以清廷支援寻求生存空间制约。《清圣祖实录》康熙二十七年蒙古军报；《清史稿·噶尔丹传》；《蒙古源流》与俄罗斯使团材料留下可核的线索；可辨的后果是清廷获得介入喀尔喀政治契机，但须承担难民安置、盟旗关系和对准噶尔战争成本。原有置信注记：入侵与迁徙主干可靠；部众规模、附属关系和“归附”表述须并读蒙文与满汉文

\paragraph{康熙二十七年（1688年）：噶尔丹进攻喀尔喀，喀尔喀诸部南徙并向清廷求援，北方草原秩序改变}
\eventanchor{EQ-1688-GALDAN-KHALKHA-AFTERMATH}{噶尔丹进攻喀尔喀，喀尔喀诸部南徙并向清廷求援，北方草原秩序改变} 置于边地道路、驻防与多方社群交错的尺度中阅读，本节点追踪「噶尔丹进攻喀尔喀，喀尔喀诸部南徙并向清廷求援，北方草原秩序改变」之后可见的余波。草原盟主竞争、牧地与贸易路线冲突促使噶尔丹东进，喀尔喀领主以清廷支援寻求生存空间构成行动者准噶尔、喀尔喀蒙古与清必须面对的条件。取证从《清圣祖实录》康熙二十七年蒙古军报；《清史稿·噶尔丹传》；《蒙古源流》与俄罗斯使团材料出发，因而只能把变化谨慎地连到清廷获得介入喀尔喀政治契机，但须承担难民安置、盟旗关系和对准噶尔战争成本。原有置信注记：入侵与迁徙主干可靠；部众规模、附属关系和“归附”表述须并读蒙文与满汉文

\paragraph{康熙二十七年（1688年）：噶尔丹进攻喀尔喀，喀尔喀诸部南徙并向清廷求援，北方草原秩序改变}
\eventanchor{EQ-1688-GALDAN-KHALKHA-DECISION}{噶尔丹进攻喀尔喀，喀尔喀诸部南徙并向清廷求援，北方草原秩序改变} 标记本节点定位围绕「噶尔丹进攻喀尔喀，喀尔喀诸部南徙并向清廷求援，北方草原秩序改变」形成的处置与调度记录，涉及准噶尔、喀尔喀蒙古与清。边地道路、驻防与多方社群交错的尺度不是抽象背景：草原盟主竞争、牧地与贸易路线冲突促使噶尔丹东进，喀尔喀领主以清廷支援寻求生存空间。《清圣祖实录》康熙二十七年蒙古军报；《清史稿·噶尔丹传》；《蒙古源流》与俄罗斯使团材料所见的顺序随后指向清廷获得介入喀尔喀政治契机，但须承担难民安置、盟旗关系和对准噶尔战争成本，但原有置信注记：入侵与迁徙主干可靠；部众规模、附属关系和“归附”表述须并读蒙文与满汉文

\paragraph{康熙二十七年（1688年）：噶尔丹进攻喀尔喀，喀尔喀诸部南徙并向清廷求援，北方草原秩序改变}
在边地道路、驻防与多方社群交错的尺度，\eventanchor{EQ-1688-GALDAN-KHALKHA-PRELUDE}{噶尔丹进攻喀尔喀，喀尔喀诸部南徙并向清廷求援，北方草原秩序改变} 让准噶尔、喀尔喀蒙古与清的一个选择或压力有了可查的坐标。本节点回看「噶尔丹进攻喀尔喀，喀尔喀诸部南徙并向清廷求援，北方草原秩序改变」之前已可见的条件。前因可写为草原盟主竞争、牧地与贸易路线冲突促使噶尔丹东进，喀尔喀领主以清廷支援寻求生存空间；《清圣祖实录》康熙二十七年蒙古军报；《清史稿·噶尔丹传》；《蒙古源流》与俄罗斯使团材料限定了记录，后续变化是清廷获得介入喀尔喀政治契机，但须承担难民安置、盟旗关系和对准噶尔战争成本。原有置信注记：入侵与迁徙主干可靠；部众规模、附属关系和“归附”表述须并读蒙文与满汉文

\paragraph{康熙二十九年（1690年）：清军在乌兰布通与噶尔丹军交战，阻止其向内蒙古腹地深入但未立即消灭其军队}
\eventanchor{EQ-1690-ULAN-BUTUNG}{清军在乌兰布通与噶尔丹军交战，阻止其向内蒙古腹地深入但未立即消灭其军队} 所关联的是清与准噶尔在边地道路、驻防与多方社群交错的尺度中的行动。本节点叙述该事件本身。它从喀尔喀危机后噶尔丹持续东进，清廷须防其联络内蒙古并威胁京师北方通道展开，借《清圣祖实录》康熙二十九年八月；《清史稿·噶尔丹传》；《亲征平定朔漠方略》及俄国商旅报告进入可检验的年序，并以清军守住战略通道却未完成歼灭，战争转为更依赖补给、盟友和追击时机的长期行动影响下一段安排。原有置信注记：战役与战略结果大体可靠；胜负判断、兵力和伤亡在方略与准噶尔相关材料中不一

\paragraph{康熙二十九年（1690年）：清军在乌兰布通与噶尔丹军交战，阻止其向内蒙古腹地深入但未立即消灭其军队}
边地道路、驻防与多方社群交错的尺度中的\eventanchor{EQ-1690-ULAN-BUTUNG-AFTERMATH}{清军在乌兰布通与噶尔丹军交战，阻止其向内蒙古腹地深入但未立即消灭其军队} 不能离开清与准噶尔来解释。本节点追踪「清军在乌兰布通与噶尔丹军交战，阻止其向内蒙古腹地深入但未立即消灭其军队」之后可见的余波。喀尔喀危机后噶尔丹持续东进，清廷须防其联络内蒙古并威胁京师北方通道说明其形成的条件；《清圣祖实录》康熙二十九年八月；《清史稿·噶尔丹传》；《亲征平定朔漠方略》及俄国商旅报告提供来源路径，清军守住战略通道却未完成歼灭，战争转为更依赖补给、盟友和追击时机的长期行动则是材料能够支持的近时结果。原有置信注记：战役与战略结果大体可靠；胜负判断、兵力和伤亡在方略与准噶尔相关材料中不一

\paragraph{康熙二十九年（1690年）：清军在乌兰布通与噶尔丹军交战，阻止其向内蒙古腹地深入但未立即消灭其军队}
\eventanchor{EQ-1690-ULAN-BUTUNG-DECISION}{清军在乌兰布通与噶尔丹军交战，阻止其向内蒙古腹地深入但未立即消灭其军队} 把清与准噶尔放回边地道路、驻防与多方社群交错的尺度的道路、城镇、边界或制度网络。本节点定位围绕「清军在乌兰布通与噶尔丹军交战，阻止其向内蒙古腹地深入但未立即消灭其军队」形成的处置与调度记录。喀尔喀危机后噶尔丹持续东进，清廷须防其联络内蒙古并威胁京师北方通道。本段采用《清圣祖实录》康熙二十九年八月；《清史稿·噶尔丹传》；《亲征平定朔漠方略》及俄国商旅报告核对其顺序，随后可追到清军守住战略通道却未完成歼灭，战争转为更依赖补给、盟友和追击时机的长期行动。原有置信注记：战役与战略结果大体可靠；胜负判断、兵力和伤亡在方略与准噶尔相关材料中不一

\paragraph{康熙二十九年（1690年）：清军在乌兰布通与噶尔丹军交战，阻止其向内蒙古腹地深入但未立即消灭其军队}
本节点回看「清军在乌兰布通与噶尔丹军交战，阻止其向内蒙古腹地深入但未立即消灭其军队」之前已可见的条件，故\eventanchor{EQ-1690-ULAN-BUTUNG-PRELUDE}{清军在乌兰布通与噶尔丹军交战，阻止其向内蒙古腹地深入但未立即消灭其军队} 不只是一个年份标签。清与准噶尔在边地道路、驻防与多方社群交错的尺度承受的压力来自喀尔喀危机后噶尔丹持续东进，清廷须防其联络内蒙古并威胁京师北方通道。《清圣祖实录》康熙二十九年八月；《清史稿·噶尔丹传》；《亲征平定朔漠方略》及俄国商旅报告可回查该节点；清军守住战略通道却未完成歼灭，战争转为更依赖补给、盟友和追击时机的长期行动是其进入后续年序的方式。原有置信注记：战役与战略结果大体可靠；胜负判断、兵力和伤亡在方略与准噶尔相关材料中不一

\paragraph{康熙三十年（1691年）：康熙帝在多伦诺尔会盟，安排喀尔喀诸部入盟旗体系并调度对噶尔丹资源}
\eventanchor{EQ-1691-DOLONNOR}{康熙帝在多伦诺尔会盟，安排喀尔喀诸部入盟旗体系并调度对噶尔丹资源} 所见的行动者为清与喀尔喀蒙古诸部，空间尺度为边地道路、驻防与多方社群交错的尺度。本节点叙述该事件本身。噶尔丹战争导致喀尔喀领主寻求清廷军事支持，清廷以礼仪、封爵和补给将安全合作制度化使它成为具体事件链的一环；《清圣祖实录》康熙三十年五月；《清史稿·藩部传》；理藩院满蒙文档案中的喀尔喀盟旗与赈济文书与清廷在蒙古影响扩大，同时承担牧地、赈济和跨部冲突调停责任，统属关系仍需谈判。原有置信注记：会盟和册封可靠；“归附”不是单向完成，盟旗编制、牧地与领主自主性须以蒙文档案追踪

\paragraph{康熙三十年（1691年）：康熙帝在多伦诺尔会盟，安排喀尔喀诸部入盟旗体系并调度对噶尔丹资源}
从边地道路、驻防与多方社群交错的尺度看，\eventanchor{EQ-1691-DOLONNOR-AFTERMATH}{康熙帝在多伦诺尔会盟，安排喀尔喀诸部入盟旗体系并调度对噶尔丹资源} 留下本节点追踪「康熙帝在多伦诺尔会盟，安排喀尔喀诸部入盟旗体系并调度对噶尔丹资源」之后可见的余波的材料坐标。清与喀尔喀蒙古诸部所面对的条件是噶尔丹战争导致喀尔喀领主寻求清廷军事支持，清廷以礼仪、封爵和补给将安全合作制度化。《清圣祖实录》康熙三十年五月；《清史稿·藩部传》；理藩院满蒙文档案中的喀尔喀盟旗与赈济文书支持的结果为清廷在蒙古影响扩大，同时承担牧地、赈济和跨部冲突调停责任，统属关系仍需谈判。原有置信注记：会盟和册封可靠；“归附”不是单向完成，盟旗编制、牧地与领主自主性须以蒙文档案追踪

\paragraph{康熙三十年（1691年）：康熙帝在多伦诺尔会盟，安排喀尔喀诸部入盟旗体系并调度对噶尔丹资源}
\eventanchor{EQ-1691-DOLONNOR-DECISION}{康熙帝在多伦诺尔会盟，安排喀尔喀诸部入盟旗体系并调度对噶尔丹资源} 记录清与喀尔喀蒙古诸部在边地道路、驻防与多方社群交错的尺度的一次变化。本节点定位围绕「康熙帝在多伦诺尔会盟，安排喀尔喀诸部入盟旗体系并调度对噶尔丹资源」形成的处置与调度记录。其发生与噶尔丹战争导致喀尔喀领主寻求清廷军事支持，清廷以礼仪、封爵和补给将安全合作制度化相连；《清圣祖实录》康熙三十年五月；《清史稿·藩部传》；理藩院满蒙文档案中的喀尔喀盟旗与赈济文书把事件系回来源，清廷在蒙古影响扩大，同时承担牧地、赈济和跨部冲突调停责任，统属关系仍需谈判显示压力如何向后转移。原有置信注记：会盟和册封可靠；“归附”不是单向完成，盟旗编制、牧地与领主自主性须以蒙文档案追踪

\paragraph{康熙三十年（1691年）：康熙帝在多伦诺尔会盟，安排喀尔喀诸部入盟旗体系并调度对噶尔丹资源}
\eventanchor{EQ-1691-DOLONNOR-PRELUDE}{康熙帝在多伦诺尔会盟，安排喀尔喀诸部入盟旗体系并调度对噶尔丹资源} 的地点与行动范围属于边地道路、驻防与多方社群交错的尺度，行动者是清与喀尔喀蒙古诸部。本节点回看「康熙帝在多伦诺尔会盟，安排喀尔喀诸部入盟旗体系并调度对噶尔丹资源」之前已可见的条件。噶尔丹战争导致喀尔喀领主寻求清廷军事支持，清廷以礼仪、封爵和补给将安全合作制度化。读《清圣祖实录》康熙三十年五月；《清史稿·藩部传》；理藩院满蒙文档案中的喀尔喀盟旗与赈济文书时可见其记录边界，最小后果只能写作清廷在蒙古影响扩大，同时承担牧地、赈济和跨部冲突调停责任，统属关系仍需谈判。原有置信注记：会盟和册封可靠；“归附”不是单向完成，盟旗编制、牧地与领主自主性须以蒙文档案追踪

\paragraph{康熙三十四年（1695年）：康熙帝筹划并发动对噶尔丹的第二阶段远征，集中军粮、驼马和多路部队}
\eventanchor{EQ-1695-GALDAN-CAMPAIGN}{康熙帝筹划并发动对噶尔丹的第二阶段远征，集中军粮、驼马和多路部队} 是清与准噶尔在边地道路、驻防与多方社群交错的尺度留下的编年标记。本节点叙述该事件本身。它从乌兰布通后噶尔丹仍有机动空间，清廷需把盟旗支持、内地粮饷和季节性草原条件组织为追击走来；《清圣祖实录》康熙三十四年军务条；《清史稿·圣祖本纪》；《康熙朝满文朱批奏折全译》噶尔丹军需奏折固定可证部分，清军以更大后勤投入缩小噶尔丹行动范围，也暴露远征对运输、疫病和环境的依赖构成下一步的可见结果。原有置信注记：远征部署可靠；实际行军人数、牲畜损耗和地方征发须以军需档而非方略夸示数字判断

\paragraph{康熙三十四年（1695年）：康熙帝筹划并发动对噶尔丹的第二阶段远征，集中军粮、驼马和多路部队}
本节点追踪「康熙帝筹划并发动对噶尔丹的第二阶段远征，集中军粮、驼马和多路部队」之后可见的余波：\eventanchor{EQ-1695-GALDAN-CAMPAIGN-AFTERMATH}{康熙帝筹划并发动对噶尔丹的第二阶段远征，集中军粮、驼马和多路部队}。清与准噶尔在边地道路、驻防与多方社群交错的尺度的处境可由乌兰布通后噶尔丹仍有机动空间，清廷需把盟旗支持、内地粮饷和季节性草原条件组织为追击说明。《清圣祖实录》康熙三十四年军务条；《清史稿·圣祖本纪》；《康熙朝满文朱批奏折全译》噶尔丹军需奏折是本段材料入口，因而只能据此追踪到清军以更大后勤投入缩小噶尔丹行动范围，也暴露远征对运输、疫病和环境的依赖。原有置信注记：远征部署可靠；实际行军人数、牲畜损耗和地方征发须以军需档而非方略夸示数字判断

\paragraph{康熙三十四年（1695年）：康熙帝筹划并发动对噶尔丹的第二阶段远征，集中军粮、驼马和多路部队}
\eventanchor{EQ-1695-GALDAN-CAMPAIGN-DECISION}{康熙帝筹划并发动对噶尔丹的第二阶段远征，集中军粮、驼马和多路部队} 发生在边地道路、驻防与多方社群交错的尺度，牵动清与准噶尔。本节点定位围绕「康熙帝筹划并发动对噶尔丹的第二阶段远征，集中军粮、驼马和多路部队」形成的处置与调度记录。乌兰布通后噶尔丹仍有机动空间，清廷需把盟旗支持、内地粮饷和季节性草原条件组织为追击是其结构性条件；《清圣祖实录》康熙三十四年军务条；《清史稿·圣祖本纪》；《康熙朝满文朱批奏折全译》噶尔丹军需奏折留下的证据把读者带到清军以更大后勤投入缩小噶尔丹行动范围，也暴露远征对运输、疫病和环境的依赖。原有置信注记：远征部署可靠；实际行军人数、牲畜损耗和地方征发须以军需档而非方略夸示数字判断

\paragraph{康熙三十四年（1695年）：康熙帝筹划并发动对噶尔丹的第二阶段远征，集中军粮、驼马和多路部队}
\eventanchor{EQ-1695-GALDAN-CAMPAIGN-PRELUDE}{康熙帝筹划并发动对噶尔丹的第二阶段远征，集中军粮、驼马和多路部队} 对应本节点回看「康熙帝筹划并发动对噶尔丹的第二阶段远征，集中军粮、驼马和多路部队」之前已可见的条件。在边地道路、驻防与多方社群交错的尺度，清与准噶尔无法绕开乌兰布通后噶尔丹仍有机动空间，清廷需把盟旗支持、内地粮饷和季节性草原条件组织为追击。依据《清圣祖实录》康熙三十四年军务条；《清史稿·圣祖本纪》；《康熙朝满文朱批奏折全译》噶尔丹军需奏折，本段把后续影响限定为清军以更大后勤投入缩小噶尔丹行动范围，也暴露远征对运输、疫病和环境的依赖。原有置信注记：远征部署可靠；实际行军人数、牲畜损耗和地方征发须以军需档而非方略夸示数字判断

\paragraph{康熙三十五年（1696年）：清军在昭莫多一带击败噶尔丹主力，噶尔丹向西遁走}
\eventanchor{EQ-1696-JAOMODO}{清军在昭莫多一带击败噶尔丹主力，噶尔丹向西遁走}。本节点叙述该事件本身。清与准噶尔在边地道路、驻防与多方社群交错的尺度的处境首先受清军多路合围并获得喀尔喀和内蒙古盟友的情报、向导与补给，噶尔丹部众已受长期迁徙消耗制约。《清圣祖实录》康熙三十五年五月；《清史稿·噶尔丹传》；《亲征平定朔漠方略》与俄国边境报告留下可核的线索；可辨的后果是噶尔丹军事核心受损，准噶尔对喀尔喀的威胁减弱，但战争未以一次会战完全结束。原有置信注记：战役结局可靠；会战位置、伤亡与各部贡献应并读满文军报、蒙古材料和俄方消息

\paragraph{康熙三十五年（1696年）：清军在昭莫多一带击败噶尔丹主力，噶尔丹向西遁走}
\eventanchor{EQ-1696-JAOMODO-AFTERMATH}{清军在昭莫多一带击败噶尔丹主力，噶尔丹向西遁走} 置于边地道路、驻防与多方社群交错的尺度中阅读，本节点追踪「清军在昭莫多一带击败噶尔丹主力，噶尔丹向西遁走」之后可见的余波。清军多路合围并获得喀尔喀和内蒙古盟友的情报、向导与补给，噶尔丹部众已受长期迁徙消耗构成行动者清与准噶尔必须面对的条件。取证从《清圣祖实录》康熙三十五年五月；《清史稿·噶尔丹传》；《亲征平定朔漠方略》与俄国边境报告出发，因而只能把变化谨慎地连到噶尔丹军事核心受损，准噶尔对喀尔喀的威胁减弱，但战争未以一次会战完全结束。原有置信注记：战役结局可靠；会战位置、伤亡与各部贡献应并读满文军报、蒙古材料和俄方消息

\paragraph{康熙三十五年（1696年）：清军在昭莫多一带击败噶尔丹主力，噶尔丹向西遁走}
\eventanchor{EQ-1696-JAOMODO-DECISION}{清军在昭莫多一带击败噶尔丹主力，噶尔丹向西遁走} 标记本节点定位围绕「清军在昭莫多一带击败噶尔丹主力，噶尔丹向西遁走」形成的处置与调度记录，涉及清与准噶尔。边地道路、驻防与多方社群交错的尺度不是抽象背景：清军多路合围并获得喀尔喀和内蒙古盟友的情报、向导与补给，噶尔丹部众已受长期迁徙消耗。《清圣祖实录》康熙三十五年五月；《清史稿·噶尔丹传》；《亲征平定朔漠方略》与俄国边境报告所见的顺序随后指向噶尔丹军事核心受损，准噶尔对喀尔喀的威胁减弱，但战争未以一次会战完全结束，但原有置信注记：战役结局可靠；会战位置、伤亡与各部贡献应并读满文军报、蒙古材料和俄方消息

\paragraph{康熙三十五年（1696年）：清军在昭莫多一带击败噶尔丹主力，噶尔丹向西遁走}
在边地道路、驻防与多方社群交错的尺度，\eventanchor{EQ-1696-JAOMODO-PRELUDE}{清军在昭莫多一带击败噶尔丹主力，噶尔丹向西遁走} 让清与准噶尔的一个选择或压力有了可查的坐标。本节点回看「清军在昭莫多一带击败噶尔丹主力，噶尔丹向西遁走」之前已可见的条件。前因可写为清军多路合围并获得喀尔喀和内蒙古盟友的情报、向导与补给，噶尔丹部众已受长期迁徙消耗；《清圣祖实录》康熙三十五年五月；《清史稿·噶尔丹传》；《亲征平定朔漠方略》与俄国边境报告限定了记录，后续变化是噶尔丹军事核心受损，准噶尔对喀尔喀的威胁减弱，但战争未以一次会战完全结束。原有置信注记：战役结局可靠；会战位置、伤亡与各部贡献应并读满文军报、蒙古材料和俄方消息

\paragraph{康熙三十六年（1697年）：噶尔丹在科布多一带去世，其部众分散或向准噶尔核心地区退却}
\eventanchor{EQ-1697-GALDAN-DEATH}{噶尔丹在科布多一带去世，其部众分散或向准噶尔核心地区退却} 所关联的是准噶尔与清在边地道路、驻防与多方社群交错的尺度中的行动。本节点叙述该事件本身。它从昭莫多败后补给、盟友和机动力枯竭，噶尔丹无法重建跨部联盟展开，借《清圣祖实录》康熙三十六年；《清史稿·噶尔丹传》；俄国西伯利亚档案与蒙古口传整理进入可检验的年序，并以清廷宣布北方战事告捷，准噶尔政权仍在，蒙古牧地与人口流动后果持续存在影响下一段安排。原有置信注记：死亡事实大体可靠；死因、最后地点和部众归属说法各异，清廷凯旋叙事须与俄蒙材料互校

\paragraph{康熙三十六年（1697年）：噶尔丹在科布多一带去世，其部众分散或向准噶尔核心地区退却}
边地道路、驻防与多方社群交错的尺度中的\eventanchor{EQ-1697-GALDAN-DEATH-AFTERMATH}{噶尔丹在科布多一带去世，其部众分散或向准噶尔核心地区退却} 不能离开准噶尔与清来解释。本节点追踪「噶尔丹在科布多一带去世，其部众分散或向准噶尔核心地区退却」之后可见的余波。昭莫多败后补给、盟友和机动力枯竭，噶尔丹无法重建跨部联盟说明其形成的条件；《清圣祖实录》康熙三十六年；《清史稿·噶尔丹传》；俄国西伯利亚档案与蒙古口传整理提供来源路径，清廷宣布北方战事告捷，准噶尔政权仍在，蒙古牧地与人口流动后果持续存在则是材料能够支持的近时结果。原有置信注记：死亡事实大体可靠；死因、最后地点和部众归属说法各异，清廷凯旋叙事须与俄蒙材料互校

\paragraph{康熙三十六年（1697年）：噶尔丹在科布多一带去世，其部众分散或向准噶尔核心地区退却}
\eventanchor{EQ-1697-GALDAN-DEATH-DECISION}{噶尔丹在科布多一带去世，其部众分散或向准噶尔核心地区退却} 把准噶尔与清放回边地道路、驻防与多方社群交错的尺度的道路、城镇、边界或制度网络。本节点定位围绕「噶尔丹在科布多一带去世，其部众分散或向准噶尔核心地区退却」形成的处置与调度记录。昭莫多败后补给、盟友和机动力枯竭，噶尔丹无法重建跨部联盟。本段采用《清圣祖实录》康熙三十六年；《清史稿·噶尔丹传》；俄国西伯利亚档案与蒙古口传整理核对其顺序，随后可追到清廷宣布北方战事告捷，准噶尔政权仍在，蒙古牧地与人口流动后果持续存在。原有置信注记：死亡事实大体可靠；死因、最后地点和部众归属说法各异，清廷凯旋叙事须与俄蒙材料互校

\paragraph{康熙三十六年（1697年）：噶尔丹在科布多一带去世，其部众分散或向准噶尔核心地区退却}
本节点回看「噶尔丹在科布多一带去世，其部众分散或向准噶尔核心地区退却」之前已可见的条件，故\eventanchor{EQ-1697-GALDAN-DEATH-PRELUDE}{噶尔丹在科布多一带去世，其部众分散或向准噶尔核心地区退却} 不只是一个年份标签。准噶尔与清在边地道路、驻防与多方社群交错的尺度承受的压力来自昭莫多败后补给、盟友和机动力枯竭，噶尔丹无法重建跨部联盟。《清圣祖实录》康熙三十六年；《清史稿·噶尔丹传》；俄国西伯利亚档案与蒙古口传整理可回查该节点；清廷宣布北方战事告捷，准噶尔政权仍在，蒙古牧地与人口流动后果持续存在是其进入后续年序的方式。原有置信注记：死亡事实大体可靠；死因、最后地点和部众归属说法各异，清廷凯旋叙事须与俄蒙材料互校

\paragraph{康熙三十七年（1698年）：清廷在噶尔丹败亡后调整喀尔喀盟旗、台站与赈济安排，将战时同盟转为日常治理}
\eventanchor{EQ-1698-KHALKHA-ADMINISTRATION}{清廷在噶尔丹败亡后调整喀尔喀盟旗、台站与赈济安排，将战时同盟转为日常治理} 所见的行动者为清与喀尔喀蒙古诸部，空间尺度为边地道路、驻防与多方社群交错的尺度。本节点叙述该事件本身。战争遗留难民、牲畜损失和权力真空，迫使清廷将军事保护转为行政、礼仪和物资安排使它成为具体事件链的一环；《清圣祖实录》康熙三十七年蒙古条；《清史稿·藩部传》；理藩院满蒙文档案中的喀尔喀盟旗与赈济文书与盟旗体系获得更常态化清廷介入，但地方领主和牧民资源选择不能简化为被动纳入。原有置信注记：制度调整可由谕令和理藩院文书确认；实际牧地使用、移徙和领主关系需长期追踪

\paragraph{康熙三十七年（1698年）：清廷在噶尔丹败亡后调整喀尔喀盟旗、台站与赈济安排，将战时同盟转为日常治理}
从边地道路、驻防与多方社群交错的尺度看，\eventanchor{EQ-1698-KHALKHA-ADMINISTRATION-AFTERMATH}{清廷在噶尔丹败亡后调整喀尔喀盟旗、台站与赈济安排，将战时同盟转为日常治理} 留下本节点追踪「清廷在噶尔丹败亡后调整喀尔喀盟旗、台站与赈济安排，将战时同盟转为日常治理」之后可见的余波的材料坐标。清与喀尔喀蒙古诸部所面对的条件是战争遗留难民、牲畜损失和权力真空，迫使清廷将军事保护转为行政、礼仪和物资安排。《清圣祖实录》康熙三十七年蒙古条；《清史稿·藩部传》；理藩院满蒙文档案中的喀尔喀盟旗与赈济文书支持的结果为盟旗体系获得更常态化清廷介入，但地方领主和牧民资源选择不能简化为被动纳入。原有置信注记：制度调整可由谕令和理藩院文书确认；实际牧地使用、移徙和领主关系需长期追踪

\paragraph{康熙三十七年（1698年）：清廷在噶尔丹败亡后调整喀尔喀盟旗、台站与赈济安排，将战时同盟转为日常治理}
\eventanchor{EQ-1698-KHALKHA-ADMINISTRATION-DECISION}{清廷在噶尔丹败亡后调整喀尔喀盟旗、台站与赈济安排，将战时同盟转为日常治理} 记录清与喀尔喀蒙古诸部在边地道路、驻防与多方社群交错的尺度的一次变化。本节点定位围绕「清廷在噶尔丹败亡后调整喀尔喀盟旗、台站与赈济安排，将战时同盟转为日常治理」形成的处置与调度记录。其发生与战争遗留难民、牲畜损失和权力真空，迫使清廷将军事保护转为行政、礼仪和物资安排相连；《清圣祖实录》康熙三十七年蒙古条；《清史稿·藩部传》；理藩院满蒙文档案中的喀尔喀盟旗与赈济文书把事件系回来源，盟旗体系获得更常态化清廷介入，但地方领主和牧民资源选择不能简化为被动纳入显示压力如何向后转移。原有置信注记：制度调整可由谕令和理藩院文书确认；实际牧地使用、移徙和领主关系需长期追踪

\paragraph{康熙三十七年（1698年）：清廷在噶尔丹败亡后调整喀尔喀盟旗、台站与赈济安排，将战时同盟转为日常治理}
\eventanchor{EQ-1698-KHALKHA-ADMINISTRATION-PRELUDE}{清廷在噶尔丹败亡后调整喀尔喀盟旗、台站与赈济安排，将战时同盟转为日常治理} 的地点与行动范围属于边地道路、驻防与多方社群交错的尺度，行动者是清与喀尔喀蒙古诸部。本节点回看「清廷在噶尔丹败亡后调整喀尔喀盟旗、台站与赈济安排，将战时同盟转为日常治理」之前已可见的条件。战争遗留难民、牲畜损失和权力真空，迫使清廷将军事保护转为行政、礼仪和物资安排。读《清圣祖实录》康熙三十七年蒙古条；《清史稿·藩部传》；理藩院满蒙文档案中的喀尔喀盟旗与赈济文书时可见其记录边界，最小后果只能写作盟旗体系获得更常态化清廷介入，但地方领主和牧民资源选择不能简化为被动纳入。原有置信注记：制度调整可由谕令和理藩院文书确认；实际牧地使用、移徙和领主关系需长期追踪

\section{拉萨危机与高原道路}

1705年拉藏汗推翻第巴桑结嘉措的政权安排，拉萨的政教权力结构发生剧变。《清圣祖实录》所见的是清廷收到的西藏信息；《西藏王臣记》、驻藏来文和旅行者材料则保存不同的时间线与行动者。清廷是否知情、支持到何种程度，不能从后来的军事介入倒推。第六世达赖喇嘛仓央嘉措于1706年被押离拉萨，途中死亡还是失踪仍有相互冲突的传统；这一不确定性不是叙事的空白，而是继承合法性本身被各方争夺的证据条件。

1708年，卡尔藏嘉措在青海塔尔寺附近被认定为转世之一（EQ-1708-KELZANG）。出生和早期认定大体可靠，具体仪式、时间与承认范围不能以晚出的转世谱系倒灌。青海寺院、蒙古关系和清廷册封语言共同构成了一个跨区域的合法性网络。1717年准噶尔军进入拉萨、杀拉藏汗（EQ-1717-DZUNGAR-LHASA），暴力规模、地方态度和准噶尔治理应并读藏文、满文与旅行记录，而不能由任何一方的“平定”或“解救”词汇覆盖。

清军1718年自四川方向入藏，在喀喇乌苏一带覆败（EQ-1718-QING-TIBET-EXPEDITION）。这一失败说明高原道路、气候、兵站和信息不足不是战报中的附注，而是跨区域权力的条件。满文奏折和藏文地名、地方记忆需彼此核对，特别不能把清方报告中的向导与伤亡解释为无争议事实。1720年清军与青海蒙古盟军进入拉萨、驱逐准噶尔力量并迎立卡尔藏嘉措为第七世达赖喇嘛（EQ-1720-LHASA-RESTORATION）。进入拉萨和迎立的主干可靠；盟军、西藏地方行动者与清军的角色不能被凯旋叙事抹平。军事存在的增强也不等于西藏政教制度被一道中央命令取代，后续治理仍须在多方关系中重新安排。


\subsection{事件链与材料限度}
青海与西藏的政教、道路和驻防关系需要多语材料相互限制。

\paragraph{康熙四十四年（1705年）：拉藏汗推翻第巴桑结嘉措的政权安排，西藏政教权力结构剧烈变化}
\eventanchor{EQ-1705-LHAZANG-TAKEOVER}{拉藏汗推翻第巴桑结嘉措的政权安排，西藏政教权力结构剧烈变化}。本节点叙述该事件本身。和硕特汗国、西藏噶厦与清在边地道路、驻防与多方社群交错的尺度的处境首先受第巴与和硕特汗权长期紧张，达赖喇嘛继承和对外联络加深拉萨政治危机制约。《清圣祖实录》康熙四十四年西藏条；《清史稿·西藏传》；《西藏王臣记》藏文与清廷驻藏来文留下可核的线索；可辨的后果是西藏内部权力重组为准噶尔介入和清廷后续军事行动创造条件，清方命令不等于拉萨实际控制。原有置信注记：政变主干可靠；清廷知情、支持程度和地方反应须并读藏文、满文与旅行者材料

\paragraph{康熙四十四年（1705年）：拉藏汗推翻第巴桑结嘉措的政权安排，西藏政教权力结构剧烈变化}
\eventanchor{EQ-1705-LHAZANG-TAKEOVER-AFTERMATH}{拉藏汗推翻第巴桑结嘉措的政权安排，西藏政教权力结构剧烈变化} 置于边地道路、驻防与多方社群交错的尺度中阅读，本节点追踪「拉藏汗推翻第巴桑结嘉措的政权安排，西藏政教权力结构剧烈变化」之后可见的余波。第巴与和硕特汗权长期紧张，达赖喇嘛继承和对外联络加深拉萨政治危机构成行动者和硕特汗国、西藏噶厦与清必须面对的条件。取证从《清圣祖实录》康熙四十四年西藏条；《清史稿·西藏传》；《西藏王臣记》藏文与清廷驻藏来文出发，因而只能把变化谨慎地连到西藏内部权力重组为准噶尔介入和清廷后续军事行动创造条件，清方命令不等于拉萨实际控制。原有置信注记：政变主干可靠；清廷知情、支持程度和地方反应须并读藏文、满文与旅行者材料

\paragraph{康熙四十四年（1705年）：拉藏汗推翻第巴桑结嘉措的政权安排，西藏政教权力结构剧烈变化}
\eventanchor{EQ-1705-LHAZANG-TAKEOVER-DECISION}{拉藏汗推翻第巴桑结嘉措的政权安排，西藏政教权力结构剧烈变化} 标记本节点定位围绕「拉藏汗推翻第巴桑结嘉措的政权安排，西藏政教权力结构剧烈变化」形成的处置与调度记录，涉及和硕特汗国、西藏噶厦与清。边地道路、驻防与多方社群交错的尺度不是抽象背景：第巴与和硕特汗权长期紧张，达赖喇嘛继承和对外联络加深拉萨政治危机。《清圣祖实录》康熙四十四年西藏条；《清史稿·西藏传》；《西藏王臣记》藏文与清廷驻藏来文所见的顺序随后指向西藏内部权力重组为准噶尔介入和清廷后续军事行动创造条件，清方命令不等于拉萨实际控制，但原有置信注记：政变主干可靠；清廷知情、支持程度和地方反应须并读藏文、满文与旅行者材料

\paragraph{康熙四十四年（1705年）：拉藏汗推翻第巴桑结嘉措的政权安排，西藏政教权力结构剧烈变化}
在边地道路、驻防与多方社群交错的尺度，\eventanchor{EQ-1705-LHAZANG-TAKEOVER-PRELUDE}{拉藏汗推翻第巴桑结嘉措的政权安排，西藏政教权力结构剧烈变化} 让和硕特汗国、西藏噶厦与清的一个选择或压力有了可查的坐标。本节点回看「拉藏汗推翻第巴桑结嘉措的政权安排，西藏政教权力结构剧烈变化」之前已可见的条件。前因可写为第巴与和硕特汗权长期紧张，达赖喇嘛继承和对外联络加深拉萨政治危机；《清圣祖实录》康熙四十四年西藏条；《清史稿·西藏传》；《西藏王臣记》藏文与清廷驻藏来文限定了记录，后续变化是西藏内部权力重组为准噶尔介入和清廷后续军事行动创造条件，清方命令不等于拉萨实际控制。原有置信注记：政变主干可靠；清廷知情、支持程度和地方反应须并读藏文、满文与旅行者材料

\paragraph{康熙四十五年（1706年）：第六世达赖喇嘛仓央嘉措被押离拉萨，途中死亡或失踪，其结局与继承合法性成为争议}
\eventanchor{EQ-1706-SIXTH-DALAI}{第六世达赖喇嘛仓央嘉措被押离拉萨，途中死亡或失踪，其结局与继承合法性成为争议} 所关联的是西藏、和硕特汗国与清在边地道路、驻防与多方社群交错的尺度中的行动。本节点叙述该事件本身。它从拉藏汗以宗教合法性重组拉萨权力，清廷册封语言与西藏各派、蒙古势力认定不一致展开，借《清圣祖实录》康熙四十五年西藏条；《清史稿·西藏传》；第巴传记、藏文史书与拉达克、蒙古相关记载进入可检验的年序，并以达赖喇嘛地位争议削弱拉萨秩序，并为卡尔藏嘉措承认与清军入藏提供背景影响下一段安排。原有置信注记：被押离拉萨可靠；死亡地点、日期和是否在世的传统说法相互冲突，不能写成确定结论

\paragraph{康熙四十五年（1706年）：第六世达赖喇嘛仓央嘉措被押离拉萨，途中死亡或失踪，其结局与继承合法性成为争议}
边地道路、驻防与多方社群交错的尺度中的\eventanchor{EQ-1706-SIXTH-DALAI-AFTERMATH}{第六世达赖喇嘛仓央嘉措被押离拉萨，途中死亡或失踪，其结局与继承合法性成为争议} 不能离开西藏、和硕特汗国与清来解释。本节点追踪「第六世达赖喇嘛仓央嘉措被押离拉萨，途中死亡或失踪，其结局与继承合法性成为争议」之后可见的余波。拉藏汗以宗教合法性重组拉萨权力，清廷册封语言与西藏各派、蒙古势力认定不一致说明其形成的条件；《清圣祖实录》康熙四十五年西藏条；《清史稿·西藏传》；第巴传记、藏文史书与拉达克、蒙古相关记载提供来源路径，达赖喇嘛地位争议削弱拉萨秩序，并为卡尔藏嘉措承认与清军入藏提供背景则是材料能够支持的近时结果。原有置信注记：被押离拉萨可靠；死亡地点、日期和是否在世的传统说法相互冲突，不能写成确定结论

\paragraph{康熙四十五年（1706年）：第六世达赖喇嘛仓央嘉措被押离拉萨，途中死亡或失踪，其结局与继承合法性成为争议}
\eventanchor{EQ-1706-SIXTH-DALAI-DECISION}{第六世达赖喇嘛仓央嘉措被押离拉萨，途中死亡或失踪，其结局与继承合法性成为争议} 把西藏、和硕特汗国与清放回边地道路、驻防与多方社群交错的尺度的道路、城镇、边界或制度网络。本节点定位围绕「第六世达赖喇嘛仓央嘉措被押离拉萨，途中死亡或失踪，其结局与继承合法性成为争议」形成的处置与调度记录。拉藏汗以宗教合法性重组拉萨权力，清廷册封语言与西藏各派、蒙古势力认定不一致。本段采用《清圣祖实录》康熙四十五年西藏条；《清史稿·西藏传》；第巴传记、藏文史书与拉达克、蒙古相关记载核对其顺序，随后可追到达赖喇嘛地位争议削弱拉萨秩序，并为卡尔藏嘉措承认与清军入藏提供背景。原有置信注记：被押离拉萨可靠；死亡地点、日期和是否在世的传统说法相互冲突，不能写成确定结论

\paragraph{康熙四十五年（1706年）：第六世达赖喇嘛仓央嘉措被押离拉萨，途中死亡或失踪，其结局与继承合法性成为争议}
本节点回看「第六世达赖喇嘛仓央嘉措被押离拉萨，途中死亡或失踪，其结局与继承合法性成为争议」之前已可见的条件，故\eventanchor{EQ-1706-SIXTH-DALAI-PRELUDE}{第六世达赖喇嘛仓央嘉措被押离拉萨，途中死亡或失踪，其结局与继承合法性成为争议} 不只是一个年份标签。西藏、和硕特汗国与清在边地道路、驻防与多方社群交错的尺度承受的压力来自拉藏汗以宗教合法性重组拉萨权力，清廷册封语言与西藏各派、蒙古势力认定不一致。《清圣祖实录》康熙四十五年西藏条；《清史稿·西藏传》；第巴传记、藏文史书与拉达克、蒙古相关记载可回查该节点；达赖喇嘛地位争议削弱拉萨秩序，并为卡尔藏嘉措承认与清军入藏提供背景是其进入后续年序的方式。原有置信注记：被押离拉萨可靠；死亡地点、日期和是否在世的传统说法相互冲突，不能写成确定结论

\paragraph{康熙四十七年（1708年）：卡尔藏嘉措在青海塔尔寺附近被认定为第六世达赖喇嘛转世之一，围绕其身份形成跨区域支持网络}
\eventanchor{EQ-1708-KELZANG}{卡尔藏嘉措在青海塔尔寺附近被认定为第六世达赖喇嘛转世之一，围绕其身份形成跨区域支持网络} 所见的行动者为青海、蒙古诸部与西藏宗教网络，空间尺度为边地道路、驻防与多方社群交错的尺度。本节点叙述该事件本身。第六世达赖喇嘛结局未决和拉藏汗所立人选缺乏广泛承认，青海寺院与蒙古关系提供替代合法性中心使它成为具体事件链的一环；《清圣祖实录》康熙四十七年青海西藏条；《清史稿·西藏传》；塔尔寺藏文档案与第七世达赖喇嘛传与卡尔藏嘉措成为清廷、青海蒙古和西藏各派互动关键人物，宗教承认与军事保护相互牵连。原有置信注记：出生与早期认定大体可靠；具体认定程序、时间和各方承认不可由后世转世谱系倒推

\paragraph{康熙四十七年（1708年）：卡尔藏嘉措在青海塔尔寺附近被认定为第六世达赖喇嘛转世之一，围绕其身份形成跨区域支持网络}
从边地道路、驻防与多方社群交错的尺度看，\eventanchor{EQ-1708-KELZANG-AFTERMATH}{卡尔藏嘉措在青海塔尔寺附近被认定为第六世达赖喇嘛转世之一，围绕其身份形成跨区域支持网络} 留下本节点追踪「卡尔藏嘉措在青海塔尔寺附近被认定为第六世达赖喇嘛转世之一，围绕其身份形成跨区域支持网络」之后可见的余波的材料坐标。青海、蒙古诸部与西藏宗教网络所面对的条件是第六世达赖喇嘛结局未决和拉藏汗所立人选缺乏广泛承认，青海寺院与蒙古关系提供替代合法性中心。《清圣祖实录》康熙四十七年青海西藏条；《清史稿·西藏传》；塔尔寺藏文档案与第七世达赖喇嘛传支持的结果为卡尔藏嘉措成为清廷、青海蒙古和西藏各派互动关键人物，宗教承认与军事保护相互牵连。原有置信注记：出生与早期认定大体可靠；具体认定程序、时间和各方承认不可由后世转世谱系倒推

\paragraph{康熙四十七年（1708年）：卡尔藏嘉措在青海塔尔寺附近被认定为第六世达赖喇嘛转世之一，围绕其身份形成跨区域支持网络}
\eventanchor{EQ-1708-KELZANG-DECISION}{卡尔藏嘉措在青海塔尔寺附近被认定为第六世达赖喇嘛转世之一，围绕其身份形成跨区域支持网络} 记录青海、蒙古诸部与西藏宗教网络在边地道路、驻防与多方社群交错的尺度的一次变化。本节点定位围绕「卡尔藏嘉措在青海塔尔寺附近被认定为第六世达赖喇嘛转世之一，围绕其身份形成跨区域支持网络」形成的处置与调度记录。其发生与第六世达赖喇嘛结局未决和拉藏汗所立人选缺乏广泛承认，青海寺院与蒙古关系提供替代合法性中心相连；《清圣祖实录》康熙四十七年青海西藏条；《清史稿·西藏传》；塔尔寺藏文档案与第七世达赖喇嘛传把事件系回来源，卡尔藏嘉措成为清廷、青海蒙古和西藏各派互动关键人物，宗教承认与军事保护相互牵连显示压力如何向后转移。原有置信注记：出生与早期认定大体可靠；具体认定程序、时间和各方承认不可由后世转世谱系倒推

\paragraph{康熙四十七年（1708年）：卡尔藏嘉措在青海塔尔寺附近被认定为第六世达赖喇嘛转世之一，围绕其身份形成跨区域支持网络}
\eventanchor{EQ-1708-KELZANG-PRELUDE}{卡尔藏嘉措在青海塔尔寺附近被认定为第六世达赖喇嘛转世之一，围绕其身份形成跨区域支持网络} 的地点与行动范围属于边地道路、驻防与多方社群交错的尺度，行动者是青海、蒙古诸部与西藏宗教网络。本节点回看「卡尔藏嘉措在青海塔尔寺附近被认定为第六世达赖喇嘛转世之一，围绕其身份形成跨区域支持网络」之前已可见的条件。第六世达赖喇嘛结局未决和拉藏汗所立人选缺乏广泛承认，青海寺院与蒙古关系提供替代合法性中心。读《清圣祖实录》康熙四十七年青海西藏条；《清史稿·西藏传》；塔尔寺藏文档案与第七世达赖喇嘛传时可见其记录边界，最小后果只能写作卡尔藏嘉措成为清廷、青海蒙古和西藏各派互动关键人物，宗教承认与军事保护相互牵连。原有置信注记：出生与早期认定大体可靠；具体认定程序、时间和各方承认不可由后世转世谱系倒推

\paragraph{康熙五十六年（1717年）：准噶尔军进入拉萨并杀拉藏汗，拉萨政局再次易手}
\eventanchor{EQ-1717-DZUNGAR-LHASA}{准噶尔军进入拉萨并杀拉藏汗，拉萨政局再次易手} 是准噶尔、和硕特汗国与西藏在边地道路、驻防与多方社群交错的尺度留下的编年标记。本节点叙述该事件本身。它从拉藏汗政权缺乏广泛稳定支持，准噶尔借宗教合法性与军事力量越过高原介入走来；《清圣祖实录》康熙五十六年西藏军报；《清史稿·西藏传》；藏文《西藏王臣记》续编与准噶尔相关材料固定可证部分，清廷面临保护卡尔藏嘉措、维护边疆通道和回应准噶尔扩张的复合压力，入藏决策加速构成下一步的可见结果。原有置信注记：入拉萨和拉藏汗死亡可靠；暴力规模、民众态度与准噶尔治理需并读藏文、满文和旅行记录

\paragraph{康熙五十六年（1717年）：准噶尔军进入拉萨并杀拉藏汗，拉萨政局再次易手}
本节点追踪「准噶尔军进入拉萨并杀拉藏汗，拉萨政局再次易手」之后可见的余波：\eventanchor{EQ-1717-DZUNGAR-LHASA-AFTERMATH}{准噶尔军进入拉萨并杀拉藏汗，拉萨政局再次易手}。准噶尔、和硕特汗国与西藏在边地道路、驻防与多方社群交错的尺度的处境可由拉藏汗政权缺乏广泛稳定支持，准噶尔借宗教合法性与军事力量越过高原介入说明。《清圣祖实录》康熙五十六年西藏军报；《清史稿·西藏传》；藏文《西藏王臣记》续编与准噶尔相关材料是本段材料入口，因而只能据此追踪到清廷面临保护卡尔藏嘉措、维护边疆通道和回应准噶尔扩张的复合压力，入藏决策加速。原有置信注记：入拉萨和拉藏汗死亡可靠；暴力规模、民众态度与准噶尔治理需并读藏文、满文和旅行记录

\paragraph{康熙五十六年（1717年）：准噶尔军进入拉萨并杀拉藏汗，拉萨政局再次易手}
\eventanchor{EQ-1717-DZUNGAR-LHASA-DECISION}{准噶尔军进入拉萨并杀拉藏汗，拉萨政局再次易手} 发生在边地道路、驻防与多方社群交错的尺度，牵动准噶尔、和硕特汗国与西藏。本节点定位围绕「准噶尔军进入拉萨并杀拉藏汗，拉萨政局再次易手」形成的处置与调度记录。拉藏汗政权缺乏广泛稳定支持，准噶尔借宗教合法性与军事力量越过高原介入是其结构性条件；《清圣祖实录》康熙五十六年西藏军报；《清史稿·西藏传》；藏文《西藏王臣记》续编与准噶尔相关材料留下的证据把读者带到清廷面临保护卡尔藏嘉措、维护边疆通道和回应准噶尔扩张的复合压力，入藏决策加速。原有置信注记：入拉萨和拉藏汗死亡可靠；暴力规模、民众态度与准噶尔治理需并读藏文、满文和旅行记录

\paragraph{康熙五十六年（1717年）：准噶尔军进入拉萨并杀拉藏汗，拉萨政局再次易手}
\eventanchor{EQ-1717-DZUNGAR-LHASA-PRELUDE}{准噶尔军进入拉萨并杀拉藏汗，拉萨政局再次易手} 对应本节点回看「准噶尔军进入拉萨并杀拉藏汗，拉萨政局再次易手」之前已可见的条件。在边地道路、驻防与多方社群交错的尺度，准噶尔、和硕特汗国与西藏无法绕开拉藏汗政权缺乏广泛稳定支持，准噶尔借宗教合法性与军事力量越过高原介入。依据《清圣祖实录》康熙五十六年西藏军报；《清史稿·西藏传》；藏文《西藏王臣记》续编与准噶尔相关材料，本段把后续影响限定为清廷面临保护卡尔藏嘉措、维护边疆通道和回应准噶尔扩张的复合压力，入藏决策加速。原有置信注记：入拉萨和拉藏汗死亡可靠；暴力规模、民众态度与准噶尔治理需并读藏文、满文和旅行记录

\paragraph{康熙五十七年（1718年）：清军自四川方向入藏，在喀喇乌苏一带覆败，首批救援行动受挫}
\eventanchor{EQ-1718-QING-TIBET-EXPEDITION}{清军自四川方向入藏，在喀喇乌苏一带覆败，首批救援行动受挫}。本节点叙述该事件本身。清、准噶尔与西藏在边地道路、驻防与多方社群交错的尺度的处境首先受清军低估高原道路、气候、补给与准噶尔防御，跨区域指挥链难以及时校正制约。《清圣祖实录》康熙五十七年川藏军报；《清史稿·西藏传》；《康熙朝满文朱批奏折全译》入藏军需奏折及藏文地方记载留下可核的线索；可辨的后果是失败迫使清廷重组由青海、四川等方向的兵站与联盟，为1720年行动积累经验。原有置信注记：远征失利可靠；路线、伤亡和藏地向导问题主要见清方军报，需与藏文地名和地方传统核对

\paragraph{康熙五十七年（1718年）：清军自四川方向入藏，在喀喇乌苏一带覆败，首批救援行动受挫}
\eventanchor{EQ-1718-QING-TIBET-EXPEDITION-AFTERMATH}{清军自四川方向入藏，在喀喇乌苏一带覆败，首批救援行动受挫} 置于边地道路、驻防与多方社群交错的尺度中阅读，本节点追踪「清军自四川方向入藏，在喀喇乌苏一带覆败，首批救援行动受挫」之后可见的余波。清军低估高原道路、气候、补给与准噶尔防御，跨区域指挥链难以及时校正构成行动者清、准噶尔与西藏必须面对的条件。取证从《清圣祖实录》康熙五十七年川藏军报；《清史稿·西藏传》；《康熙朝满文朱批奏折全译》入藏军需奏折及藏文地方记载出发，因而只能把变化谨慎地连到失败迫使清廷重组由青海、四川等方向的兵站与联盟，为1720年行动积累经验。原有置信注记：远征失利可靠；路线、伤亡和藏地向导问题主要见清方军报，需与藏文地名和地方传统核对

\paragraph{康熙五十七年（1718年）：清军自四川方向入藏，在喀喇乌苏一带覆败，首批救援行动受挫}
\eventanchor{EQ-1718-QING-TIBET-EXPEDITION-DECISION}{清军自四川方向入藏，在喀喇乌苏一带覆败，首批救援行动受挫} 标记本节点定位围绕「清军自四川方向入藏，在喀喇乌苏一带覆败，首批救援行动受挫」形成的处置与调度记录，涉及清、准噶尔与西藏。边地道路、驻防与多方社群交错的尺度不是抽象背景：清军低估高原道路、气候、补给与准噶尔防御，跨区域指挥链难以及时校正。《清圣祖实录》康熙五十七年川藏军报；《清史稿·西藏传》；《康熙朝满文朱批奏折全译》入藏军需奏折及藏文地方记载所见的顺序随后指向失败迫使清廷重组由青海、四川等方向的兵站与联盟，为1720年行动积累经验，但原有置信注记：远征失利可靠；路线、伤亡和藏地向导问题主要见清方军报，需与藏文地名和地方传统核对

\paragraph{康熙五十七年（1718年）：清军自四川方向入藏，在喀喇乌苏一带覆败，首批救援行动受挫}
在边地道路、驻防与多方社群交错的尺度，\eventanchor{EQ-1718-QING-TIBET-EXPEDITION-PRELUDE}{清军自四川方向入藏，在喀喇乌苏一带覆败，首批救援行动受挫} 让清、准噶尔与西藏的一个选择或压力有了可查的坐标。本节点回看「清军自四川方向入藏，在喀喇乌苏一带覆败，首批救援行动受挫」之前已可见的条件。前因可写为清军低估高原道路、气候、补给与准噶尔防御，跨区域指挥链难以及时校正；《清圣祖实录》康熙五十七年川藏军报；《清史稿·西藏传》；《康熙朝满文朱批奏折全译》入藏军需奏折及藏文地方记载限定了记录，后续变化是失败迫使清廷重组由青海、四川等方向的兵站与联盟，为1720年行动积累经验。原有置信注记：远征失利可靠；路线、伤亡和藏地向导问题主要见清方军报，需与藏文地名和地方传统核对

\paragraph{康熙五十九年（1720年）：清军及青海蒙古盟军进入拉萨，驱逐准噶尔力量并迎立卡尔藏嘉措为第七世达赖喇嘛}
\eventanchor{EQ-1720-LHASA-RESTORATION}{清军及青海蒙古盟军进入拉萨，驱逐准噶尔力量并迎立卡尔藏嘉措为第七世达赖喇嘛} 所关联的是清、准噶尔、西藏与青海蒙古在边地道路、驻防与多方社群交错的尺度中的行动。本节点叙述该事件本身。它从1718年失败后清廷改进补给并取得青海蒙古支持，准噶尔在拉萨的统治基础也不稳展开，借《清圣祖实录》康熙五十九年西藏军报；《清史稿·西藏传》；满文军机前身档、藏文第七世达赖传与青海蒙古文书进入可检验的年序，并以清廷在拉萨的军事政治存在增强，但西藏政教制度未被单一中央命令取代，后续治理仍待安排影响下一段安排。原有置信注记：进入拉萨和迎立主干可靠；清军、蒙古盟军与西藏地方行动者角色及伤亡不能由凯旋叙事抹平

\paragraph{康熙五十九年（1720年）：清军及青海蒙古盟军进入拉萨，驱逐准噶尔力量并迎立卡尔藏嘉措为第七世达赖喇嘛}
边地道路、驻防与多方社群交错的尺度中的\eventanchor{EQ-1720-LHASA-RESTORATION-AFTERMATH}{清军及青海蒙古盟军进入拉萨，驱逐准噶尔力量并迎立卡尔藏嘉措为第七世达赖喇嘛} 不能离开清、准噶尔、西藏与青海蒙古来解释。本节点追踪「清军及青海蒙古盟军进入拉萨，驱逐准噶尔力量并迎立卡尔藏嘉措为第七世达赖喇嘛」之后可见的余波。1718年失败后清廷改进补给并取得青海蒙古支持，准噶尔在拉萨的统治基础也不稳说明其形成的条件；《清圣祖实录》康熙五十九年西藏军报；《清史稿·西藏传》；满文军机前身档、藏文第七世达赖传与青海蒙古文书提供来源路径，清廷在拉萨的军事政治存在增强，但西藏政教制度未被单一中央命令取代，后续治理仍待安排则是材料能够支持的近时结果。原有置信注记：进入拉萨和迎立主干可靠；清军、蒙古盟军与西藏地方行动者角色及伤亡不能由凯旋叙事抹平

\paragraph{康熙五十九年（1720年）：清军及青海蒙古盟军进入拉萨，驱逐准噶尔力量并迎立卡尔藏嘉措为第七世达赖喇嘛}
\eventanchor{EQ-1720-LHASA-RESTORATION-DECISION}{清军及青海蒙古盟军进入拉萨，驱逐准噶尔力量并迎立卡尔藏嘉措为第七世达赖喇嘛} 把清、准噶尔、西藏与青海蒙古放回边地道路、驻防与多方社群交错的尺度的道路、城镇、边界或制度网络。本节点定位围绕「清军及青海蒙古盟军进入拉萨，驱逐准噶尔力量并迎立卡尔藏嘉措为第七世达赖喇嘛」形成的处置与调度记录。1718年失败后清廷改进补给并取得青海蒙古支持，准噶尔在拉萨的统治基础也不稳。本段采用《清圣祖实录》康熙五十九年西藏军报；《清史稿·西藏传》；满文军机前身档、藏文第七世达赖传与青海蒙古文书核对其顺序，随后可追到清廷在拉萨的军事政治存在增强，但西藏政教制度未被单一中央命令取代，后续治理仍待安排。原有置信注记：进入拉萨和迎立主干可靠；清军、蒙古盟军与西藏地方行动者角色及伤亡不能由凯旋叙事抹平

\paragraph{康熙五十九年（1720年）：清军及青海蒙古盟军进入拉萨，驱逐准噶尔力量并迎立卡尔藏嘉措为第七世达赖喇嘛}
本节点回看「清军及青海蒙古盟军进入拉萨，驱逐准噶尔力量并迎立卡尔藏嘉措为第七世达赖喇嘛」之前已可见的条件，故\eventanchor{EQ-1720-LHASA-RESTORATION-PRELUDE}{清军及青海蒙古盟军进入拉萨，驱逐准噶尔力量并迎立卡尔藏嘉措为第七世达赖喇嘛} 不只是一个年份标签。清、准噶尔、西藏与青海蒙古在边地道路、驻防与多方社群交错的尺度承受的压力来自1718年失败后清廷改进补给并取得青海蒙古支持，准噶尔在拉萨的统治基础也不稳。《清圣祖实录》康熙五十九年西藏军报；《清史稿·西藏传》；满文军机前身档、藏文第七世达赖传与青海蒙古文书可回查该节点；清廷在拉萨的军事政治存在增强，但西藏政教制度未被单一中央命令取代，后续治理仍待安排是其进入后续年序的方式。原有置信注记：进入拉萨和迎立主干可靠；清军、蒙古盟军与西藏地方行动者角色及伤亡不能由凯旋叙事抹平

\subsection{1705—1722：高原危机、赋役承诺与继承}

1705年拉藏汗推翻第巴桑结嘉措的政权安排，拉萨的政教权力结构发生变动；1706年仓央嘉措被押离拉萨，其死亡或失踪的说法至今并不一致。1708年卡尔藏嘉措在青海塔尔寺附近被认定为转世之一，合法性因而同时联结青海寺院、蒙古关系和清廷册封语言。1717年准噶尔军进入拉萨、杀拉藏汗；清军1718年由四川方向入藏，在喀喇乌苏一带覆败。道路、气候、兵站与信息不足不是战报的附注，而是跨高原行动的条件。1720年清军与青海蒙古盟军进入拉萨、驱逐准噶尔力量并迎立卡尔藏嘉措为第七世达赖喇嘛；军事进入不等于一道命令已经取代西藏的政教制度。

帝国的另一端也显示出治理并非纯粹的胜负问题。1711年康熙宣布以当年丁数为定额，次年重申固定丁额、要求地方遵行。诏旨和会典能够确证国家作出的制度承诺，却不能直接换算为全国减负程度；造册、丁银、地丁并征和实际征收仍需户部题本及各省赋役档案来检验。1721年台湾朱一贵起事、一度进入府城，清廷随后派兵恢复控制，说明1684年后的跨海设治仍须面对垦殖、租佃、治安和征收的摩擦。官方“逆匪”的分类不能替代参与者的社会构成或诉求。

1722年康熙帝去世，胤禛即位。死期、即位和诏书的制度效力可由《清圣祖实录》、内阁大库遗诏档和《雍正朝起居注册》互校；围绕传位的传说不能仅因流传甚广便成为证据。新皇帝接手的是尚待处理的赋役、旗务、边疆和宗室问题。至此，现有清初编年材料完成了1644年至1722年的连续阅读；它没有授权本书跳过核验而把其后的年代写成既成事实。

\paragraph{材料位置。} 《清圣祖实录》呈现的是清廷收到信息、发出谕旨和编排凯旋的方式；藏文传统、满文奏折、寺院材料、旅行记录、地方志和赋役档案分别补足或挑战其视野。对于高原行动、丁额及台湾起事，材料所能证明的范围各不相同，不能用一类文书取代所有地方经验。

\section{定额的承诺与继承的不确定}

康熙五十年宣布以当年丁数为定额，后生人丁不再加征丁税；次年又重申固定丁额、要求地方遵行（EQ-1711-TAX-FREEZE、EQ-1712-TAX-EDICT）。《清圣祖实录》的谕旨可以确证中央规范，《大清会典事例》界定制度语言，户部题本及江苏、直隶等省的赋役档案才可能检验造册、丁银、地丁并征和实际征收。因而“滋生人丁，永不加赋”首先是一项政策承诺，不是可直接换算为全国减负程度的数字。各省执行时间、旧额、缺额和摊派的差异，要求把诏令、部议、地方账册和民众负担分开书写。它改变了丁额增长的制度预期，并为雍正时期的赋役调整留下条件，却没有消灭地方财政压力。

康熙六十一年十一月，康熙帝去世，皇四子胤禛即位（EQ-1722-KANGXI-DEATH）。死期、即位与诏书的制度效力，可由《清圣祖实录》、内阁大库遗诏档和《雍正朝起居注册》互校。《清史稿》提供了后出的编年索引，却不能使传位过程中的每一个传说自动成立。康熙晚年缺少公开稳定的继承机制，宗室、旗人和官僚围绕信息形成竞争和猜测；这一背景可以说明何以出现争议，不能把争议本身当成阴谋的证据。新皇帝继承的，是尚待处理的财政、旗务、边疆和宗室问题，而非一套已经没有摩擦的盛世机器。


\subsection{事件链与材料限度}
赋役、继承与地方执行之间的联系，须在具体账册、诏令和社会实践中逐项追踪。

\paragraph{康熙五十年（1711年）：康熙帝宣布以康熙五十年丁数为定额，后生人丁不再加征丁税}
\eventanchor{EQ-1711-TAX-FREEZE}{康熙帝宣布以康熙五十年丁数为定额，后生人丁不再加征丁税}。本节点叙述该事件本身。清在京师与多地军政线路相接的尺度的处境首先受长期战争后清廷寻求稳定赋役预期，人口登记与丁税征收的摩擦已显现制约。《清圣祖实录》康熙五十年丁银谕；《清史稿·食货志》；户部档案与《大清会典事例》赋役条留下可核的线索；可辨的后果是政策改变丁额增长的制度预期，为雍正朝摊丁入亩等改革留下条件；实际减负须按省县账册验证。原有置信注记：诏令可靠；各省执行时间、丁银与地丁并征口径和实际负担不能由诏令推定

\paragraph{康熙五十年（1711年）：康熙帝宣布以康熙五十年丁数为定额，后生人丁不再加征丁税}
\eventanchor{EQ-1711-TAX-FREEZE-AFTERMATH}{康熙帝宣布以康熙五十年丁数为定额，后生人丁不再加征丁税} 置于京师与多地军政线路相接的尺度中阅读，本节点追踪「康熙帝宣布以康熙五十年丁数为定额，后生人丁不再加征丁税」之后可见的余波。长期战争后清廷寻求稳定赋役预期，人口登记与丁税征收的摩擦已显现构成行动者清必须面对的条件。取证从《清圣祖实录》康熙五十年丁银谕；《清史稿·食货志》；户部档案与《大清会典事例》赋役条出发，因而只能把变化谨慎地连到政策改变丁额增长的制度预期，为雍正朝摊丁入亩等改革留下条件；实际减负须按省县账册验证。原有置信注记：诏令可靠；各省执行时间、丁银与地丁并征口径和实际负担不能由诏令推定

\paragraph{康熙五十年（1711年）：康熙帝宣布以康熙五十年丁数为定额，后生人丁不再加征丁税}
\eventanchor{EQ-1711-TAX-FREEZE-DECISION}{康熙帝宣布以康熙五十年丁数为定额，后生人丁不再加征丁税} 标记本节点定位围绕「康熙帝宣布以康熙五十年丁数为定额，后生人丁不再加征丁税」形成的处置与调度记录，涉及清。京师与多地军政线路相接的尺度不是抽象背景：长期战争后清廷寻求稳定赋役预期，人口登记与丁税征收的摩擦已显现。《清圣祖实录》康熙五十年丁银谕；《清史稿·食货志》；户部档案与《大清会典事例》赋役条所见的顺序随后指向政策改变丁额增长的制度预期，为雍正朝摊丁入亩等改革留下条件；实际减负须按省县账册验证，但原有置信注记：诏令可靠；各省执行时间、丁银与地丁并征口径和实际负担不能由诏令推定

\paragraph{康熙五十年（1711年）：康熙帝宣布以康熙五十年丁数为定额，后生人丁不再加征丁税}
在京师与多地军政线路相接的尺度，\eventanchor{EQ-1711-TAX-FREEZE-PRELUDE}{康熙帝宣布以康熙五十年丁数为定额，后生人丁不再加征丁税} 让清的一个选择或压力有了可查的坐标。本节点回看「康熙帝宣布以康熙五十年丁数为定额，后生人丁不再加征丁税」之前已可见的条件。前因可写为长期战争后清廷寻求稳定赋役预期，人口登记与丁税征收的摩擦已显现；《清圣祖实录》康熙五十年丁银谕；《清史稿·食货志》；户部档案与《大清会典事例》赋役条限定了记录，后续变化是政策改变丁额增长的制度预期，为雍正朝摊丁入亩等改革留下条件；实际减负须按省县账册验证。原有置信注记：诏令可靠；各省执行时间、丁银与地丁并征口径和实际负担不能由诏令推定

\paragraph{康熙五十一年（1712年）：清廷再次谕令固定丁额并要求地方遵行，围绕丁银征收与造册形成持续行政问题}
\eventanchor{EQ-1712-TAX-EDICT}{清廷再次谕令固定丁额并要求地方遵行，围绕丁银征收与造册形成持续行政问题} 所关联的是清在省、府县与地方社会相互牵动的尺度中的行动。本节点叙述该事件本身。它从上一年诏令需经户部、督抚和州县造册落实，旧有缺额、摊派和财政压力造成执行摩擦展开，借《清圣祖实录》康熙五十一年赋役谕；《清史稿·食货志》；户部题本与江苏、直隶等省赋役档进入可检验的年序，并以“永不加赋”成为政策语言与后续改革依据，但登记人口、实际税负和地区差异不能被单一口号遮蔽影响下一段安排。原有置信注记：重复申谕可靠；它证明中央规范而非各地立即停止加派，需追踪地方账册与奏报

\paragraph{康熙五十一年（1712年）：清廷再次谕令固定丁额并要求地方遵行，围绕丁银征收与造册形成持续行政问题}
省、府县与地方社会相互牵动的尺度中的\eventanchor{EQ-1712-TAX-EDICT-AFTERMATH}{清廷再次谕令固定丁额并要求地方遵行，围绕丁银征收与造册形成持续行政问题} 不能离开清来解释。本节点追踪「清廷再次谕令固定丁额并要求地方遵行，围绕丁银征收与造册形成持续行政问题」之后可见的余波。上一年诏令需经户部、督抚和州县造册落实，旧有缺额、摊派和财政压力造成执行摩擦说明其形成的条件；《清圣祖实录》康熙五十一年赋役谕；《清史稿·食货志》；户部题本与江苏、直隶等省赋役档提供来源路径，“永不加赋”成为政策语言与后续改革依据，但登记人口、实际税负和地区差异不能被单一口号遮蔽则是材料能够支持的近时结果。原有置信注记：重复申谕可靠；它证明中央规范而非各地立即停止加派，需追踪地方账册与奏报

\paragraph{康熙五十一年（1712年）：清廷再次谕令固定丁额并要求地方遵行，围绕丁银征收与造册形成持续行政问题}
\eventanchor{EQ-1712-TAX-EDICT-DECISION}{清廷再次谕令固定丁额并要求地方遵行，围绕丁银征收与造册形成持续行政问题} 把清放回省、府县与地方社会相互牵动的尺度的道路、城镇、边界或制度网络。本节点定位围绕「清廷再次谕令固定丁额并要求地方遵行，围绕丁银征收与造册形成持续行政问题」形成的处置与调度记录。上一年诏令需经户部、督抚和州县造册落实，旧有缺额、摊派和财政压力造成执行摩擦。本段采用《清圣祖实录》康熙五十一年赋役谕；《清史稿·食货志》；户部题本与江苏、直隶等省赋役档核对其顺序，随后可追到“永不加赋”成为政策语言与后续改革依据，但登记人口、实际税负和地区差异不能被单一口号遮蔽。原有置信注记：重复申谕可靠；它证明中央规范而非各地立即停止加派，需追踪地方账册与奏报

\paragraph{康熙五十一年（1712年）：清廷再次谕令固定丁额并要求地方遵行，围绕丁银征收与造册形成持续行政问题}
本节点回看「清廷再次谕令固定丁额并要求地方遵行，围绕丁银征收与造册形成持续行政问题」之前已可见的条件，故\eventanchor{EQ-1712-TAX-EDICT-PRELUDE}{清廷再次谕令固定丁额并要求地方遵行，围绕丁银征收与造册形成持续行政问题} 不只是一个年份标签。清在省、府县与地方社会相互牵动的尺度承受的压力来自上一年诏令需经户部、督抚和州县造册落实，旧有缺额、摊派和财政压力造成执行摩擦。《清圣祖实录》康熙五十一年赋役谕；《清史稿·食货志》；户部题本与江苏、直隶等省赋役档可回查该节点；“永不加赋”成为政策语言与后续改革依据，但登记人口、实际税负和地区差异不能被单一口号遮蔽是其进入后续年序的方式。原有置信注记：重复申谕可靠；它证明中央规范而非各地立即停止加派，需追踪地方账册与奏报

\paragraph{康熙六十一年十一月（1722年）：康熙帝去世，皇四子胤禛即位为雍正帝，早期清朝的长期统治转入新的继承阶段}
\eventanchor{EQ-1722-KANGXI-DEATH}{康熙帝去世，皇四子胤禛即位为雍正帝，早期清朝的长期统治转入新的继承阶段} 所见的行动者为清，空间尺度为京师与多地军政线路相接的尺度。本节点叙述该事件本身。康熙晚年未公开稳定继承机制，宗室、旗人和官僚围绕皇位信息形成竞争与猜测使它成为具体事件链的一环；《清圣祖实录》康熙六十一年十一月；《清史稿·圣祖本纪、世宗本纪》；《雍正朝起居注册》与内阁大库遗诏档与雍正即位后将处理财政、旗务、边疆和宗室问题；传位细节不能以阴谋叙事代替可核文书。原有置信注记：去世与即位日期可靠；遗诏、传位过程和储位争议受政治清算与后世传说影响，须以档案版本互校

\paragraph{康熙六十一年十一月（1722年）：康熙帝去世，皇四子胤禛即位为雍正帝，早期清朝的长期统治转入新的继承阶段}
从京师与多地军政线路相接的尺度看，\eventanchor{EQ-1722-KANGXI-DEATH-AFTERMATH}{康熙帝去世，皇四子胤禛即位为雍正帝，早期清朝的长期统治转入新的继承阶段} 留下本节点追踪「康熙帝去世，皇四子胤禛即位为雍正帝，早期清朝的长期统治转入新的继承阶段」之后可见的余波的材料坐标。清所面对的条件是康熙晚年未公开稳定继承机制，宗室、旗人和官僚围绕皇位信息形成竞争与猜测。《清圣祖实录》康熙六十一年十一月；《清史稿·圣祖本纪、世宗本纪》；《雍正朝起居注册》与内阁大库遗诏档支持的结果为雍正即位后将处理财政、旗务、边疆和宗室问题；传位细节不能以阴谋叙事代替可核文书。原有置信注记：去世与即位日期可靠；遗诏、传位过程和储位争议受政治清算与后世传说影响，须以档案版本互校

\paragraph{康熙六十一年十一月（1722年）：康熙帝去世，皇四子胤禛即位为雍正帝，早期清朝的长期统治转入新的继承阶段}
\eventanchor{EQ-1722-KANGXI-DEATH-DECISION}{康熙帝去世，皇四子胤禛即位为雍正帝，早期清朝的长期统治转入新的继承阶段} 记录清在京师与多地军政线路相接的尺度的一次变化。本节点定位围绕「康熙帝去世，皇四子胤禛即位为雍正帝，早期清朝的长期统治转入新的继承阶段」形成的处置与调度记录。其发生与康熙晚年未公开稳定继承机制，宗室、旗人和官僚围绕皇位信息形成竞争与猜测相连；《清圣祖实录》康熙六十一年十一月；《清史稿·圣祖本纪、世宗本纪》；《雍正朝起居注册》与内阁大库遗诏档把事件系回来源，雍正即位后将处理财政、旗务、边疆和宗室问题；传位细节不能以阴谋叙事代替可核文书显示压力如何向后转移。原有置信注记：去世与即位日期可靠；遗诏、传位过程和储位争议受政治清算与后世传说影响，须以档案版本互校

\paragraph{康熙六十一年十一月（1722年）：康熙帝去世，皇四子胤禛即位为雍正帝，早期清朝的长期统治转入新的继承阶段}
\eventanchor{EQ-1722-KANGXI-DEATH-PRELUDE}{康熙帝去世，皇四子胤禛即位为雍正帝，早期清朝的长期统治转入新的继承阶段} 的地点与行动范围属于京师与多地军政线路相接的尺度，行动者是清。本节点回看「康熙帝去世，皇四子胤禛即位为雍正帝，早期清朝的长期统治转入新的继承阶段」之前已可见的条件。康熙晚年未公开稳定继承机制，宗室、旗人和官僚围绕皇位信息形成竞争与猜测。读《清圣祖实录》康熙六十一年十一月；《清史稿·圣祖本纪、世宗本纪》；《雍正朝起居注册》与内阁大库遗诏档时可见其记录边界，最小后果只能写作雍正即位后将处理财政、旗务、边疆和宗室问题；传位细节不能以阴谋叙事代替可核文书。原有置信注记：去世与即位日期可靠；遗诏、传位过程和储位争议受政治清算与后世传说影响，须以档案版本互校

\begin{figure}[htbp]
  \centering
  \begin{tikzpicture}[
    x=1cm,y=1cm,
    place/.style={draw=timeline, fill=paper, rounded corners=2pt, align=center,
      inner sep=3pt, font=\small},
    qing/.style={place, draw=dynasty, fill=dynasty!13},
    frontier/.style={place, draw=timeline, fill=timeline!13},
    sea/.style={place, draw=source, fill=source!10},
    note/.style={align=center, font=\scriptsize\color{source}}
  ]
    \fill[paper] (-.7,-4.0) rectangle (7.5,3.8);
    \draw[dynasty, line width=.8pt] (-.7,3.8) -- (7.5,3.8);
    \node[font=\bfseries\large, text=dynasty] at (3.4,3.35)
      {早期清朝的边疆治理：蒙古、西藏与台湾};

    \node[qing] (beijing) at (3.35,2.55) {北京\\军政、礼仪与文书枢纽};

    \node[frontier] (dzungar) at (.25,1.25) {准噶尔\\1688后};
    \node[frontier] (khalkha) at (2.15,.75) {喀尔喀诸部\\南徙、求援};
    \node[frontier] (dolonnor) at (3.85,.9) {多伦诺尔\\1691会盟};
    \node[frontier] (banner) at (5.65,1.15) {盟旗、台站、赈济\\战后日常治理};

    \node[frontier] (qinghai) at (.65,-1.55) {青海蒙古\\寺院网络};
    \node[frontier] (lhasa) at (2.75,-2.2) {拉萨\\1717--1720};
    \node[frontier] (tibetnote) at (5.25,-2.2) {政教、军行与\\地方协商的多重网络};

    \node[sea] (penghu) at (4.85,-.55) {澎湖\\1683};
    \node[sea] (taiwan) at (6.35,-.85) {台湾\\府县、驻防\\1684后};

    \draw[->, timeline, line width=1pt] (dzungar) -- node[above,sloped,note]
      {战争与迁徙} (khalkha);
    \draw[->, dynasty, line width=1.15pt] (khalkha) -- node[above,sloped,note]
      {会盟、编旗的过程} (dolonnor);
    \draw[->, dynasty, line width=1.15pt] (dolonnor) -- node[above,sloped,note]
      {理藩院事务（示意）} (banner);
    \draw[dynasty, dashed] (beijing) -- (dolonnor);

    \draw[->, timeline, dashed, line width=1pt] (qinghai) -- node[below,sloped,note]
      {盟军、寺院与交通网络} (lhasa);
    \draw[->, dynasty, line width=1.15pt] (beijing) to[out=-150,in=70] node[left,note]
      {1720年军行与迎立（示意）} (lhasa);
    \draw[timeline, dashed] (lhasa) -- (tibetnote);

    \draw[->, source, line width=1pt] (penghu) -- node[above,sloped,note]
      {海战、投降与接收} (taiwan);
    \draw[dynasty, dashed] (beijing) to[out=-30,in=110] node[right,note]
      {设治、驻防的文书路径（示意）} (taiwan);

    \node[draw=source, fill=paper, rounded corners=3pt, align=center,
      text=source, font=\small] at (5.7,-3.25)
      {没有连续国界：节点、虚线与箭头\\只表示档案可见的事务、路线或关系};
  \end{tikzpicture}
  \caption{早期清朝的蒙古、西藏与台湾治理（约1688--1721，示意）。图以会盟、盟旗／台站、军行、寺院网络、海战及设治等可辨事务为单位，而非以色块制造连续疆界。蒙古部分据《清圣祖实录》康熙二十七至三十七年蒙古军报和理藩院满、汉、蒙文盟旗及赈济档，并参照《蒙古源流》与俄国使团／边境材料；西藏部分据《清圣祖实录》康熙五十九年军报、满文军机前身档、藏文《西藏王臣记》、塔尔寺藏文档与青海蒙古文书；台湾部分据《清圣祖实录》康熙二十二、二十三年条、施琅《飞报澎湖大捷疏》、《台湾府志》、福建督抚奏折及台湾契约、港口与荷兰／日文海洋材料。范围只标出少数关系节点和约略路线，未标按比例的道路、牧地、海域或行政分界；箭头不等同主权、常态化管辖或各地社会的接受，尤其不能由盟旗、寺院或府县名目推导现代边界。}
  \label{fig:early-qing-frontier-governance}
\end{figure}

\section{战争、财用与省级治理的缓慢成形}

一六四四年北京易手后，摆在新政权面前的并非一张可以立即征税的版图，而是一座要供养军队、收发文书并安顿官署的城市。仓廒中的粮食、运河上的漕船、熟悉钱粮与刑名的旧吏，以及城外仍在移动的部队，彼此并不相属，却都决定北京能否从行军终点变成政令的出发处。摄政者可以招纳旧官、发布征发和安民的命令；命令能否越过关卡、在府县被改写成粮差、夫役或留仓，则取决于地方官、绅士、书吏、船户和军队之间当时能够维持的关系。接收明代的部院与名目，因而不是取得现成财源，而是开始学习怎样在战争尚未结束时使用一套已被危机撕扯过的行政手脚。

这层困难沿着运河和长江向南展开。一六四五年南京失守，改变了江南的政治中心与运输枢纽，却没有使江浙、闽粤和西南同时变成可以稳定征解的地区。军队要沿水陆继续推进，城池的守备者和居民则要在围城、征粮、逃亡与报复之间决定是否留下。江阴等地的战事提醒人们：军报里的“恢复”首先是一个军事节点；对县城周围的村落而言，它还可能意味着田地无人耕作、船只被征用、家族分散，或必须向新来的驻军重新交纳物资。战争中的粮饷不能只从户部或将领的角度理解。每一次催运都要经过堤岸、渡口、仓门和人力，而这些条件既限制军队，也改变地方社会保存生计的办法。

清廷正是在这种不齐整的地面上组织省级行政。督抚、道府州县和将领并不是一条笔直的传令线：上报军情与钱粮的人会考虑考成、请饷和自保，转运的人会受到水位、道路、船力和地方库存的约束，承担差役的人又未必能在文书中留下名字。北京得到的奏报可以显示某项调拨已经被请求、批准或催办，却不能让我们断言粮食已经抵达哪一支队伍，更不能由一次报解推定一省居民共同承担了多少。省份在这里不是预先完成的财政容器，而是中央命令、驻军需要和地方中介反复碰撞出来的治理空间。

西南留下了最清楚、也最难用一个结局收束的例子。永历朝廷在滇西与缅甸方向的流转，使道路、边境通行和供给同时成为政治问题；一六六二年永历帝被处死，消除了一个明朝皇帝的名义中心，却没有消除云南的军队、流亡者和商路。吴三桂等军镇原本承担征服与驻守的任务，也掌握同地方钱粮、人事和旧部相连的渠道。对朝廷而言，它们既是维持西南的工具，也是长期无法核算和收回的负担；对州县、商人和乡绅而言，它们则是必须打交道的实际权力。把这种安排称作单纯的中央放权，会抹去它从战争中生出、又依赖地方供给才得以维持的性质。

沿海的经验使“供给”有了另一种尺度。一六六一年以来的迁界，意在削弱郑氏同闽粤岸上社会的联系；渔盐、耕作、渡海、亲族往来和港口买卖却不会按同一速度停止或恢复。有人被迫迁离，有人在限制变动后设法返乡，也有人把劳作和关系移到别处。官府的禁令能够说明国家试图截断什么，地方志的复业叙述则常在事后以恢复为线索；两者之间，究竟有哪些村落失去田地、哪些船户仍能航行、哪些家庭重新组合，必须由契约、港口记录和地方材料逐处追问。海禁与复界不是陆地战争的附注，而是把军费、海防和人口流动同时压到沿岸社会身上的政策。

康熙十二年撤藩时，这些临时安排的代价终于集中显露。朝廷想收回三藩的军政与饷源，面对的不是三块可以从地图上擦除的领地，而是云南、广东、福建已有的军队、仓储、家属、债务、运输路线和地方往来。吴三桂起兵以后，战线把西南山路、长江节点、福建港口和西北通道牵在一起；哪一方能够守住城池，往往取决于粮草、人夫、舟船和投降者是否真的能够安置。清军一六八一年进入昆明，结束了吴周政权，却不能使多年征发留下的道路损耗、荒地、驻军开销与重新招徕人力同时消失。战后所谓垦复和复业，首先是官府与地方试图重建可耕、可运、可征关系的过程，不是人口和田赋已经恢复如初的证明。

一六八三年澎湖之战以后，郑氏政权降服，次年台湾设府，海峡两岸的军事关系随之改写。行政链可以向东延伸，驻防、渡海限制、土地开垦和港口贸易的实际形状却仍须在岛内村社、移民、船户与原住民社群的接触中逐步形成。设府同撤藩一样，表明朝廷能够重新安排名义上的权责；它不能替我们决定每一处土地如何取得、每一艘船为何出航，也不能把不同社群的处境写成一种接受或拒绝。战争结束后仍不断发生的渡海与迁居，正说明财政、军事和地方生计并没有各自留在边界清楚的领域里。

康熙晚年固定丁额的承诺，可以放在这段经历的尾声来读。它试图给地方赋役一个较可预期的尺度，也为随后调整地丁关系留下条件；但定额写入谕旨，不等于各省旧额、欠额、摊派和实收立即一致。清初战争所塑造的国家，获得了把命令、军队和钱粮更远地联结起来的能力，代价则由道路上的劳力、仓储中的粮食、迁居者的家计和地方中介的筹措共同承受。省级治理的力量正在于它能把这些资源接续起来；它的限度也正在于，任何一道来自北京的命令，都还必须穿过具体的山路、河道、港湾和人群，才可能取得有限而不均匀的结果。


\section{回看：为编年服务的政治与边地问题}

以下两节不另行取代上述时间线，而是回到其中反复出现的接收、供给、迁徙和多语交涉，说明同一批事件在制度与地方关系上留下的压力。
\subsection{从山海关到北京：一座城如何变成朝廷}

顺治元年四月，山海关把三种彼此利用、也彼此不信任的力量挤到了一起。李自成进入北京后，关宁军同大顺在兵饷和统属上已经决裂；多尔衮则把吴三桂的处境变成八旗越关的机会。会战打开的是通往华北的道路，不是天下已经归属的证明。五月入北京之后，清军首先要接收的是仓储、官署、粮道和等待观望的官员；冬季的招纳、征发与整顿官署，才试图把一支行军的军队变成能收发文书的政权。

这份接收有很具体的物质尺度。驻军的口粮、马匹的草料、南下所需的船只和车夫，都不能由北京城门已经易帜而自动出现；而搜集这些资源又会把新政权带到州县的仓廒、渡口和里甲面前。对朝廷而言，旧有的赋役名目、漕运节点和能办案的书吏，是让军令继续向外走的条件；对被征发的人和被留任的官吏而言，它们首先意味着饷额、差役与人身安全尚未有定数。故而“接收明制”既是制度上的借用，也是一次围绕粮、船、人力与地方信誉不断谈判的过程。

北京并非唯一的政治空间。顺治帝迁驻北京，使皇帝、部院和礼制中心汇合；盛京仍保有旗务、皇室与王朝记忆的重量。与此同时，江南的水路、两广和西南的山道、福建沿海的港口与岛屿，各有不同的通行条件、供给办法和对手。摄政政权能在京师发出同一格式的命令，却不能使这些区域以同一速度进入其财政和军政网络。把清初写成由北京向四方均匀铺开的统治，会把多区域的战争误写成一张已经完成的行政地图。

摄政结束、顺治亲政，也没有消除这种多重结构。北京需要熟悉赋役、漕运与刑名的书吏和降官，前线持续索要粮饷、船只和兵丁；皇帝的诏令必须经过宗室、旗务、满汉官僚、将领以及地方中介，才可能抵达一州一县。康熙幼年时的四辅政大臣、鳌拜被拘捕和皇帝重新掌握人事，显示宫廷权力的重新分配；它们并不证明命令从宫门笔直地变成地方的服从。每一次任命与调度仍受道路、军费、旧属关系和战区局势牵制。

征服的代价也写在城池和身体上。南京失守以后，弘光、隆武、绍武、永历等政治中心仍在东南和西南移动；一座府城易手不等于乡村、岛屿、山道和税源同时易手。薙发易服的命令把政治归属做成可见的标记，却在各地同守城、征发、逃亡和报复纠缠。有人借服从求得暂时的秩序，有人因旧主、亲属、粮饷或城防而延宕、离散或抵抗；这些选择不能被压成一条从命到不从命的整齐线。

扬州、江阴和广州的杀戮与损毁不能被胜利叙事抹去，但也不宜以一个貌似精确的总数代替不同地方的经历。城守记述、幸存者叙事和后修方志使围城、屠戮、逃亡与毁坏可见；它们又常带有传抄、追忆、纪念共同体和数字修辞的痕迹。清廷战报所记的攻取与招降，也不是伤亡普查。能够确定的是征服暴力深刻改变了城市和周边社会；不能由现存文本合算一项代表全境的死亡数字。直到永历政权覆亡、夔东据点被攻破，郑氏仍据海上与台湾，所谓统一始终是多战区逐段取得的结果。

\subsection{材料与争议}

《清世祖实录》、内国史院满文档案和题本能追踪北京政权怎样制作命令、接收官署和调兵筹饷；它们首先是新政权的文书，而非命令在远方已经实现的证明。满、汉文本也不能未经核对便互作替身：官号、地名和分类用语的对应，会影响我们所见的行动者与行政范围。题本所留下的是呈报、批示和案卷得以进入中枢的链条，尤其难以覆盖未能上报的逃散、拖延和私下交涉。

南明奏疏、地方记事、城守记录与后出的方志保存败亡朝廷和城市共同体的不同时间，也有传抄、追忆和地方立场。它们与官书相对读，不是为了拼出一种全知的中性叙事，而是为了分别检验城池何时易手、何种征发可见、哪些人被写入或留在沉默中。山海关的求援文书、各军兵数与攻城伤亡因此不能被写成无争议的精确场景；同样，中央命令不能代替地方接受度的统计。能够确定的是征服的暴力、反复与制度接收同时发生；各地的服从、损失与恢复程度则须随材料的保存状况分别判断。

\subsection{交叉阅读窗口：一纸薙发谕如何把征服写进身体}

顺治二年六月谕旨以“遵依者为我国之民，迟疑者同逆命之寇”划分服从者与敌人。短句出自《清世祖实录》的朝廷命令，和内国史院档案一样，首先可证明新政权如何发出、归类和威吓；它的对偶和二分法把发式压缩成政治归属，因而具有比行政术语更强的逼迫力。它不是江南每一处执行情形的逐日记录，更不能说明每个遵从者都出于同一种判断。命令越过城门、渡口和驿路时，会同驻军补给、城防崩解、地方差役与家庭避难交织在一起。

钱穆会追问的不是这句谕令是否“有效”，而是摄政政权怎样借明代官署、旗务与礼制把命令变成可持续的政治关系；吕思勉式的问题则把视线移向军饷、赋役、家庭逃散与城镇社会如何限制或改变这种关系。两种提问都不能替代证据：实录和档案证明命令与中央分类，地方记录仅能局部检验收到、抵抗或服从的过程。这个窗口因此不把一句严厉的国家声音误作征服已经完成的证明。

\subsection{把胜利留在地方：西南的军镇与海峡的往返}

一六六一年，\eventanchor{EQ-1661-KANGXI-ACCESSION}{康熙即位（1661）}时，北京交到八岁玄烨和辅政大臣手中的，不只是宫廷的人事安排，还有仍在西南、海峡与各省军镇之间迟缓传递的军情。北京能够发出任命和处分，仍须经旗务、部院、督抚与将领把命令送进不同战区；接到的回报又会被行军、季节与呈报者的利害改变。次年，\eventanchor{EQ-1662-YONGLI-EXECUTION}{永历帝被处死（1662）}，南明最后一个皇帝名义中心消失，但朱由榔自缅甸方向被押回的路线，已把吴三桂集团、清军、缅甸宫廷与边境通行牵在一起。处决是一个政治节点，不是云南西部的军队、流亡者、贸易与安置问题同时终结的时刻。

康熙十二年春，撤藩的题本从北京发出时，云南、广东和福建早已不是可以在部议中随意腾挪的空格。吴三桂、尚可喜、耿氏家族所维系的，是入关征服后留在当地的军队、粮饷渠道、亲属与旧部的前途，也包括州县官、商人和乡绅不得不同这些军镇打交道的日常。中央希望把长期的驻军财政与任命权收回，首先碰到的却是很具体的问题：哪一处仓粮可供转运，哪一路驿道能让新调来的部队通过，原先领饷的人由谁安置。因而\eventanchor{EQ-1673-FEUDATORY-ABOLITION}{撤藩决定（1673）}不是一道命令从北京落到三处属地的单线故事，而是把这些已被嵌入地方的安排骤然推入不确定之中。

吴三桂在云南举兵后，西南的战争也不能只按朝廷与一位将领的胜负来读。滇、黔的山路、渡口和城镇把军队、运夫、商旅、逃难者与地方守备卷在一起；一支部队能否继续向前，取决于粮草是否赶得上、地方是否还能征到人力，也取决于守城者、降将和州县官在消息尚未明朗时如何选择。清军最终进入昆明，改变了吴氏政权的军事前景，却没有自动消除驻军的开销、垦复的压力和战后重建税源的困难。把战争写成若干城池颜色的转换，会遗漏最持久的一层后果：新政权仍须在被扰乱的道路上运输，在经历征发的人群中重新建立可征可管的关系。

官修的战事汇编和实录擅长保存北京收到的捷报、招降、任免与奖惩，因而能让我们追踪朝廷怎样把战局编排为凯旋的次序。它们不能替代军需题本、府县志和地方契约所偶然留下的运粮、逃徙、复业和讼案，更不能据此计算每一地的损失或所有人的立场。所谓“降附”在军报中是一种行政结论；对一名改换主人的军官、一个要保护田土的家族或一个被军队经过的村落而言，它可能分别意味着领饷、避祸、被征发或暂时等待。材料能使这些差异局部可见，却不足以把西南居民写成整齐的一种反应。

海峡的难题则由船、风和港口把同样的行政愿望拆散。迁界意在截断沿岸对郑氏的接济，改变的却是闽粤渔户出海、盐货转运、耕地利用与亲族往来的条件；复界和海禁的调整也不能仅凭一纸上谕推定为同时发生的生活转变。沿海村落有人迁离，有人在限制松动后设法回返，也有人把生计移往别处。地方志常在后来用恢复的语言回顾这些变化，官府命令则记录其规定而非每一次迁移；两者都应与港口账册、契约和航海记录并读，才可知道哪些地点、哪些群体留下了可说的痕迹。

\subsection{从热兰遮到澎湖：岛内土地、海上补给与接收的限度}

一六六一年，\eventanchor{EQ-1661-ZHENG-TAIWAN-LANDING}{郑成功登陆台湾（1661）}，热兰遮城外的围困把海上争夺压缩为船只何时到港、粮食能否支撑、港湾是否可用的连续计算。郑氏需要一个能供给军队并容纳其政治组织的基地；荷兰东印度公司的据点同时嵌在贸易、税收与台湾西南既有的居民关系里。登陆和围城可以确定为军事行动的开端，却不能由此替汉人移民、不同原住民社群、港口商人或荷兰雇员指定同一立场。荷兰城堡日志、清方海防军报和后出的郑氏叙事各自看见不同的敌人、盟友与损失；兵数、伤亡以及各社如何卷入，仍须在这些不相称的记录之间谨慎停留。

郑成功去世后，大陆据点与台湾之间的取舍变得更急迫。\eventanchor{EQ-1663-XIAMEN-JINMEN}{厦门、金门易手（1663）}，并不等于清廷已经在海面上建立无缝的控制；它使郑氏主力退向台湾，也使闽南的岛屿、港口和沿岸村落面对新的航路与征发压力。熟悉水面的降将、地方船户和能否取得船粮，都比地图上两座岛的归属更接近当时的胜负条件。清方军报和荷兰方面的消息可以互相钉住据点的易手，却难以替沿海居民说明一次禁令如何影响其出海，或替郑氏内部还原每一笔补给从何而来。

这种不确定性在\eventanchor{EQ-1681-ZHENG-JING-DEATH}{郑经去世（1681）}后进入东宁的继承问题。郑克塽取得位置，意味着郑氏的宗族、将领与文官已在岛内重新排列；北京得到的情报能够显示清廷怎样估量这个窗口，却不能把郑氏后出的叙事和清方侦报拼接成一场没有空白的宫廷政变。继承冲突削弱了对手协调船队与人力的能力，但不足以单独解释两年后澎湖的结局。福建水师还必须筹集舰船、等待适航的季节、维持兵粮，并判断沿海据点能否供接应。郑氏的战败和投降改变了海峡两岸的军事关系；清廷在台湾设府，是把福建的行政链向东延伸的起点，而非一枚印信立刻重写岛内土地、渡海限制与不同社群关系的终点。

因此，所谓海疆的“接收”必须拆成几条速度不同的线来看。清廷可以任命官员、部署驻防、规定渡台与贸易；船户和移民仍要面对风候、航程、债务与港口机会，岛内村社则在垦地、水源、税役和与官府交涉的具体处境中作出反应。原住民社群留下的直接自述尤为稀少，殖民官员、郑氏与清朝文书又都带有土地和治理的诉求。它们足以提示行政接收改变了谁可以征税、谁能够通航，却不能让我们替每一处聚落断定接受、抵抗或融入。海峡不是陆地政令的末端，而是财政、人群流动和军队补给不断相互改写的空间。

\subsection{草原、河流与高原：多语关系不是一张完成的疆图}

北方的秩序也以迁徙和供给开始，而不是以会盟仪式结束。噶尔丹进攻喀尔喀之后，部众、牲畜与领主向南移动，清廷所面对的既是军事求援，也是安置、赈济、水草和道路的问题。台站要转运消息与物资，草原上的关系却不能由北京单方面排定：请求保护的领主、原有的牧地使用者、寺院网络和负责调度的人，都在同一片季节性移动的空间里。兵力与迁徙人数在不同战报中常有夸张或省略；把它们化为一个精确规模，反而会掩盖记录者究竟是在报功、请饷还是争取援助。

在这样的背景下，\eventanchor{EQ-1691-DOLONNOR}{多伦诺尔会盟（1691）}把清廷的军事保护、封爵礼仪、盟旗编制和物资调度连接起来。会盟能够确证清廷希望以何种名义组织与喀尔喀诸部的关系；理藩院的满、蒙文档案又能使后续的赈济、编制和事务往来局部可见。但“归附”不是可以套在所有人身上的完成式。牧地怎样使用、部际纠纷由谁居中、领主保留何种选择，仍须沿着不同语言的文书和地方记忆逐项追踪。多伦的仪式若脱离这些问题，只会把一段谈判关系误读为地图上已经涂好的边界。

一六八九年尼布楚的交涉提供了另一种限度。河名、山名和地名被写入满文、拉丁文、俄文等文本，是使双方能够暂时结束阿尔巴津危机的办法；翻译、地物对应和使团行程本身就是谈判的一部分。条约的文字可证明签署者承认了怎样的程序和措辞，不能自行量出一条现代精确线，也不能证明河岸的猎人、商旅或巡查者已经知道并接受同一种分类。清方实录所记的批准与赏赐、俄方使团文书所记的交涉、黑龙江地方志及驿路和贸易记录所见的实际往来，应当并置：前两者说明两个国家怎样叙述一次外交成果，后两者才可能让边界在何时、何地被实践的问题显出轮廓。

通往青海、西藏的关系同样不能由册封、驻军或后来的地图倒推为稳定控制。寺院、地方政教力量、蒙古盟军、商旅与沿途的向导、驮运者共同决定道路在某个季节是否可通；军队进入拉萨前后，粮秣、病损与地方合作都限制着它能实行什么。清廷驻防和理藩院文书留下的是一条行政视线，藏文寺院文书和地方传统则有自己的时间线，且彼此保存并不均衡。把它们并读，不是为了制造一个无缝的“真实全景”，而是为了不把其中任何一种制度语言误作所有行动者的经验。

\subsection{把材料的视野留在它能够到达的地方}

从北京望出去，早期清朝似乎由诏令、任命、盟约和凯旋组成；从军需与地方材料望回去，它又是许多不稳定的接口。部院可以下令撤藩或设府，却要通过地方官、将领、旗务机构和中介人员获得粮、船、马匹与消息。边地的行政不是命令距离京城较远的简化版：在云南，它受山路、仓储和战后安置牵制；在台湾，它受季风、港湾、渡海限制和土地关系牵制；在草原与高原，它受牧地、寺院、翻译和季节性道路牵制。不同限制不是中央权力的注脚，而是决定政策在何处改变形状的条件。

《平定三逆方略》和《清圣祖实录》使人看见朝廷如何排出招降、军功与处分的顺序。题本、地方志、郑氏材料以及满、蒙、藏文档案则让供给、迁徙和地方关系在若干节点露出侧面。任何一类材料都不能独自解决另一个层次的问题：一份上谕证明命令曾被发布，不证明村社如何执行；一份条约证明文字曾被同意，不证明地名和界线已经被共同实践；一条地方记录能够照亮一处道路或港口，也不能替整个区域发言。兵额、战果、人口损失和航运规模尤其需要逐项注明记录者、地域与目的，并与对手、地方和跨境材料互校。

这也改变了我们阅读“平定”“归附”或“安民”这类官府语言的方式。它们首先说明国家试图把人、地和关系纳入何种行政分类，而不是居民、盟旗、寺院或绿洲社会已经毫无保留地接受了这种分类。钱穆关于制度如何把中央意图送入边地的提问，能帮助我们追问这些接口；吕思勉关于社会经济条件的提问，则提醒我们回到河流两岸的贸易、牧猎、移徙与驻防。两种提问都不能代替史料的限度：在直接证词缺失处，叙事应保留不知道谁怎样选择的空间，而不以帝国的文书替他们补写回答。

\subsection{合卷：征服的所得与未结的账}

清初的所得，是把军队占领、皇帝驻跸、部院文书和部分地方合作接成一个可以持续运作的朝廷。北京提供了旧帝国的行政中心；撤藩削除了将领家族长期据有军政的危险；海峡、蒙古和西南的经营又扩大了清廷能够动员资源的范围。

这些收益带有不能消失的成本。降官、旗人、降军、民夫、移民、盟旗与海疆居民并不承担同一种征发，也不在同一时刻接受同一种秩序。中央以册封、驻防、海禁、赈济和设治解决的，常是下一轮牧地、饷源、迁徙或身份争议的开始。帝国的边界因而不是一条完成线，而是无数地方关系被反复接入的结果。

\subsection{制度得失}

\textbf{所得：} 摄政、部院接收与旗军调度使新王朝能够将一次军事突破转化为跨区域的行政架构；撤藩削弱了独立军政中心。

\textbf{代价：} 财政必须长期供养驻军、转运与身份体系；地方社会承担战后复垦、移民、海禁和征发的差异性风险。

\textbf{留下的压力：} 多语边地与海峡需要的不是单一命令，而是持续的中介、补给和谈判。这些压力将贯穿所谓“盛世”的扩张。

